% Lectura y comentario de poema ; figuras retóricas ; debate : alimentación sana vs comida basura — Espagnol Tle A — Cours signé OnBuch+
% Programme officiel MINESEC. Compiler avec : tectonic E11-poema-alimentacion-sana-comida-basura.tex
\newcommand{\DOCMATIERE}{Espagnol}
\newcommand{\DOCNIVEAU}{Tle A}
\newcommand{\DOCMODULE}{Módulo 2 : Environnement, santé et bien-être}
\newcommand{\DOCLECON}{E11}
\newcommand{\DOCTITRE}{Lectura y comentario de poema ; figuras retóricas ; debate : alimentación sana vs comida basura}
% OnBuch+ — préambule « cours signés » ENRICHI (compilé avec tectonic / XeLaTeX)
% Système visuel « Pop » d'OnBuch+ : fond crème (jamais blanc), bordures encre
% épaisses, ombres « dures » décalées, titres Archivo Black en capitales.
% Version MATHÉMATIQUES : graphes (pgfplots/tikz), boîtes pédagogiques
% (définition, propriété, méthode, attention, le savais-tu, exercice résolu,
% exercice, TP, bilan), page de garde et signature OnBuch+ ancrée en bas.
%
% Macros attendues dans main.tex AVANT \input{preamble} :
%   \DOCMATIERE  \DOCNIVEAU  \DOCMODULE  \DOCLECON (numéro)  \DOCTITRE
\documentclass[11pt]{article}

\usepackage{fontspec}
\usepackage[spanish,french]{babel} % français = langue principale ; l'espagnol via \esp{..} / espagnol
\usepackage[a4paper,margin=2cm,top=3.4cm,bottom=2.8cm,headheight=1.4cm,footskip=1.3cm]{geometry}
\usepackage{amsmath,amssymb,amsfonts}
\usepackage{mathtools} % \xmapsto, \coloneqq... souvent utilisés par les modèles
\usepackage{cancel} % simplifications barrées (bilans de forces)
% \abs{z} (module, valeur absolue) et \norm{u} (norme), utilisés par les modèles
\providecommand{\abs}{}\let\abs\relax\DeclarePairedDelimiter{\abs}{\lvert}{\rvert}
\providecommand{\norm}{}\let\norm\relax\DeclarePairedDelimiter{\norm}{\lVert}{\rVert}
\usepackage[table]{xcolor} % [table] : fournit aussi \rowcolors (alternance de lignes)
\usepackage{pagecolor}
\usepackage{tikz}
\usetikzlibrary{arrows.meta,positioning,calc,shapes.geometric,shapes.misc,shapes.arrows,decorations.pathmorphing,decorations.pathreplacing,decorations.markings,patterns,shadows,mindmap,trees,fit,backgrounds,matrix,intersections,angles,quotes}
\usepackage{pgfplots}
\usepackage[version=4]{mhchem} % équations nucléaires \ce{^{235}_{92}U}
\usepackage[european,siunitx]{circuitikz} % schémas électriques
\pgfplotsset{compat=1.18}
\usepgfplotslibrary{fillbetween,groupplots} % aires entre courbes (calcul intégral)
\usepackage{enumitem}
\usepackage{fancyhdr}
% [most] charge toutes les bibliothèques tcolorbox (listings, minted, raster,
% documentation...) et gonfle inutilement la mémoire TeX sur un document long
% (~300 boîtes) : on ne charge que ce qui est réellement utilisé.
\usepackage[skins,breakable]{tcolorbox}
\usepackage{graphicx}
\usepackage{colortbl}
\usepackage{booktabs}
\usepackage{tabularx}
\usepackage{diagbox} % case d'en-tête diagonale des tableaux à double entrée
% Colonnes extensibles centrée / gauche / droite (C, L, R), souvent utilisées par les modèles
\newcolumntype{C}{>{\centering\arraybackslash}X}
\newcolumntype{L}{>{\raggedright\arraybackslash}X}
\newcolumntype{R}{>{\raggedleft\arraybackslash}X}
\usepackage{array}
\usepackage{multicol}
\usepackage{siunitx}
% parse-numbers=false : accepte aussi \SI{2,5 \times 10^{-3}}{...}
\sisetup{locale=FR,output-decimal-marker={,},group-minimum-digits=4,parse-numbers=false}
\DeclareSIUnit\ohm{\ensuremath{\symup{\Omega}}} % U+2126 absent de la police
\DeclareSIUnit\division{div} % réglages d'oscilloscope (ms/div, V/div)
\DeclareSIUnit\tr{tr} % tours (vitesse de rotation, ex. \SI{1500}{\tr\per\minute})
\DeclareSIUnit\molar{M} % concentration molaire (ex. \SI{3}{\molar}, \SI{1}{\micro\molar})
\DeclareSIUnit\an{an} % année (le modèle confond parfois avec \year, un compteur TeX) — ex. \SI{5}{\kilo\meter\per\an}
\DeclareSIUnit\FCFA{FCFA} % franc CFA (ex. \SI{500}{\FCFA\per\kilo\gram})
\DeclareSIUnit\degreeBrix{\textdegree Brix} % teneur en sucre d'un sirop/jus (réfractométrie)
\DeclareSIUnit\month{mois} % le modèle utilise parfois \month, un compteur TeX, comme unité de durée
\usepackage{qrcode}
% \degree : commande fréquemment utilisée par les modèles pour un angle ou une
% température en dehors de \SI{}{} (non fournie par les paquets chargés).
\providecommand{\degree}{^{\circ}}
% \qed (fin de démonstration), utilisé par les modèles sans amsthm
\providecommand{\qed}{\hfill$\square$}
% \barre : double barre des valeurs interdites dans un tableau de variations (tkz-tab)
\providecommand{\barre}{\ensuremath{\|}}
% environnement proof (amsthm non chargé) utilisé par les modèles
\ifdefined\proof\else\newenvironment{proof}[1][Démonstration]{\par\noindent\textit{#1.}\ }{\qed\par}\fi
\usepackage{lastpage}
\usepackage{hyperref}
\hypersetup{hidelinks}

% ============================================================
% Palette « Pop » OnBuch+ — mêmes hex que la classe P (plus_ui.dart)
% ============================================================
\definecolor{poporange}{HTML}{F59321}
\definecolor{popdark}{HTML}{E07A0C}
\definecolor{popcream}{HTML}{FFF6E8}
\definecolor{popink}{HTML}{1C1714}
\definecolor{poppurple}{HTML}{7A5AE0}
\definecolor{popgreen}{HTML}{2E9E6A}
\definecolor{poppink}{HTML}{E0457B}
\definecolor{popblue}{HTML}{2F7DE1}
\definecolor{popgold}{HTML}{C98A12}
\definecolor{popmuted}{HTML}{8A7F76}
\definecolor{popline}{HTML}{EADFD2}
\definecolor{popgray}{HTML}{8A7F76}
\colorlet{popgrey}{popgray}
% Alias tolérants (le modèle nomme parfois les couleurs différemment)
\colorlet{popwhite}{white}
\colorlet{popblack}{popink}
\colorlet{popred}{poppink}
\colorlet{popyellow}{popgold}
% Teintes claires (fonds de tableaux/encadrés) — le modèle nomme parfois
% les couleurs avec un suffixe L (ex. popblueL) sans les avoir définies.
\colorlet{poporangeL}{poporange!20}
\colorlet{popdarkL}{popdark!20}
\colorlet{poppurpleL}{poppurple!20}
\colorlet{popgreenL}{popgreen!20}
\colorlet{poppinkL}{poppink!20}
\colorlet{popblueL}{popblue!20}
\colorlet{popgoldL}{popgold!20}
\colorlet{popredL}{popred!20}
\colorlet{popgrayL}{popgray!20}
% Teintes claires pour fonds de tableaux / zones
\colorlet{popblueL}{popblue!10!white}
\colorlet{poporangeL}{poporange!14!white}
\colorlet{popgreenL}{popgreen!12!white}
\colorlet{poppurpleL}{poppurple!12!white}
\colorlet{poppinkL}{poppink!12!white}
\colorlet{popgoldL}{popgold!14!white}

\pagecolor{popcream}

% ============================================================
% Typographie
% ============================================================
\defaultfontfeatures{Ligatures=TeX}
\setmainfont{PlusJakartaSans-Medium.ttf}[Path=../../../fonts/,BoldFont=PlusJakartaSans-Bold.ttf]
% Maths en sans-serif (Fira Math) avec chiffres et lettres droites de Plus
% Jakarta, pour que les formules se fondent dans le texte.
\usepackage[math-style=ISO,bold-style=ISO]{unicode-math}
\setmathfont{FiraMath-Regular.otf}[Path=../../../fonts/]
\setmathfont{PlusJakartaSans-Medium.ttf}[Path=../../../fonts/,range={up/num,up/{Latin,latin},\mathup}]
\usepackage{icomma} % virgule décimale sans espace en mode math
% Caractères mathématiques Unicode absents des polices de texte (Plus Jakarta,
% Archivo Black) : rendus par leur équivalent mathématique, en texte comme en
% formule — sinon ils s'affichent en carré vide (ex. « Arithmétique dans ℤ »).
\usepackage{newunicodechar}
\newunicodechar{ℝ}{\ensuremath{\mathbb{R}}}
\newunicodechar{ℤ}{\ensuremath{\mathbb{Z}}}
\newunicodechar{ℂ}{\ensuremath{\mathbb{C}}}
\newunicodechar{ℕ}{\ensuremath{\mathbb{N}}}
\newunicodechar{ℚ}{\ensuremath{\mathbb{Q}}}
\newunicodechar{𝔻}{\ensuremath{\mathbb{D}}}
\newunicodechar{α}{\ensuremath{\alpha}}
\newunicodechar{β}{\ensuremath{\beta}}
\newunicodechar{γ}{\ensuremath{\gamma}}
\newunicodechar{δ}{\ensuremath{\delta}}
\newunicodechar{ε}{\ensuremath{\varepsilon}}
\newunicodechar{θ}{\ensuremath{\theta}}
\newunicodechar{λ}{\ensuremath{\lambda}}
\newunicodechar{μ}{\ensuremath{\mu}}
\newunicodechar{σ}{\ensuremath{\sigma}}
\newunicodechar{τ}{\ensuremath{\tau}}
\newunicodechar{φ}{\ensuremath{\varphi}}
\newunicodechar{ψ}{\ensuremath{\psi}}
\newunicodechar{ω}{\ensuremath{\omega}}
\newunicodechar{Σ}{\ensuremath{\Sigma}}
% majuscules produites par \MakeUppercase dans les titres (α-aminés -> Α)
\newunicodechar{Α}{\ensuremath{\alpha}}
\newunicodechar{Β}{\ensuremath{\beta}}
\newunicodechar{Γ}{\ensuremath{\Gamma}}
\newunicodechar{Δ}{\ensuremath{\Delta}}
\newunicodechar{Ε}{\ensuremath{\varepsilon}}
\newunicodechar{Θ}{\ensuremath{\Theta}}
\newunicodechar{Λ}{\ensuremath{\Lambda}}
\newunicodechar{Μ}{\ensuremath{\mu}}
\newunicodechar{Τ}{\ensuremath{\tau}}
\newunicodechar{Φ}{\ensuremath{\Phi}}
\newunicodechar{Ψ}{\ensuremath{\Psi}}
\newunicodechar{Ω}{\ensuremath{\Omega}}
\newunicodechar{↦}{\ensuremath{\mapsto}}
\newunicodechar{∈}{\ensuremath{\in}}
\newunicodechar{∉}{\ensuremath{\notin}}
\newunicodechar{∧}{\ensuremath{\wedge}}
\newunicodechar{⇔}{\ensuremath{\Leftrightarrow}}
\newunicodechar{⟺}{\ensuremath{\Longleftrightarrow}}
\newunicodechar{⇒}{\ensuremath{\Rightarrow}}
\newunicodechar{⟹}{\ensuremath{\Longrightarrow}}
\newunicodechar{∘}{\ensuremath{\circ}}
\newunicodechar{∪}{\ensuremath{\cup}}
\newunicodechar{∩}{\ensuremath{\cap}}
\newunicodechar{≡}{\ensuremath{\equiv}}
\newunicodechar{⊂}{\ensuremath{\subset}}
\newunicodechar{⊥}{\ensuremath{\perp}}
\newunicodechar{∃}{\ensuremath{\exists}}
\newunicodechar{∀}{\ensuremath{\forall}}
\newunicodechar{∛}{\ensuremath{\sqrt[3]{\,}}}
\newunicodechar{∜}{\ensuremath{\sqrt[4]{\,}}}
\newunicodechar{⃗}{\ensuremath{{}^{\to}}}
\newunicodechar{ⁿ}{\ifmmode^{n}\else\textsuperscript{\ensuremath{n}}\fi}
\newunicodechar{ˣ}{\ifmmode^{x}\else\textsuperscript{\ensuremath{x}}\fi}
\newunicodechar{ᵇ}{\ifmmode^{b}\else\textsuperscript{\ensuremath{b}}\fi}
\newunicodechar{ʳ}{\ifmmode^{r}\else\textsuperscript{\ensuremath{r}}\fi}
\newunicodechar{ᵘ}{\ifmmode^{u}\else\textsuperscript{\ensuremath{u}}\fi}
\newunicodechar{ʸ}{\ifmmode^{y}\else\textsuperscript{\ensuremath{y}}\fi}
\newunicodechar{ᵃ}{\ifmmode^{a}\else\textsuperscript{\ensuremath{a}}\fi}
\newunicodechar{ᵉ}{\ifmmode^{e}\else\textsuperscript{\ensuremath{e}}\fi}
\newunicodechar{ᵏ}{\ifmmode^{k}\else\textsuperscript{\ensuremath{k}}\fi}
\newunicodechar{ᵖ}{\ifmmode^{p}\else\textsuperscript{\ensuremath{p}}\fi}
\newunicodechar{ᴺ}{\ifmmode^{N}\else\textsuperscript{\ensuremath{N}}\fi}
\newunicodechar{ᶜ}{\ifmmode^{c}\else\textsuperscript{\ensuremath{c}}\fi}
\newunicodechar{ʰ}{\ifmmode^{h}\else\textsuperscript{\ensuremath{h}}\fi}
\newunicodechar{ᵠ}{\ifmmode^{q}\else\textsuperscript{\ensuremath{q}}\fi}
\newunicodechar{ₙ}{\ifmmode_{n}\else\textsubscript{\ensuremath{n}}\fi}
\newunicodechar{ₐ}{\ifmmode_{a}\else\textsubscript{\ensuremath{a}}\fi}
\newunicodechar{ᵢ}{\ifmmode_{i}\else\textsubscript{\ensuremath{i}}\fi}
\newunicodechar{ᵣ}{\ifmmode_{r}\else\textsubscript{\ensuremath{r}}\fi}
\newunicodechar{ₖ}{\ifmmode_{k}\else\textsubscript{\ensuremath{k}}\fi}
\newunicodechar{ᵦ}{\ifmmode_{\beta}\else\textsubscript{\ensuremath{\beta}}\fi}
\newunicodechar{ₓ}{\ifmmode_{x}\else\textsubscript{\ensuremath{x}}\fi}
\newfontfamily\popdisplay{ArchivoBlack-Regular.ttf}[Path=../../../fonts/]
\newfontfamily\poplogo{SpaceGrotesk-Bold.ttf}[Path=../../../fonts/]
\newfontfamily\popsemi{PlusJakartaSans-SemiBold.ttf}[Path=../../../fonts/]
\newfontfamily\popxbold{PlusJakartaSans-ExtraBold.ttf}[Path=../../../fonts/]
\newfontfamily\popxboldls{PlusJakartaSans-ExtraBold.ttf}[Path=../../../fonts/,LetterSpace=3.0]

\setlist[enumerate]{leftmargin=1.7em,itemsep=5pt,topsep=4pt}
\setlist[itemize]{leftmargin=1.7em,itemsep=4pt,topsep=4pt,label={\color{poporange}\textbullet}}
\setlength{\parindent}{0pt}
\setlength{\parskip}{5pt}
\renewcommand{\arraystretch}{1.25}

% pgfplots : style Pop par défaut
\pgfplotsset{
  every axis/.append style={axis line style={popink,line width=1pt},tick style={popink},
    grid style={popline,line width=0.5pt},label style={font=\small},tick label style={font=\footnotesize}},
  popcurve/.style={poporange,line width=1.8pt,no markers,smooth},
}
% Même style disponible dans un \draw TikZ simple (hors pgfplots)
\tikzset{popcurve/.style={poporange,line width=1.8pt}}
\tikzset{popgrid/.style={popline,very thin}}
\colorlet{popcurve}{poporange} % le nom du style est parfois utilisé comme couleur

% ============================================================
% Pill / squiggle
% ============================================================
\newcommand{\poppill}[3]{%
  \tikz[baseline=-0.55ex]\node[draw=popink,line width=1.2pt,rounded corners=2.6pt,%
    fill=#2,inner xsep=6.5pt,inner ysep=3pt]{%
    \popxboldls\fontsize{7}{7}\selectfont\color{#3}\MakeUppercase{#1}};%
}
\newcommand{\popsquiggle}[1]{%
  \tikz[baseline=-0.4ex]{
    \draw[#1,line width=2.6pt,line cap=round]
      plot [smooth] coordinates {(0,0.09) (0.34,0.01) (0.68,0.21) (1.02,0.01) (1.36,0.10)};
  }%
}
% Mot-clé surligné (marqueur orange sous le texte)
\newcommand{\cle}[1]{{\popxbold\color{popdark}#1}}

% ============================================================
% En-tête / pied
% ============================================================
\pagestyle{fancy}
\fancyhf{}
\renewcommand{\headrulewidth}{0pt}
\renewcommand{\footrulewidth}{0pt}
\lhead{\raisebox{-0.15ex}{\poplogo\fontsize{13}{13}\selectfont\color{popink}OnBuch\color{poporange}+}}
\rhead{\poppill{\DOCMATIERE\ \textbullet\ \DOCNIVEAU\ \textbullet\ Leçon \DOCLECON}{white}{popink}}
\lfoot{\popxbold\fontsize{7.6}{7.6}\selectfont\color{popmuted}\MakeUppercase{Cours signé OnBuch+}\ \textbullet\ {\popsemi\DOCTITRE}}
\rfoot{\tikz[baseline=-0.6ex]\node[rounded corners=8.5pt,fill=popink,minimum height=17pt,minimum width=17pt,inner xsep=5pt,inner sep=0]{\popxbold\fontsize{8}{8}\selectfont\color{popcream}\thepage\,/\,\pageref*{LastPage}};}

% ============================================================
% Titres
% ============================================================
\newcommand{\chaptitre}[2]{%
  \begin{tcolorbox}[enhanced,colback=poporange,colframe=popink,boxrule=2.6pt,arc=11pt,
      drop shadow=popink,left=15pt,right=15pt,top=12pt,bottom=11pt,before skip=3pt,after skip=18pt]
    \poppill{#2}{popink}{popcream}\par\vspace{9pt}
    {\raggedright\hyphenpenalty=10000\exhyphenpenalty=10000\popdisplay\fontsize{22}{24}\selectfont\color{popcream}\MakeUppercase{#1}\par}
  \end{tcolorbox}
}
% Section numérotée : pastille orange avec le numéro + titre Archivo
\newcounter{popsec}
\newcommand{\coursec}[1]{%
  \stepcounter{popsec}\setcounter{popsub}{0}%
  \par\vspace{16pt}\noindent
  \begin{minipage}{\linewidth}
  \tikz[baseline=-0.7ex]\node[draw=popink,line width=1.6pt,fill=poporange,rounded corners=4pt,minimum size=22pt,inner sep=0]{\popdisplay\fontsize{12}{12}\selectfont\color{popcream}\thepopsec};%
  \hspace{8pt}{\popdisplay\fontsize{15}{17}\selectfont\color{popink}\MakeUppercase{#1}}\par
  \vspace{4pt}\noindent{\color{poporange}\rule{36pt}{3.6pt}}
  \end{minipage}\par\nopagebreak\vspace{8pt}
}
\newcounter{popsub}
\newcommand{\courssub}[1]{%
  \stepcounter{popsub}%
  \par\vspace{9pt}\noindent{\popxbold\fontsize{11.5}{13}\selectfont\color{popink}{\color{poporange}\thepopsec.\thepopsub}\hspace{6pt}#1}\par\nopagebreak\vspace{4pt}
}

% ============================================================
% Boîtes pédagogiques — même recette : fond blanc, bordure épaisse colorée,
% ombre dure encre, badge plein. Titre optionnel [#1].
% ============================================================
\tcbset{popbase/.style={breakable,enhanced,colback=white,boxrule=2.3pt,arc=9pt,
  shadow={3pt}{-3pt}{0pt}{popink},left=14pt,right=14pt,top=11pt,bottom=11pt,before skip=11pt,after skip=15pt,
  fontupper=\popsemi\fontsize{10.6}{14.5}\selectfont\color{popink},
  fontlower=\popsemi\fontsize{10.6}{14.5}\selectfont\color{popink}}}
% Coche dessinée (le glyphe ✓ n'existe pas dans les polices du document)
\newcommand{\popcheck}[1]{\tikz[baseline=-0.5ex]\draw[#1,line width=1.6pt,line cap=round,line join=round](-0.13,0.0)--(-0.04,-0.09)--(0.13,0.1);}
\providecommand{\coche}{\popcheck{popgreen}}
\providecommand{\resultat}[1]{\par\smallskip\noindent\fcolorbox{popgreen}{popgreenL}{\parbox{\dimexpr\linewidth-2\fboxsep-2\fboxrule\relax}{\textbf{Résultat.} #1}}\par\smallskip}
\newcommand{\popboxhead}[3]{\poppill{#1}{#2}{white}\ifx\relax#3\relax\else\hspace{6pt}{\popxbold\fontsize{10.6}{12}\selectfont\color{popink}#3}\fi\par\vspace{7pt}}

\newtcolorbox{aretenir}[1][]{popbase,colframe=popgreen,before upper={\popboxhead{À retenir}{popgreen}{#1}}}
% Texte en espagnol (document de lecture, dialogue, poème, extrait) : cadre
% d'étude. Un titre optionnel (source, auteur) s'affiche dans l'en-tête.
\newtcolorbox{textebox}[1][]{popbase,colframe=popblue,before upper={\popboxhead{Texto}{popblue}{#1}\begin{otherlanguage}{spanish}},after upper={\end{otherlanguage}}}
% Marques ✓ / ✗ (forme correcte / fautive) souvent utilisées par les modèles
\providecommand{\cross}{\textcolor{poppink}{\ensuremath{\times}}}
\providecommand{\xmark}{\cross}\providecommand{\crossmark}{\cross}\providecommand{\cmark}{\ensuremath{\checkmark}}
% Ligne de réponse / justification à compléter (inventée par les modèles)
\providecommand{\justification}[1]{\par\noindent\textit{Justification~:} \dotfill\par}
% \textipa (package tipa non chargé) : la police ne couvre pas l'API -> simple passage
\providecommand{\textipa}[1]{#1}
% Soulignement (ulem non chargé) : \uline, \ul
\providecommand{\uline}[1]{\underline{#1}}\providecommand{\ul}[1]{\underline{#1}}
% Mot ou phrase en espagnol à l'intérieur d'un paragraphe français
\newcommand{\esp}[1]{\ifmmode\text{\textit{#1}}\else\begingroup\selectlanguage{spanish}\itshape #1\endgroup\fi}
\newtcolorbox{exemplebox}[1][]{popbase,colframe=poporange,before upper={\popboxhead{Exemple}{poporange}{#1}}}
\newtcolorbox{definition}[1][]{popbase,colframe=popblue,before upper={\popboxhead{Définition}{popblue}{#1}}}
\newtcolorbox{propriete}[1][]{popbase,colframe=poppurple,before upper={\popboxhead{Propriété}{poppurple}{#1}}}
\newtcolorbox{methode}[1][]{popbase,colframe=popgold,before upper={\popboxhead{Méthode}{popgold}{#1}}}
\newtcolorbox{attention}[1][]{popbase,colframe=poppink,before upper={\popboxhead{Attention}{poppink}{#1}}}
\newtcolorbox{experience}[1][]{popbase,colframe=popblue,colback=popblueL,before upper={\popboxhead{Expérience}{popblue}{#1}}}
% « Le savais-tu ? » : carte violette pleine (contexte, culture, Cameroun)
\newtcolorbox{savaistu}[1][]{popbase,colback=poppurple,colframe=popink,
  fontupper=\popsemi\fontsize{10.4}{14.2}\selectfont\color{white},
  before upper={\poppill{Le savais-tu\,?}{white}{poppurple}\ifx\relax#1\relax\else\hspace{6pt}{\popxbold\color{white}#1}\fi\par\vspace{7pt}}}
% Exercice résolu : énoncé en haut, \tcblower puis la solution sur fond crème
\newcounter{popexo}\newcounter{popexr}
\newtcolorbox{exoresolu}[1][]{popbase,colframe=poporange,colbacklower=popcream,
  segmentation style={popink,line width=1.2pt,dashed},
  before upper={\stepcounter{popexr}\popboxhead{Exercice résolu \thepopexr}{poporange}{#1}},
  before lower={\poppill{Solution}{popink}{popcream}\par\vspace{6pt}}}
% Exercice (non résolu en place) — la correction est dans la partie Corrigés
\newtcolorbox{exercice}[1][]{popbase,colframe=popink,
  before upper={\stepcounter{popexo}\popboxhead{Exercice \thepopexo}{popink}{#1}}}
% Corrigé
\newtcolorbox{corrige}[1][]{popbase,colframe=popgreen,colback=popgreenL,
  before upper={\popboxhead{Corrigé}{popgreen}{#1}}}
% Objectifs / prérequis / bilan
\newtcolorbox{objectifs}{popbase,colframe=popink,colback=poporangeL,
  before upper={\popboxhead{Objectifs de la leçon}{popink}{}}}
\newtcolorbox{prerequis}{popbase,colframe=popink,colback=popblueL,
  before upper={\popboxhead{Prérequis}{popblue}{}}}
\newtcolorbox{bilan}{popbase,colframe=popink,colback=white,boxrule=2.8pt,
  before upper={\popboxhead{Fiche bilan}{poporange}{L'essentiel à connaître pour l'examen}}}
% Figure encadrée avec légende
\newtcolorbox{popfigure}[1][]{enhanced,breakable=false,colback=white,colframe=popline,boxrule=1.4pt,arc=8pt,
  left=10pt,right=10pt,top=10pt,bottom=8pt,before skip=10pt,after skip=12pt,halign=center,
  lower separated=false,halign lower=center,
  fontlower=\popsemi\fontsize{9}{11.5}\selectfont\color{popmuted},
  #1}
\newcommand{\legende}[1]{\par\vspace{6pt}{\popsemi\fontsize{9}{11.5}\selectfont\color{popmuted}#1}}

% ============================================================
% Signature OnBuch+
% ============================================================
\newcommand{\poplogoinline}[1]{{\poplogo\fontsize{#1}{#1}\selectfont\color{popink}OnBuch\color{poporange}+}}
\newcommand{\signatureonbuch}{%
  \begin{tcolorbox}[enhanced,colback=popink,colframe=popink,boxrule=2.2pt,arc=10pt,
      drop shadow=poporange,left=14pt,right=14pt,top=10pt,bottom=10pt,before skip=0pt,after skip=0pt]
    \tikz[baseline=-0.7ex]\node[circle,fill=popgreen,draw=popcream,line width=1.3pt,minimum size=22pt,inner sep=0]{\popcheck{white}};%
    \hspace{9pt}\begin{minipage}[c]{0.62\linewidth}
      {\popxbold\fontsize{12}{13}\selectfont\color{popcream}Signé\ \textbullet\ {\poplogo OnBuch\color{poporange}+}}\par\vspace{2pt}
      {\raggedright\popsemi\fontsize{8}{10.5}\selectfont\color{popline}\MakeUppercase{Équipe pédagogique OnBuch+}\\\MakeUppercase{Conforme au programme officiel MINESEC}\par}
    \end{minipage}\hfill
    {\poplogo\fontsize{22}{22}\selectfont\color{popcream}OnBuch\color{poporange}+}
  \end{tcolorbox}%
}

% ============================================================
% Page de garde
% #1 titre · #2 matière · #3 niveau · #4 module · #5 n° leçon · #6 durée ·
% #7 description / objectifs (texte ou itemize)
% ============================================================
\newcommand{\pagegarde}[7]{%
  \thispagestyle{empty}
  \newgeometry{a4paper,margin=1.7cm,top=1.6cm,bottom=1.4cm}%
  \begin{tikzpicture}[remember picture,overlay]
    \fill[poporange!18] ([xshift=-3.2cm,yshift=-3.6cm]current page.north east) circle (4.6cm);
    \fill[poppurple!14] ([xshift=2.4cm,yshift=7.6cm]current page.south west) circle (3.8cm);
  \end{tikzpicture}%
  {\poplogo\fontsize{30}{30}\selectfont\color{popink}OnBuch\color{poporange}+}\hfill
  \raisebox{0.6ex}{\poppill{Cours signé \textbullet\ Premium}{popink}{popcream}}\par
  \vspace{1pt}\popsquiggle{poppurple}\par
  \vspace{8pt}
  \begin{tcolorbox}[enhanced,colback=poporange,colframe=popink,boxrule=2.8pt,arc=13pt,
      drop shadow=popink,left=18pt,right=18pt,top=12pt,bottom=13pt,after skip=10pt]
    \poppill{#2 \textbullet\ #3}{popink}{popcream}\hspace{5pt}\poppill{#4}{white}{popink}\par\vspace{8pt}
    {\popdisplay\fontsize{14}{14}\selectfont\color{popink}LEÇON #5}\par\vspace{4pt}
    {\raggedright\hyphenpenalty=10000\exhyphenpenalty=10000\popdisplay\fontsize{25}{27}\selectfont\color{popcream}\MakeUppercase{#1}\par}
    \vspace{8pt}
    {\popxbold\fontsize{9.5}{9.5}\selectfont\color{popink}\MakeUppercase{Durée conseillée : #6}}
  \end{tcolorbox}
  \begin{tcolorbox}[enhanced,colback=white,colframe=popink,boxrule=2.2pt,arc=10pt,
      drop shadow=popink,left=16pt,right=16pt,top=10pt,bottom=10pt,after skip=8pt]
    \poppill{Dans cette leçon}{poppurple}{white}\par\vspace{5pt}
    \popsemi\fontsize{9.3}{12.6}\selectfont\color{popink}#7
  \end{tcolorbox}
  \noindent\begin{minipage}[t]{0.487\linewidth}
    \begin{tcolorbox}[enhanced,colback=popgreen,colframe=popink,boxrule=2.2pt,arc=10pt,
        drop shadow=popink,left=11pt,right=11pt,top=10pt,bottom=10pt,height=3.1cm,valign=center]
      \tcbox[colback=white,colframe=popink,boxrule=1.3pt,arc=5pt,boxsep=0pt,nobeforeafter,
        left=3pt,right=3pt,top=3pt,bottom=3pt,valign=center]{\qrcode[height=1.9cm]{https://play.google.com/store/apps/details?id=com.luvvix.onbuch&hl=fr}}%
      \hspace{8pt}\begin{minipage}[c]{0.5\linewidth}
        {\popdisplay\fontsize{11}{12.5}\selectfont\color{white}\MakeUppercase{Télécharge OnBuch}}\par\vspace{4pt}
        {\raggedright\popsemi\fontsize{8.4}{11}\selectfont\color{white}Gratuit \textbullet\ Google Play\\Tuteur IA, annales, concours, résultats\par}
      \end{minipage}
    \end{tcolorbox}
  \end{minipage}\hfill
  \begin{minipage}[t]{0.487\linewidth}
    \begin{tcolorbox}[enhanced,colback=poppurple,colframe=popink,boxrule=2.2pt,arc=10pt,
        drop shadow=popink,left=11pt,right=11pt,top=10pt,bottom=10pt,height=3.1cm,valign=center]
      \tikz[baseline=-0.7ex]\node[circle,fill=white,draw=popink,line width=1.3pt,minimum size=1.6cm,inner sep=0]{\popdisplay\fontsize{24}{24}\selectfont\color{poppurple}+};%
      \hspace{8pt}\begin{minipage}[c]{0.6\linewidth}
        {\popdisplay\fontsize{11}{12.5}\selectfont\color{white}\MakeUppercase{Passe à OnBuch+}}\par\vspace{4pt}
        {\raggedright\popsemi\fontsize{8.4}{11}\selectfont\color{white}Tous les cours signés, Tuteur IA illimité, fiches PDF — onglet OnBuch+ de l'app\par}
      \end{minipage}
    \end{tcolorbox}
  \end{minipage}
  \vspace{8pt}
  \popcontenu
  \vfill
  \signatureonbuch
  \restoregeometry
  \newpage
}

% Rangée « Ce cours contient » (page de garde)
\newcommand{\poptile}[3]{% #1 couleur #2 gros texte #3 légende
  \begin{tcolorbox}[enhanced,colback=#1,colframe=popink,boxrule=2pt,arc=9pt,shadow={2.5pt}{-2.5pt}{0pt}{popink},
      left=6pt,right=6pt,top=9pt,bottom=9pt,halign=center,valign=center,height=1.75cm,width=\linewidth,nobeforeafter]
    {\popdisplay\fontsize{12.5}{13}\selectfont\color{white}\MakeUppercase{#2}}\par\vspace{3pt}
    {\centering\popsemi\fontsize{7.8}{9.5}\selectfont\color{white}#3\par}
  \end{tcolorbox}}
\newcommand{\popcontenu}{%
  \par\noindent{\popxboldls\fontsize{7.5}{7.5}\selectfont\color{popmuted}CE COURS SIGNÉ CONTIENT}\par\vspace{7pt}
  \noindent\begin{minipage}[t]{0.235\linewidth}\poptile{popblue}{Cours}{complet, illustré, pas à pas}\end{minipage}\hfill
  \begin{minipage}[t]{0.235\linewidth}\poptile{poporange}{Méthodes}{et exercices résolus}\end{minipage}\hfill
  \begin{minipage}[t]{0.235\linewidth}\poptile{poppink}{Exercices}{type examen + corrigés}\end{minipage}\hfill
  \begin{minipage}[t]{0.235\linewidth}\poptile{popgreen}{Fiche bilan}{l'essentiel à retenir}\end{minipage}\par}

% Sommaire
\newtcolorbox{sommaire}{enhanced,colback=white,colframe=popink,boxrule=2.2pt,arc=10pt,
  drop shadow=popink,left=18pt,right=18pt,top=13pt,bottom=13pt,before skip=6pt,after skip=20pt,
  before upper={\poppill{Sommaire}{poppurple}{white}\par\vspace{9pt}\popsemi\fontsize{10.6}{14.5}\selectfont\color{popink}}}

% Signature de fin de cours — poussée tout en bas de la dernière page
\newcommand{\findecours}{%
  \par\vfill\nopagebreak
  \begin{center}\popsquiggle{poporange}\end{center}\vspace{4pt}
  \signatureonbuch
}
\AtBeginDocument{\def\llbracket{\mathopen{[\mkern-3.5mu[}}\def\rrbracket{\mathclose{]\mkern-3.5mu]}}} % intervalles d'entiers (glyphe ⟦ absent de la police)
% Glyphes absents de Fira Math : redéfinis à partir de glyphes disponibles
\AtBeginDocument{%
  \def\setminus{\mathbin{\backslash}}\let\smallsetminus\setminus
  \def\vdots{\mathord{\vbox{\baselineskip3.2pt\lineskiplimit0pt\kern4pt\hbox{.}\hbox{.}\hbox{.}}}}%
  \def\ddots{\mathinner{\mkern1mu\raise6.5pt\hbox{.}\mkern2mu\raise3.6pt\hbox{.}\mkern2mu\raise0.7pt\hbox{.}\mkern1mu}}%
  \def\triangle{\mathord{\scalebox{0.9}{$\symup{\Delta}$}}}%
  \def\nabla{\mathord{\raisebox{\depth}{\scalebox{1}[-1]{$\symup{\Delta}$}}}}%
  \def\checkmark{\popcheck{popdark}}%
  \def\star{\mathbin{\ast}}\def\bigstar{\mathord{\ast}}%
  \def\diamond{\mathbin{\scalebox{0.6}{$\Diamond$}}}%
  \def\bigcup{\mathop{\mathchoice{\raisebox{-0.3em}{\scalebox{1.7}{$\cup$}}}{\scalebox{1.25}{$\cup$}}{\cup}{\cup}}\displaylimits}%
  \def\bigcap{\mathop{\mathchoice{\raisebox{-0.3em}{\scalebox{1.7}{$\cap$}}}{\scalebox{1.25}{$\cap$}}{\cap}{\cap}}\displaylimits}%
  \def\lesseqqgtr{\mathrel{\lessgtr}}\def\bigoplus{\mathop{\oplus}\displaylimits}%
}
\providecommand{\ssi}{si et seulement si} % abréviation fréquente
\AtBeginDocument{\providecommand{\wideparen}{\overparen}} % arc de cercle

%% FIGFIX-BEGIN v4 — correctifs globaux des figures (docs/onbuch-generation/tools/figfix)
\usepackage{adjustbox}\usepackage{etoolbox}
% toute figure TikZ/pgfplots plus large que la ligne est réduite à la largeur disponible
\BeforeBeginEnvironment{tikzpicture}{\begin{adjustbox}{max width=\linewidth}}
\AfterEndEnvironment{tikzpicture}{\end{adjustbox}}
% légendes pgfplots lisibles par-dessus les courbes
\pgfplotsset{every axis/.append style={legend style={fill=white,fill opacity=0.92,text opacity=1,draw=black!15,inner sep=3pt}}}
% étiquettes d'arêtes (syntaxe edge["..."]) avec fond blanc
\tikzset{every edge quotes/.append style={fill=white,inner sep=1.2pt}}
%% FIGFIX-END
\begin{document}

\pagegarde{Lectura y comentario de poema ; figuras retóricas ; debate : alimentación sana vs comida basura}{Espagnol}{Tle A}{Módulo 2 : Environnement, santé et bien-être}{E11}{2 h}{Cette leçon mixte propose l'étude complète d'un poème original sur l'alimentation, l'analyse de sa structure et de ses figures de style, l'acquisition du lexique poétique et du vocabulaire du débat nutritionnel, ainsi que la maîtrise des structures grammaticales pour argumenter (subjonctif, connecteurs, comparatifs). Elle prépare aux épreuves écrites et orales du Baccalauréat.
\vspace{4pt}
\begin{multicols}{2}\small
\begin{itemize}
\item Lire, comprendre et commenter un poème en espagnol en identifiant son thème, sa structure et sa versification.
\item Reconnaître et expliquer l'effet de six figures de style : métaphore, comparaison, personification, hyperbole, allitération, anaphore.
\item Utiliser le lexique spécifique de la poésie (verso, estrofa, rima, poeta, musa) et du débat alimentaire (alimentos procesados, nutrientes, dieta equilibrada, comida chatarra).
\item Appliquer la méthode en quatre étapes du commentaire littéraire (contexte, lecture, analyse, synthèse).
\item Produire un commentaire écrit structuré du poème support.
\item Maîtriser les structures grammaticales de l'argumentation : subjonctif après expressions d'opinion, connecteurs logiques, comparatifs d'inégalité et d'égalité.
\end{itemize}\end{multicols}}

\begin{sommaire}
\begin{multicols}{2}
\begin{enumerate}
\item 1. Découverte du poème support : « El duelo del plato »
\item 2. Lexique de la poésie et analyse structurelle détaillée
\item 3. Figures de style : identification, définition et effets
\item 4. Méthodologie du commentaire de poème et commentaire modèle
\item 5. Lexique du débat : alimentation saine vs malbouffe
\item 6. Grammaire \& Traduction : structures d'argumentation et d'opinion
\item 7. Production guidée : débat oral et rédaction d'un article d'opinion
\item Activité d'intégration
\item Méthodes et exercices résolus
\item Exercices
\item Corrigés des exercices
\item Fiche bilan
\end{enumerate}
\end{multicols}
\end{sommaire}

\begin{objectifs}
\begin{itemize}
\item Lire, comprendre et commenter un poème en espagnol en identifiant son thème, sa structure et sa versification.
\item Reconnaître et expliquer l'effet de six figures de style : métaphore, comparaison, personification, hyperbole, allitération, anaphore.
\item Utiliser le lexique spécifique de la poésie (verso, estrofa, rima, poeta, musa) et du débat alimentaire (alimentos procesados, nutrientes, dieta equilibrada, comida chatarra).
\item Appliquer la méthode en quatre étapes du commentaire littéraire (contexte, lecture, analyse, synthèse).
\item Produire un commentaire écrit structuré du poème support.
\item Maîtriser les structures grammaticales de l'argumentation : subjonctif après expressions d'opinion, connecteurs logiques, comparatifs d'inégalité et d'égalité.
\item Traduire du français vers l'espagnol (thème) et de l'espagnol vers le français (version) des phrases sur le thème de l'alimentation.
\item Participer à un débat oral structuré en prenant position et en justifiant son point de vue avec du vocabulaire précis.
\end{itemize}
\end{objectifs}

\begin{prerequis}
\begin{itemize}
\item Connaissance des temps de l'indicatif (présent, passé composé, imparfait) et du subjonctif présent (formation régulière et irrégulière).
\item Notions de base sur la phrase complexe (subordonnées complétives, relatives, adverbiales).
\item Vocabulaire de base sur l'alimentation (frutas, verduras, carne, pescado, cocinar, comer) acquis en Première.
\item Capacité à identifier le sujet, le verbe et les compléments dans une phrase simple.
\item Expérience de la lecture à voix haute et de la prise de parole en continu courte.
\end{itemize}
\end{prerequis}

\newpage

\coursec{Découverte du poème support : « El duelo del plato »}

\courssub{Situation d'accroche et problématique}

Au Cameroun, comme dans tout le monde hispanophone, l'alimentation est au cœur des préoccupations quotidiennes : entre le \emph{ndolé} fait maison et les burgers des fast-foods de Douala ou de Yaoundé, le choix influence notre santé. Un poète espagnol contemporain, Juan García, illustre ce dilemme dans un poème où la salade et la frite s'affrontent en vers. Imaginez un instant : dans une assiette, deux camps se regardent en chiens de paella — d'un côté, la laitue croquante, fière de ses vitamines ; de l'autre, la frite dorée, irrésistible et coupable. Qui gagnera ce duel ?

\textbf{Problématique :} comment un poème peut-il transformer un sujet quotidien — manger sainement ou se faire plaisir — en un véritable combat poétique, et quelles armes formelles (vers, rimes, images) le poète utilise-t-il pour nous faire réfléchir ?

Pour répondre à cette question, nous allons d'abord découvrir le poème, le lire à voix haute, le comprendre globalement, puis en analyser la structure externe : strophes, vers, syllabes et rimes.

\courssub{Présentation de l'auteur et du contexte}

Juan García est un poète espagnol contemporain (personnage fictif créé pour cette leçon). Il appartient à une génération de poètes qui écrivent sur les problèmes de la société actuelle : l'environnement, la santé, les habitudes de consommation. Son style est simple, proche de la poésie populaire : des vers courts, des rimes régulières, un ton tantôt humoristique, tantôt critique.

\begin{savaistu}[La poésie engagée du quotidien]
En Espagne et en Amérique latine, de nombreux poètes contemporains abandonnent les grands thèmes abstraits (l'amour éternel, la mort) pour parler de la vie quotidienne : le supermarché, le fast-food, le sport, les réseaux sociaux. Cette poésie « du quotidien » utilise souvent l'humour et l'ironie pour critiquer nos habitudes. Au Cameroun, on pourrait imaginer un poème sur le duel entre le \emph{ndolé} traditionnel et la pizza livrée à domicile !
\end{savaistu}

\courssub{Lecture du poème : « El duelo del plato »}

L'enseignant lit d'abord le poème à voix haute, avec une intonation dramatique : c'est un \emph{duel}, un combat ! Les élèves écoutent, puis lisent à leur tour, en respectant les pauses (fin de vers, fin de strophe) et en marquant les rimes.

\begin{textebox}[« El duelo del plato », Juan García]
1\quad En mi plato hay una guerra\\
2\quad crujiente y dorada frita,\\
3\quad verde hoja de la tierra,\\
4\quad la papa que nunca grita.\\[6pt]
5\quad La lechuga da la vida\\
6\quad con vitaminas y salud;\\
7\quad hamburguesa prometida,\\
8\quad dulce engaño, no salud.\\[6pt]
9\quad Yo miro mi propio plato:\\
10\quad el tomate, rojo y fuerte,\\
11\quad es de mi salud retrato;\\
12\quad la frita pierde su suerte.\\[6pt]
13\quad Comer sano, comer sano,\\
14\quad es quererse a uno bien;\\
15\quad la verdura es mi hermano\\
16\quad y mi cuerpo dice también.
\end{textebox}

\textbf{Glossaire (mots difficiles) :}

\begin{center}
\begin{tabularx}{\linewidth}{|X|X|}\hline
\rowcolor{popblueL} \textbf{Español} & \textbf{Français} \\ \hline
\esp{crujiente} & croquant, croustillant \\ \hline
\esp{la papa} & la pomme de terre (Amérique latine ; en Espagne : \esp{la patata}) \\ \hline
\esp{la lechuga} & la laitue \\ \hline
\esp{la salud} & la santé \\ \hline
\esp{antojarse} & avoir très envie de, se faire désirer (aliment) \\ \hline
\esp{el engaño} & le piège, la tromperie \\ \hline
\esp{el retrato} & le portrait, le reflet \\ \hline
\esp{la suerte} & la chance \\ \hline
\esp{comer sano} & manger sainement \\ \hline
\esp{la verdura} & les légumes verts \\ \hline
\esp{engordar} & grossir, faire grossir \\ \hline
\esp{el nutriente} & le nutriment \\ \hline
\end{tabularx}
\end{center}

\textbf{Questions de première compréhension :}
\begin{enumerate}
\item ¿De qué habla el poema ? (De quoi parle le poème ?)
\item ¿Qué dos « personajes » se enfrentan ? (Quels deux « personnages » s'affrontent ?)
\item ¿Cuál es el tono : serio, humorístico, crítico ? (Quel est le ton ?)
\item ¿Qué consejo da el poeta al final ? (Quel conseil le poète donne-t-il à la fin ?)
\end{enumerate}

\courssub{Première compréhension globale}

\begin{exoresolu}[Compréhension globale du poème]
Énoncé : répondez en espagnol aux questions suivantes : (a) ¿Cuál es el tema del poema ? (b) ¿Qué tono utiliza el poeta ? (c) ¿Quién « gana » el duelo ?
\tcblower
Solution détaillée :\\
(a) \esp{El tema es la lucha entre la comida sana (ensalada, verduras) y la comida basura (fritas, hamburguesa).} « Le thème est la lutte entre l'alimentation saine et la malbouffe. »\\
(b) \esp{El tono es humorístico y crítico: el poeta convierte los alimentos en guerreros, pero quiere hacernos reflexionar sobre nuestra salud.} « Le ton est humoristique et critique : le poète transforme les aliments en guerriers, mais veut nous faire réfléchir sur notre santé. »\\
(c) \esp{Gana la comida sana: « la frita pierde su suerte » y el poeta aconseja « comer sano ».} « C'est l'alimentation saine qui gagne : la frite perd sa chance et le poète conseille de manger sainement. »
\end{exoresolu}

\courssub{Le lexique de la poésie}

Pour parler d'un poème, il faut d'abord maîtriser son vocabulaire technique.

\begin{definition}[Les mots-clés de la poésie]
\begin{itemize}
\item \cle{Verso} (le vers) : chaque ligne d'un poème. Exemple : \esp{« En mi plato hay una guerra »} est le premier \esp{verso} du poème.
\item \cle{Estrofa} (la strophe) : un groupe de vers séparé des autres par un espace. Notre poème a quatre \esp{estrofas} de quatre vers.
\item \cle{Rima} (la rime) : la répétition de sons identiques à la fin des vers. Exemple : \esp{guerra / tierra}.
\item \cle{Poeta} (le poète) : la personne qui écrit le poème. Ici, Juan García.
\item \cle{Musa} (la muse) : la source d'inspiration du poète. Ici, sa « muse » est... l'assiette quotidienne !
\end{itemize}
\end{definition}

\begin{attention}[Ne pas confondre \esp{verso} et \esp{estrofa}]
Les francophones confondent souvent ces deux mots. \esp{Un verso} = UNE ligne ; \esp{una estrofa} = un GROUPE de lignes. On ne dit jamais « le poème a seize strophes » pour un poème de seize vers : ici, le poème a \esp{dieciséis versos} répartis en \esp{cuatro estrofas}. Autre piège : \esp{rima} peut être \esp{consonante} (voyelles ET consonnes identiques : \esp{guerra/tierra}) ou \esp{asonante} (seules les voyelles identiques : \esp{guerra/fiesta}, a-e).
\end{attention}

\courssub{Analyse de la structure externe}

Comptons maintenant : le poème contient \textbf{4 strophes} de \textbf{4 vers} chacune (une strophe de quatre vers s'appelle \esp{cuarteto}), soit \textbf{16 vers} au total.

\begin{propriete}[La métrique espagnole : compter les syllabes]
Pour compter les syllabes d'un vers espagnol (\esp{medir el verso}), on applique trois règles :
\begin{enumerate}
\item \textbf{La sinalefa} (\esp{sinalefa}) : quand un mot se termine par une voyelle et que le suivant commence par une voyelle, les deux voyelles fusionnent en une seule syllabe. Exemple : \esp{En-mi-pla-to-hay\_u-na-gue-rra} : « hay una » se prononce \esp{ha-yu-na} (2 syllabes au lieu de 3).
\item \textbf{Mot aigu final} (\esp{palabra aguda}) : si le dernier mot est accentué sur la dernière syllabe (\esp{salud, corazón}), on AJOUTE une syllabe au compte. \esp{con vitaminas y salud} : 7 syllabes + 1 = 8.
\item \textbf{Mot esdrújulo final} (\esp{palabra esdrújula}) : si le dernier mot est accentué sur l'antépénultième syllabe (\esp{médico}), on RETRANCHE une syllabe.
\end{enumerate}
Un mot accentué sur l'avant-dernière syllabe (\esp{palabra llana} : \esp{plato, guerra}) ne change rien au compte.
\end{propriete}

Vérifions avec le premier vers : \esp{En mi plato hay una guerra} = En(1)-mi(2)-pla(3)-to(4)-hay\_u(5)-na(6)-gue(7)-rra(8). Le mot final \esp{guerra} est \esp{llano} : le compte reste 8. C'est un \cle{verso octosílabo} (vers de 8 syllabes), comme tous les vers du poème.

\begin{savaistu}[Le vers de huit syllabes, roi de la poésie populaire]
Le \esp{verso octosílabo} est le vers le plus fréquent de la poésie populaire espagnole : on le trouve dans les \esp{romances} médiévaux, les \esp{coplas} et les chansons traditionnelles. Facile à mémoriser et à chanter, il ressemble au rythme naturel de la parole. Les grands poètes comme Federico García Lorca l'ont souvent utilisé.
\end{savaistu}

\courssub{Les rimes : consonantes et schéma ABAB}

Observons les derniers mots de chaque vers : \esp{guerra / frita / tierra / grita}. Les sons finaux \esp{-erra} et \esp{-ita} se répètent exactement (voyelles et consonnes identiques) : ce sont des \cle{rimas consonantes} (rimes consonantiques, appelées « rimes riches » en français). Le schéma est \textbf{ABAB} : on dit que les rimes sont \esp{cruzadas} (croisées).

\begin{center}
\begin{tabularx}{\linewidth}{|X|X|X|}\hline
\rowcolor{popblueL} \textbf{Français} & \textbf{Español} & \textbf{Exemple du poème} \\ \hline
Rime consonantique (riche) & \esp{rima consonante} & \esp{guerra / tierra} \\ \hline
Rime assonantique (pauvre) & \esp{rima asonante} & \esp{vida / prometida} (voyelles i-a) \\ \hline
Rimes croisées & \esp{rimas cruzadas} (ABAB) & strophes 1 à 4 \\ \hline
Rimes embrassées & \esp{rimas abrazadas} (ABBA) & — \\ \hline
Rimes suivies & \esp{rimas pareadas} (AABB) & — \\ \hline
\end{tabularx}
\end{center}

\begin{popfigure}
\begin{tikzpicture}[
  verso/.style={rectangle, rounded corners=2pt, draw=popblue, fill=popblueL, minimum height=0.62cm, font=\small, align=center},
  rimaA/.style={circle, draw=poporange, fill=poporangeL, inner sep=1.5pt, font=\footnotesize\bfseries},
  rimaB/.style={circle, draw=popgreen, fill=popgreenL, inner sep=1.5pt, font=\footnotesize\bfseries},
  estrofa/.style={draw=poppurple, dashed, rounded corners=6pt, inner sep=6pt},
  lbl/.style={font=\footnotesize\itshape, text=popmuted}
]
% Strophe 1
\node[verso] (v1) at (0,0) {En mi plato hay una guerra};
\node[verso] (v2) at (0,-0.8) {crujiente y dorada frita};
\node[verso] (v3) at (0,-1.6) {verde hoja de la tierra};
\node[verso] (v4) at (0,-2.4) {la papa que nunca grita};
\node[rimaA] (a1) at (4.6,0) {A};
\node[rimaB] (b1) at (4.6,-0.8) {B};
\node[rimaA] (a2) at (4.6,-1.6) {A};
\node[rimaB] (b2) at (4.6,-2.4) {B};
\draw[->, poporange, thick] (a1) -- (a2);
\draw[->, popgreen, thick] (b1) -- (b2);
\begin{scope}[on background layer]
\node[estrofa, fit=(v1)(v2)(v3)(v4)(a1)(b2), label={[lbl]above:Estrofa 1 : cuarteto}] {};
\end{scope}
% Flèche vers résumé
\node[rectangle, draw=popdark, fill=popcream, rounded corners=4pt, align=center, font=\small] (res) at (2.3,-3.9) {4 estrofas $\times$ 4 versos\\ versos octosílabos (8 sílabas)\\ rima consonante ABAB (cruzada)};
\draw[->, popdark, thick] (2.3,-2.95) -- (res.north);
\end{tikzpicture}
\legende{Structure externe du poème : 4 strophes de 4 vers octosyllabes, rimes croisées ABAB.}
\end{popfigure}

\begin{exoresolu}[Analyse formelle d'une strophe]
Énoncé : analysez la strophe 3 du poème : nombre de vers, mesure des vers 9 et 12, type de rime et schéma.
\tcblower
Solution détaillée :\\
1. Nombre de vers : la strophe 3 contient \esp{cuatro versos} (vers 9 à 12) : c'est un \esp{cuarteto}.\\
2. Mesure du vers 9 : \esp{Yo mi-ro mi pro-pio pla-to} = 8 syllabes ; \esp{plato} est \esp{llano}, donc le compte reste 8 : \esp{verso octosílabo}.\\
3. Mesure du vers 12 : \esp{la fri-ta pier-de su suer-te} = 8 syllabes, mot final \esp{llano} : \esp{verso octosílabo}.\\
4. Rimes : \esp{plato / retrato} (A) et \esp{fuerte / suerte} (B) : voyelles et consonnes identiques (\esp{-ato}, \esp{-erte}), donc \esp{rima consonante}.\\
5. Schéma : ABAB, c'est-à-dire des \esp{rimas cruzadas} (rimes croisées).
\end{exoresolu}

\begin{methode}[Première approche d'un poème en quatre étapes]
\begin{enumerate}
\item \textbf{Contexte et auteur} : qui écrit ? quand ? sur quel sujet ? (Ici : Juan García, poète contemporain, thème de l'alimentation.)
\item \textbf{Lecture à voix haute} : lire plusieurs fois, marquer les pauses, écouter les rimes et le rythme.
\item \textbf{Compréhension globale} : identifier le sujet (\esp{el tema}), le ton (\esp{el tono} : humorístico, crítico, reflexivo) et l'idée principale.
\item \textbf{Analyse formelle} : compter les strophes et les vers, mesurer les syllabes (attention à la \esp{sinalefa} et aux mots \esp{agudos}), relever les rimes et établir le schéma (ABAB, ABBA...).
\end{enumerate}
\end{methode}

\begin{exercice}[À vous de jouer]
1. Comptez les syllabes des vers 5 et 6 du poème en appliquant les règles de la métrique.\\
2. Les rimes \esp{sano / hermano} (vers 13 et 15) sont-elles consonantes ou asonantes ? Justifiez.\\
3. Relevez dans le poème deux mots \esp{agudos} et deux mots \esp{llanos}.
\end{exercice}

\begin{corrige}[Exercice : métrique et rimes]
1. Vers 5 : \esp{La le-chu-ga da la vi-da} = 8 syllabes, mot final \esp{llano} : octosyllabe. Vers 6 : \esp{con vi-ta-mi-nas y sa-lud} = 7 syllabes, mais \esp{salud} est \esp{agudo} : 7 + 1 = 8, octosyllabe.\\
2. \esp{sano / hermano} : les sons finaux sont \esp{-ano} dans les deux mots, voyelles et consonnes identiques : c'est une \esp{rima consonante}.\\
3. Mots \esp{agudos} : \esp{salud}, \esp{también}. Mots \esp{llanos} : \esp{plato}, \esp{guerra} (aussi \esp{frita, suerte, hermano}).
\end{corrige}

\begin{aretenir}[L'essentiel de cette étape]
\begin{itemize}
\item Le poème « El duelo del plato » met en scène un combat humoristique entre alimentation saine et malbouffe, sur un ton critique.
\item Structure externe : 4 \esp{estrofas} (\esp{cuartetos}) de 4 \esp{versos}, soit 16 vers.
\item Tous les vers sont \esp{octosílabos} (8 syllabes) : penser à la \esp{sinalefa} et à la règle du mot final (\esp{agudo} : +1 ; \esp{esdrújulo} : $-$1).
\item Les rimes sont \esp{consonantes} et suivent le schéma ABAB (\esp{rimas cruzadas}).
\end{itemize}
\end{aretenir}

\coursec{Lexique de la poésie et analyse structurelle détaillée}

Dans la section précédente, nous avons découvert le poème \emph{El duelo del plato}, ce combat symbolique entre les aliments sains et la malbouffe. Nous l'avons lu, écouté, compris globalement. Mais un poème ne se contente pas de « raconter » : il est une \cle{construction} savante, faite de vers mesurés, de rimes calculées, de strophes ordonnées. Pour le commenter comme un hispanophone cultivé, il faut d'abord posséder le \cle{vocabulaire technique} de la poésie, puis savoir \cle{découper} le texte en unités de sens et en unités de son. C'est l'objet de cette section : apprendre à nommer, mesurer et analyser.

\courssub{Le lexique de la poésie : nommer pour comprendre}

Observons d'abord deux vers de notre poème :

\esp{La lechuga verde susurra.} \esp{La patata frita ríe.}

« La laitue verte chuchote. » « La frite rit. »

Chacune de ces lignes est un \esp{verso} ; le groupe de quatre vers qui les contient forme une \esp{estrofa} ; la ressemblance sonore entre les fins de vers est une \esp{rima}. En espagnol comme en français, la poésie possède son propre vocabulaire, mais attention : certains mots se ressemblent sans se recouvrir exactement.

\begin{center}
\begin{tabularx}{\linewidth}{|l|l|X|}
\hline
\rowcolor{popblueL} \textbf{Español} & \textbf{Français} & \textbf{Exemple tiré du poème} \\
\hline
\esp{verso} & vers & \esp{La lechuga verde susurra} \\
\hline
\esp{estrofa} & strophe & le poème compte 4 strophes de 4 vers \\
\hline
\esp{rima consonante} & rime suffisante & \esp{canta / levanta} : voyelles et consonnes identiques \\
\hline
\esp{rima asonante} & assonance & \esp{susurra / lucha} : seules les voyelles (u-a) se répètent \\
\hline
\esp{métrica} & métrique & vers de 8 syllabes = \esp{octosílabo} \\
\hline
\esp{sinalefa} & synalèphe & \esp{la\_ensalada} : 2 voyelles fusionnent \\
\hline
\esp{césura} & césure & pause intérieure : \esp{La lechuga // verde susurra} \\
\hline
\esp{poeta / poetisa} & poète / poétesse & l'auteur du poème \\
\hline
\esp{musa} & muse & source d'inspiration du poète \\
\hline
\esp{lira} & lyre & symbole de la poésie lyrique \\
\hline
\esp{copla} & couplet, strophe populaire & strophe courte de la tradition orale \\
\hline
\esp{poema / poesía} & poème / poésie & \esp{poesía} = l'art ; \esp{poema} = l'œuvre concrète \\
\hline
\end{tabularx}
\end{center}

\begin{attention}[Faux amis et pièges de vocabulaire]
Ne confondez jamais \esp{poesía} et \esp{poema} : la \esp{poesía} est l'art poétique en général (« la poésie »), le \esp{poema} est un texte précis (« le poème »). De même, \esp{estrofa} se dit « strophe » et non « estrophe » en français soigné. Enfin, \esp{rima} est féminin : \esp{la rima}, comme « la rime ».
\end{attention}

\begin{savaistu}[La copla, mémoire du peuple]
La \esp{copla} est une strophe brève, souvent de quatre vers octosyllabes, née de la tradition orale espagnole et transmise par les chansons, les romances et les corridos latino-américains. Au Cameroun, on peut la comparer aux refrains des chanteurs de makossa ou de bikutsi : courts, rythmés, faciles à mémoriser. Notre poème \emph{El duelo del plato} hérite directement de cette forme populaire.
\end{savaistu}

\begin{popfigure}
\begin{tikzpicture}[mindmap, grow cyclic, every node/.style={concept, text=white, align=center, font=\small}, concept color=popblue, level 1/.append style={level distance=3.2cm, sibling angle=90}]
\node[concept color=popdark, font=\small\bfseries, align=center]{La poesía\\(léxico poético)}
  child[concept color=poporange] { node {Forma}
    child { node[concept color=poporange!80] {verso} }
    child { node[concept color=poporange!80] {estrofa} }
    child { node[concept color=poporange!80] {copla} }
  }
  child[concept color=popgreen] { node {Sonido}
    child { node[concept color=popgreen!80, align=center]{rima\\consonante} }
    child { node[concept color=popgreen!80, align=center]{rima\\asonante} }
    child { node[concept color=popgreen!80] {aliteración} }
  }
  child[concept color=poppurple] { node {Estructura}
    child { node[concept color=poppurple!80] {métrica} }
    child { node[concept color=poppurple!80] {sinalefa} }
    child { node[concept color=poppurple!80] {césura} }
  }
  child[concept color=poppink] { node {Actores}
    child { node[concept color=poppink!80] {poeta} }
    child { node[concept color=poppink!80] {musa} }
    child { node[concept color=poppink!80] {lira} }
  };
\end{tikzpicture}
\legende{Carte mentale : les champs lexicaux de la poésie (forme, son, structure, acteurs).}
\end{popfigure}

\courssub{La versification : mesurer les vers}

Compter les syllabes d'un vers espagnol ne se fait pas comme en français : il faut appliquer deux règles spécifiques, la \cle{sinalefa} et la \cle{loi de l'accent final}.

\begin{definition}[La sinalefa]
La \esp{sinalefa} (synalèphe) est la fusion de deux voyelles en une seule syllabe métrique lorsqu'un mot se termine par une voyelle et que le mot suivant commence par une voyelle (ou un h muet). Exemple tiré du poème : \esp{la\_ensalada} se prononce métriquement \emph{laen-sa-la-da} = 4 syllabes métriques au lieu de 5 phonétiques. La synalèphe peut enchaîner trois voyelles : \esp{que\_a\_otros} = \emph{queao-tros}.
\end{definition}

\begin{propriete}[La loi de l'accent final du vers]
En espagnol, la dernière syllabe métrique du vers dépend de l'accentuation du dernier mot :
\begin{itemize}
\item mot \esp{agudo} (accentué sur la dernière syllabe : \esp{ríe}, \esp{comió}) : on \textbf{ajoute 1} syllabe au compte ;
\item mot \esp{llano} (accentué sur l'avant-dernière : \esp{susurra}, \esp{plato}) : le compte reste \textbf{inchangé} ;
\item mot \esp{esdrújulo} (accentué sur la troisième syllabe en partant de la fin : \esp{rápido}) : on \textbf{retire 1} syllabe.
\end{itemize}
Application sur deux vers du poème :
\begin{itemize}
\item \esp{La lechuga verde susurra} : La-le-chu-ga-ver-de-su-sur-ra = 8 syllabes ; dernier mot \esp{susurra} (llano) $\rightarrow$ 8 syllabes métriques. C'est un \esp{octosílabo}.
\item \esp{La patata frita ríe} : La-pa-ta-ta-fri-ta-ríe = 7 syllabes phonétiques ; dernier mot \esp{ríe} (agudo, hiatus í-e = 2 syllabes) $\rightarrow$ 7 + 1 = 8 syllabes métriques. C'est un \esp{octosílabo} agudo.
\end{itemize}
\end{propriete}

\begin{exemplebox}[Deux vers analysés et traduits]
\esp{La lechuga verde susurra} : 8 syllabes, vers llano, césure après la 4\textsuperscript{e} syllabe (\esp{La lechuga // verde susurra}). Traduction : « La laitue verte chuchote ».

\esp{La patata frita ríe} : 7 syllabes phonétiques, mais vers agudo (\esp{ríe}) donc 8 syllabes métriques. Traduction : « La frite rit ».

Remarquez que le français ne connaît ni la synalèphe systématique ni la loi de l'accent final : le vers français se compte simplement en syllabes prononcées, avec le e muet qui joue un rôle particulier. C'est une différence majeure entre les deux systèmes métriques.
\end{exemplebox}

\begin{attention}[L'erreur classique du francophone]
L'erreur la plus fréquente consiste à compter les \cle{syllabes phonétiques} au lieu des \cle{syllabes métriques}, c'est-à-dire à ignorer la sinalefa. Ainsi \esp{la ensalada verde} ne fait pas 7 syllabes mais 6 (la\_en-sa-la-da-ver-de). Deuxième piège : oublier d'ajouter une syllabe quand le vers se termine par un mot \esp{agudo}. Entraînez-vous toujours en marquant d'abord les synalèphes d'un trait de liaison, puis en vérifiant l'accent du dernier mot.
\end{attention}

\begin{exoresolu}[Mesurer un vers]
Énoncé : mesurez le vers \esp{El aceite quema mi garganta} (« L'huile brûle ma gorge ») et indiquez s'il est agudo, llano ou esdrújulo.
\tcblower
Solution détaillée :
\begin{enumerate}
\item Marquage des synalèphes : \esp{El\_aceite} fusionne (El\_a-cei-te) ; \esp{mi\_garganta} fusionne (mi\_gar-gan-ta).
\item Comptage : El\_a-cei-te-que-ma-mi\_gar-gan-ta = 10 syllabes.
\item Dernier mot : \esp{garganta}, accentué sur l'avant-dernière syllabe (gar-GAN-ta) = mot llano $\rightarrow$ compte inchangé.
\item Conclusion : vers de 10 syllabes métriques (\esp{decasílabo}), vers llano.
\end{enumerate}
\end{exoresolu}

\courssub{Les rimes : musique et mémoire}

Reprenons les quatre vers de la strophe 2 du poème :

\esp{Tomate rojo canta,} \\
\esp{La hamburguesa llora,} \\
\esp{Brócoli firme levanta,} \\
\esp{La pizza grasa se dora.}

Les vers 1 et 3 riment en \esp{-anta} (\esp{canta / levanta}) : voyelles \emph{et} consonnes identiques après la dernière voyelle accentuée, c'est une \esp{rima consonante}. Les vers 2 et 4 (\esp{llora / dora}) riment aussi de façon consonante en \esp{-ora}. Le schéma est donc ABAB, appelé \esp{rima cruzada} (rimes croisées). En revanche, si seules les voyelles se répètent, comme dans \esp{susurra / lucha} (u-a / u-a), on parle de \esp{rima asonante}.

\begin{aretenir}[À quoi servent les rimes ?]
\begin{itemize}
\item \cle{Renforcement du sens} : les mots rimés se répondent ; ici \esp{canta / levanta} associent les légumes à la vie et à l'action victorieuse, tandis que \esp{llora / dora} associent la malbouffe à la plainte et à la fausse apparence.
\item \cle{Musicalité} : la rime crée un balancement régulier qui rapproche le poème de la chanson.
\item \cle{Mémorisation} : un texte rimé se retient facilement, c'est pourquoi la tradition orale (coplas, comptines, slogans publicitaires) l'utilise partout dans le monde hispanophone.
\end{itemize}
\end{aretenir}

\courssub{La structure interne : la progression thématique}

Un commentaire de poème exige de dégager la \cle{progression des idées}. \emph{El duelo del plato} est construit en quatre mouvements, comme un petit drame :

\begin{itemize}
\item \textbf{Strophe 1 — présentation} : les aliments entrent en scène ; le poète dresse le portrait des deux camps (légumes frais : \esp{la lechuga verde susurra, el tomate rojo brilla} ; malbouffe : \esp{la patata frita ríe, y la lechuga lucha}). Le cadre est le marché ou la cuisine, espace familier de l'élève camerounais comme de l'hispanophone.
\item \textbf{Strophe 2 — confrontation} : le « duel » proprement dit ; les verbes d'action (\esp{canta, llora, levanta, dora}) animent la bataille symbolique entre le sain et le nuisible.
\item \textbf{Strophe 3 — conséquences} : le corps paie les dégâts de la malbouffe (\esp{la garganta}, \esp{el corazón cansado}) tandis que les aliments sains redonnent des forces.
\item \textbf{Strophe 4 — conclusion et réflexion} : le poète s'adresse au lecteur et transforme le récit en leçon de vie : choisir son assiette, c'est choisir sa santé.
\end{itemize}

Cette progression suit le schéma classique \emph{exposition — nœud — dénouement — morale}, hérité de la fable. Le poème n'est donc pas une simple liste d'aliments : c'est un \cle{récit argumenté} qui prépare le débat de la fin de la leçon.

\begin{methode}[Analyser la structure d'un poème en quatre étapes]
\begin{enumerate}
\item \textbf{Compter} : nombre de strophes, nombre de vers par strophe (ici : 4 strophes de 4 vers = quatrains, \esp{cuartetos}).
\item \textbf{Mesurer} : comptez les syllabes métriques de chaque vers (sinalefa + accent final) et repérez les césures.
\item \textbf{Noter le schéma de rimes} : attribuez une lettre à chaque rime (ABAB, ABBA, AABB...) et précisez consonante ou asonante.
\item \textbf{Observer l'enchaînement des idées} : résumez chaque strophe en une phrase et montrez comment elles s'enchaînent (présentation $\rightarrow$ confrontation $\rightarrow$ conséquences $\rightarrow$ réflexion).
\end{enumerate}
\end{methode}

\courssub{Texte d'application : un regard de spécialiste}

\begin{textebox}[Cómo leer un poema, mini-ensayo original]
Leer un poema es como escuchar una canción con los ojos. Antes de buscar el sentido, el buen lector observa la forma: cuenta los versos, mide las sílabas, escucha las rimas. En la poesía española, el octosílabo es el verso del pueblo: lo usan las coplas, los romances y las canciones infantiles. Su ritmo corto y vivo se parece al latido del corazón.

La sinalefa es un fenómeno esencial: une las palabras como une el río a sus afluentes. Gracias a ella, el verso fluye sin interrupciones y la lengua se vuelve música. La césura, en cambio, introduce una pausa, un momento de silencio que permite respirar y pensar.

Las rimas no son un adorno: refuerzan el significado. Cuando dos palabras riman, dialogan entre sí, se iluminan mutuamente. En un poema sobre la comida, por ejemplo, si «ensalada» rima con «vida» y «grasa» rima con «desgracia», el poeta nos está diciendo algo sin decirlo directamente.

Finalmente, toda estructura poética esconde una intención. Un poema que avanza de la presentación al conflicto y del conflicto a la reflexión imita el camino del pensamiento humano. Por eso, analizar la forma no es un ejercicio mecánico: es descubrir cómo el poeta convierte el lenguaje en experiencia. El verso medido, la rima buscada, la estrofa ordenada: todo conspira para que el mensaje llegue al corazón del lector y se quede en su memoria para siempre.
\end{textebox}

\textbf{Glossaire} : \esp{el latido} = le battement ; \esp{el afluente} = l'affluent ; \esp{un adorno} = un ornement ; \esp{mutuamente} = mutuellement ; \esp{la desgracia} = le malheur ; \esp{esconder} = cacher ; \esp{conspirar} = conspirer, concourir.

\textbf{Questions de compréhension} :
\begin{enumerate}
\item ¿Qué debe observar el lector antes de buscar el sentido de un poema?
\item ¿Por qué el octosílabo es « el verso del pueblo »?
\item Según el texto, ¿qué efecto produce la césura?
\item ¿Cómo pueden las rimas reforzar el mensaje de un poema sobre la comida?
\end{enumerate}

\courssub{Entraînement : identifier dans le poème}

\begin{exercice}[Chasse aux formes]
Dans les vers étudiés ci-dessus (\emph{El duelo del plato}), retrouvez :
\begin{enumerate}
\item un vers octosyllabe (8 syllabes métriques) et précisez s'il est llano ou agudo ;
\item une rime asonante entre deux vers ;
\item un vers dont la césure tombe après la 4\textsuperscript{e} syllabe ;
\item un vers terminé par un mot agudo, et donnez son compte métrique exact.
\end{enumerate}
\end{exercice}

\begin{corrige}[Exercice « Chasse aux formes »]
\begin{enumerate}
\item Octosyllabe llano : \esp{La lechuga verde susurra} (8 syllabes, mot final \esp{susurra} llano, compte inchangé). Octosyllabe agudo : \esp{La patata frita ríe} (7 syllabes phonétiques + 1 pour agudo = 8 métriques).
\item Rime asonante : \esp{susurra / lucha} (strophe 1) : seules les voyelles u-a se répètent, les consonnes diffèrent (ss/r vs ch).
\item Césure après la 4\textsuperscript{e} syllabe : \esp{La lechuga // verde susurra} (La-le-chu-ga = 4 syllabes avant la pause).
\item Vers agudo : \esp{La patata frita ríe} : \esp{ríe} est agudo (accent sur le \emph{i}), compte phonétique 7, compte métrique 8.
\end{enumerate}
\end{corrige}

Nous savons désormais nommer les éléments du poème, mesurer ses vers, cartographier ses rimes et suivre sa progression thématique. Il nous reste à observer de plus près les \cle{images} qui donnent vie à ce duel — métaphores, comparaisons, personnifications : ce sera l'objet de la section suivante, consacrée aux figures de style.

\coursec{Figures de style : identification, définition et effets}

Dans la section précédente, nous avons appris à décrire la \emph{structure} du poème « El duelo del plato » : ses vers, ses strophes, ses rimes. Mais un poème ne se résume pas à son squelette. Ce qui fait vibrer le lecteur, ce qui transforme une simple salade en « forêt verte » et un hamburger en piège doré, ce sont les \cle{figuras retóricas} (ou \cle{figuras de estilo}) : des procédés de langue qui s'écartent volontairement de l'expression ordinaire pour créer une image, une émotion ou une musicalité. C'est précisément ce que l'examinateur attend de vous au commentaire de texte : non seulement \emph{repérer} la figure, mais surtout \emph{expliquer son effet}. Observons d'abord, nommons ensuite, analysons enfin.

\courssub{Observation : les figures à l'œuvre dans le poème}

Relisons quelques vers clés du poème support et posons-nous une question simple : le poète parle-t-il « normalement » ?

\begin{exemplebox}[Extraits commentés de « El duelo del plato »]
\esp{La ensalada es un bosque verde} (v. 3) — « La salade est une forêt verte » : une salade n'est \emph{littéralement} pas une forêt. Le poète identifie deux réalités.\\
\esp{La lechuga crujiente susurra} (v. 5) — « La laitue croquante chuchote » : une laitue ne parle pas. On lui prête un comportement humain.\\
\esp{La hamburguesa brilla como el oro} (v. 6) — « Le hamburger brille comme l'or » : le mot \esp{como} signale un rapprochement explicite.\\
\esp{Mil calorías gritan en mi interior} (v. 10) — « Mille calories hurlent en moi » : exagération évidente.\\
\esp{Cruje, cruje la corteza} (v. 12) — « Ça croustille, ça croustille, la croûte » : répétition du son [kr].\\
\esp{Como la fruta de la tierra, / Como el pan de cada día} (v. 13-14) — « Comme le fruit de la terre, / Comme le pain de chaque jour » : même mot en tête de deux vers.
\end{exemplebox}

Chaque vers contient un \emph{écart} par rapport au langage ordinaire. À chaque type d'écart correspond une figure. Voici le tableau récapitulatif à recopier et à apprendre.

\begin{popfigure}
\begin{tikzpicture}[
  box/.style={draw=popline, rounded corners=2pt, fill=popcream, text width=3.1cm, align=center, font=\small, inner sep=4pt},
  head/.style={draw=popblue, fill=popblueL, text width=3.1cm, align=center, font=\small\bfseries, inner sep=4pt}
]
\node[head] at (0,0) {Figura};
\node[head] at (3.4,0) {Definición};
\node[head] at (6.8,0) {Ejemplo del poema};
\node[head] at (10.2,0) {Efecto / Función};
\node[box] at (0,-1.1) {\textbf{metáfora}};
\node[box] at (3.4,-1.1) {A es B, sin palabra-enlace};
\node[box] at (6.8,-1.1) {«La ensalada es un bosque verde»};
\node[box] at (10.2,-1.1) {imagen fuerte, identificación};
\node[box] at (0,-2.3) {\textbf{comparación}};
\node[box] at (3.4,-2.3) {A es \emph{como} B};
\node[box] at (6.8,-2.3) {«brilla como el oro»};
\node[box] at (10.2,-2.3) {imagen matizada, explícita};
\node[box] at (0,-3.5) {\textbf{personificación}};
\node[box] at (3.4,-3.5) {lo inanimado actúa como humano};
\node[box] at (6.8,-3.5) {«la lechuga susurra»};
\node[box] at (10.2,-3.5) {vida, cercanía afectiva};
\node[box] at (0,-4.7) {\textbf{hipérbole}};
\node[box] at (3.4,-4.7) {exageración voluntaria};
\node[box] at (6.8,-4.7) {«mil calorías gritan»};
\node[box] at (10.2,-4.7) {emoción intensa, dramatización};
\node[box] at (0,-5.9) {\textbf{aliteración}};
\node[box] at (3.4,-5.9) {repetición de sonidos consonánticos};
\node[box] at (6.8,-5.9) {«\textbf{cr}uje, \textbf{cr}uje la \textbf{c}orteza»};
\node[box] at (10.2,-5.9) {musicalidad, onomatopeya};
\node[box] at (0,-7.1) {\textbf{anáfora}};
\node[box] at (3.4,-7.1) {repetición al inicio del verso};
\node[box] at (6.8,-7.1) {«\textbf{Como}... / \textbf{Como}...»};
\node[box] at (10.2,-7.1) {insistencia, ritmo, martilleo};
\end{tikzpicture}
\legende{Tableau récapitulatif des six figures de style au programme de Terminale A.}
\end{popfigure}

\courssub{Définitions rigoureuses des six figures}

\begin{definition}[Les six figures de style au programme]
\begin{itemize}
\item \cle{Metáfora} : identification directe de deux réalités \emph{sans mot-outil} (schéma : A es B). \esp{Tus ojos son luceros.} « Tes yeux sont des astres. »
\item \cle{Comparación} (ou \esp{símil}) : rapprochement de deux réalités \emph{avec un mot-outil} : \esp{como}, \esp{parecido a}, \esp{igual que}, \esp{semejante a}, \esp{cual}. \esp{Tus ojos brillan como luceros.} « Tes yeux brillent comme des astres. »
\item \cle{Personificación} (ou \esp{prosopopeya}) : attribution de traits, d'actions ou de sentiments humains à un être inanimé ou abstrait. \esp{La ciudad duerme.} « La ville dort. »
\item \cle{Hipérbole} : exagération volontaire et visible pour amplifier une réalité. \esp{Te he llamado mil veces.} « Je t'ai appelé mille fois. »
\item \cle{Aliteración} : répétition d'un ou plusieurs sons consonantiques dans un même vers ou groupe de mots. \esp{El ruido con que rueda la ronca tempestad} (Bécquer, domaine public).
\item \cle{Anáfora} : répétition d'un mot ou d'un groupe de mots \emph{en début} de vers ou de phrases successives. \esp{Como la fruta..., / Como el pan...}
\end{itemize}
\end{definition}

\begin{propriete}[Effets propres des figures sonores et rythmiques]
L'\esp{aliteración} crée un \textbf{effet sonore} : musicalité du vers, imitation d'un bruit (onomatopée). Dans \esp{Cruje, cruje la corteza}, la répétition du groupe [kr] imite le bruit du croquant : le vers \emph{fait entendre} ce qu'il décrit.\\
L'\esp{anáfora} crée une \textbf{insistance rythmique et thématique} : le mot répété martèle l'idée, donne au poème un souffle d'incantation ou de prière, et soude les vers entre eux.
\end{propriete}

\courssub{Le piège capital : metáfora ou comparación ?}

C'est la confusion la plus fréquente chez les francophones, car le français distingue aussi « métaphore » et « comparaison », mais les élèves oublient d'appliquer le critère en espagnol.

\begin{attention}[Le critère du mot-outil]
Une seule question suffit : \textbf{y a-t-il un mot-outil} (\esp{como}, \esp{parece}, \esp{parecido a}, \esp{igual que}, \esp{cual}) ?
\begin{itemize}
\item \esp{Tus ojos son luceros.} → \textbf{metáfora} (identification directe, plus audacieuse).
\item \esp{Tus ojos brillan como luceros.} → \textbf{comparación} (ressemblance explicite, plus prudente).
\end{itemize}
Autre paire : \esp{El examen fue un infierno} (métaphore : « l'épreuve fut un enfer ») / \esp{El examen fue como un infierno} (comparaison). La métaphore affirme ; la comparaison suggère. Dans le commentaire, dites-le : la métaphore produit une image plus \emph{condensée} et plus \emph{violente}.
\end{attention}

\begin{center}
\begin{tabularx}{\linewidth}{|X|X|X|}
\hline
\rowcolor{popblueL} Français & Español & Remarque \\
\hline
métaphore & metáfora & accent obligatoire sur le \emph{a} \\
\hline
comparaison & comparación / símil & \esp{símil} est le terme technique \\
\hline
personnification & personificación / prosopopeya & deux termes synonymes \\
\hline
hyperbole & hipérbole & accent sur le \emph{e} \\
\hline
allitération & aliteración & un seul \emph{l} en espagnol \\
\hline
anaphore & anáfora & accent sur le premier \emph{a} \\
\hline
\end{tabularx}
\end{center}

\begin{attention}[Faux amis et orthographe]
\begin{itemize}
\item \esp{hipérbole} s'écrit avec un accent (mot \esp{esdrújulo}) ; ne pas écrire \emph{hiperbole}.
\item \esp{aliteración} : un seul \emph{l}, contrairement au français « allitération ».
\item \esp{figura} signifie « figure de style » mais aussi « silhouette » ; ne pas traduire \esp{una buena figura} par « une bonne figure de style » !
\item On dit \esp{una figura retórica} ou \esp{una figura de estilo}, jamais \emph{un estilo de figura}.
\end{itemize}
\end{attention}

\courssub{Recherche guidée : surligner et étiqueter le poème}

\begin{methode}[Identifier une figure de style en quatre pas]
\begin{enumerate}
\item \textbf{Repérer l'écart} : souligner ce qui n'est pas du langage ordinaire (une salade-forêt, une laitue qui parle, un mot répété).
\item \textbf{Identifier le procédé} : s'agit-il d'un rapprochement (avec ou sans \esp{como} ?), d'une attribution humaine, d'une exagération, d'une répétition de sons ou de mots ?
\item \textbf{Nommer la figure} avec le terme espagnol exact et son accent : \esp{metáfora}, \esp{comparación}, \esp{personificación}, \esp{hipérbole}, \esp{aliteración}, \esp{anáfora}.
\item \textbf{Expliquer l'effet dans le contexte} : jamais de réponse toute faite. Demandez-vous : pourquoi \emph{ici}, dans \emph{ce} poème sur l'alimentation ? Exemple : la métaphore du « bosque verde » valorise la salade, la rend vaste et vivante ; la comparaison du hamburger à l'or suggère une séduction trompeuse, un leurre brillant.
\end{enumerate}
\end{methode}

\begin{experience}[Étiquetage collectif du poème]
Par groupes de trois : photocopiez ou recopiez le poème « El duelo del plato ». Avec trois couleurs, surlignez : en bleu les rapprochements (métaphores et comparaisons), en vert les personnifications, en rouge les répétitions (aliteración, anáfora). Dans la marge, étiquetez chaque occurrence en espagnol : \esp{v. 3 : metáfora ; v. 6 : comparación...} Puis chaque groupe choisit une figure et rédige une phrase d'analyse complète selon le modèle : \esp{En el verso 5, la personificación de la lechuga que «susurra» da vida a los alimentos sanos y crea una atmósfera de intimidad con la naturaleza.} (« Au vers 5, la personnification de la laitue qui "chuchote" donne vie aux aliments sains et crée une atmosphère d'intimité avec la nature. »)
\end{experience}

\courssub{Entraînement : un poème d'application}

Avant de classer des exemples isolés, entraînons-nous sur un court poème original, à commenter avec la méthode en quatre pas.

\begin{textebox}[« La merienda del barrio » (poème original)]
\esp{La tarde cae como un telón cansado\\
sobre el mercado de Mvog-Mbi.\\
El vendedor de frutas es un rey sin corona:\\
sus mangos son soles pequeños,\\
sus piñas, tesoros de oro verde.\\
Frente a él, la fritura chisporrotea y canta,\\
canta y llama, llama y engaña:\\
«Ven, ven, soy crujiente y barata».\\
Mil manos corren hacia el aceite,\\
mientras la naranja espera, redonda y paciente,\\
como una luna llena de vitaminas.}
\end{textebox}

\textbf{Glossaire} : \esp{la merienda} = le goûter ; \esp{el telón} = le rideau (de théâtre) ; \esp{la fritura} = la friture ; \esp{chisporrotear} = crépiter ; \esp{engañar} = tromper ; \esp{crujiente} = croustillant ; \esp{el aceite} = l'huile ; \esp{redondo/a} = rond(e).

\textbf{Questions} : 1) Relevez deux métaphores et deux comparaisons. 2) Quelle figure domine dans les vers 6-7 et quel effet produit-elle ? 3) Trouvez une hyperbole. 4) Quel contraste le poème construit-il entre les fruits et la friture ?

\begin{exoresolu}[Classification d'exemples nouveaux]
Énoncé. Indiquez la figure de style employée dans chaque phrase et justifiez en une ligne :
a) \esp{El fútbol es una religión en mi barrio.} b) \esp{Mi hermano come como un león.} c) \esp{El viento acaricia las hojas del cacao.} d) \esp{He esperado el bendskin durante siglos.} e) \esp{La lenta lluvia llora en la ladera.} f) \esp{Todo por la patria, todo por la escuela, todo por el futuro.}
\tcblower
Solution détaillée.
a) \textbf{Metáfora} : identification directe \esp{fútbol = religión}, sans \esp{como}. Effet : sacralisation du football.
b) \textbf{Comparación} : présence du mot-outil \esp{como}. Effet : image de l'appétit féroce.
c) \textbf{Personificación} : le vent « caresse », action humaine. Effet : tendresse, nature vivante.
d) \textbf{Hipérbole} : « des siècles » pour un taxi-moto, exagération comique de l'attente.
e) \textbf{Aliteración} : répétition du son [l] (\esp{lenta, lluvia, llora, ladera}) ; musicalité douce et mélancolique. On peut aussi noter la personnification de la pluie qui « pleure ».
f) \textbf{Anáfora} : répétition de \esp{Todo por} en tête de trois groupes ; effet de martèlement solennel.
\end{exoresolu}

\begin{exercice}[À vous de jouer]
1. Classez : a) \esp{Sus palabras eran puñales.} b) \esp{La sopa humea como un volcán pequeño.} c) \esp{El reloj grita las horas.} d) \esp{Te quiero más que a mil estrellas.} e) \esp{Rápido corre el río por la roca.} f) \esp{Nada de grasa, nada de azúcar, nada de prisa.}
2. Transformez la comparaison \esp{La yuca es como una batería de energía} en métaphore, puis expliquez la différence d'effet.
3. Création collective : avec votre voisin, écrivez \textbf{deux vers originaux} sur le thème de l'alimentation saine, utilisant \textbf{au moins deux figures} différentes ; nommez-les sous vos vers. Exemple attendu : \esp{El aguacate es una esmeralda cremosa / que abraza mi cuerpo como una madre.} (métaphore + comparaison).
\end{exercice}

\begin{corrige}[Exercice « À vous de jouer »]
1. a) métaphore (\esp{palabras = puñales}, sans mot-outil) ; b) comparaison (\esp{como}) ; c) personnification (l'horloge « crie ») ; d) hyperbole (\esp{mil estrellas}) ; e) allitération du son [r] (\esp{rápido, corre, río, roca}) ; f) anaphore (\esp{Nada de} répété).
2. Métaphore : \esp{La yuca es una batería de energía.} La comparaison suggère une ressemblance prudente ; la métaphore affirme l'identité, l'image est plus forte et plus directe.
3. Corrigé libre : vérifier la présence de deux figures nommées correctement et l'exactitude de l'espagnol (accents, accords).
\end{corrige}

\begin{savaistu}[L'anaphore, arme des poètes engagés]
L'anaphore est la figure favorite de la poésie engagée : en répétant le même mot en tête de vers, le poète martèle sa revendication comme un slogan. Pablo Neruda l'utilise avec force (\esp{Yo no te pido...}), et la tradition orale africaine — griots, chants de lutte, prêches — repose sur le même procédé de reprise rythmique. Quand vous débattrez de l'alimentation saine, réutilisez-la : \esp{¡Basta de grasa, basta de azúcar, basta de comida basura!} — trois \esp{basta}, et votre auditoire est conquis.
\end{savaistu}

\begin{aretenir}[L'essentiel de la section]
Six figures, trois familles : les \textbf{images} (metáfora sans mot-outil, comparación avec \esp{como}/\esp{parece}, personificación), l'\textbf{exagération} (hipérbole) et les \textbf{répétitions} (aliteración = sons, anáfora = mot en tête de vers). Pour le commentaire : repérer l'écart, nommer la figure en espagnol correctement accentué, et surtout \emph{expliquer l'effet dans le contexte} du poème. Ne jamais confondre \esp{La ensalada es un bosque} (métaphore) et \esp{La ensalada es como un bosque} (comparaison).
\end{aretenir}

\coursec{Méthodologie du commentaire de poème et commentaire modèle}

Dans la section précédente, nous avons appris à repérer et à nommer les figures de style (métaphore, comparaison, personnification, hyperbole, allitération, anaphore) présentes dans le poème \esp{« El duelo del plato »}. Mais savoir identifier une figure ne suffit pas au baccalauréat : il faut encore organiser ces observations dans un devoir structuré, rédigé et argumenté. C'est précisément l'objet du \cle{comentario de poema}. Cette section vous donne une méthode complète, en quatre étapes, puis un commentaire modèle intégral que vous pourrez imiter.

\courssub{Qu'est-ce qu'un commentaire de poème ?}

\begin{definition}[El comentario de poema]
Le \cle{comentario de poema} est un devoir écrit qui propose une \textbf{lecture organisée et argumentée} d'un poème. Il ne s'agit ni de résumer le texte, ni de donner simplement son avis : il faut montrer \textbf{comment} le poème produit son sens et ses émotions, en s'appuyant sur des citations analysées. En espagnol, on parle aussi d'\esp{análisis literario} ou d'\esp{explicación de texto}.
\end{definition}

Observez la différence entre ces deux phrases d'élève :

\begin{exemplebox}[Résumé ou analyse ?]
\esp{En el verso 3, el poeta dice que el tomate es un guerrero rojo.} — « Au vers 3, le poète dit que la tomate est un guerrier rouge. » $\rightarrow$ simple \textbf{paraphrase}, sans valeur ajoutée.

\esp{La metáfora « el tomate, guerrero rojo » (v. 3) convierte la comida en un campo de batalla y anuncia el tono épico-burlesco de todo el poema.} — « La métaphore "la tomate, guerrier rouge" (v. 3) transforme le repas en champ de bataille et annonce le ton épique et burlesque de tout le poème. » $\rightarrow$ véritable \textbf{analyse} : la citation est nommée (métaphore), localisée (v. 3) et interprétée (effet produit).
\end{exemplebox}

\courssub{Les quatre étapes de la méthode}

\begin{methode}[Les 4 étapes du commentaire de poème]
\begin{enumerate}
\item \textbf{Introduction} : rédigez une phrase d'\textbf{accroche} (généralité sur la poésie ou sur le thème), la \textbf{fiche d'identité} du poème (titre, auteur, thème), la \textbf{problématique} (la question que le commentaire va résoudre) et l'\textbf{annonce du plan} (les deux ou trois axes).
\item \textbf{Développement} : rédigez deux ou trois \textbf{axes thématiques}, chacun construit de la même façon : idée directrice, citation entre guillemets, analyse de la citation (forme + sens + effet), mini-conclusion de l'axe.
\item \textbf{Conclusion} : faites le \textbf{bilan} (réponse claire à la problématique), puis proposez une \textbf{ouverture} vers un autre texte, un autre art ou un débat actuel (par exemple l'alimentation saine contre la malbouffe).
\item \textbf{Relecture} : vérifiez l'orthographe et les accents (á, é, í, ó, ú, ñ), la cohérence des temps (présent de narration), les accords et la logique du plan. Une relecture soigneuse peut sauver deux points.
\end{enumerate}
\end{methode}

\begin{popfigure}
\begin{center}
\begin{tikzpicture}[
node distance=0.55cm,
etape/.style={rectangle, rounded corners=4pt, draw=popblue, fill=popblueL, text width=4.6cm, align=center, font=\small\bfseries, minimum height=0.9cm},
detalle/.style={rectangle, rounded corners=4pt, draw=popmuted, fill=popcream, text width=5.6cm, align=center, font=\scriptsize},
flecha/.style={-{Stealth[length=3mm]}, thick, poporange}
]
\node[etape] (intro) {1. INTRODUCCIÓN};
\node[detalle, right=0.6cm of intro] (d1) {accroche + ficha de identidad + problemática + anuncio del plan};
\node[etape, below=of intro] (eje1) {2. EJE 1};
\node[etape, below=of eje1] (eje2) {EJE 2};
\node[etape, below=of eje2] (eje3) {EJE 3};
\node[detalle, right=0.6cm of eje2] (d2) {idea + cita «...» (v. n) + análisis (forma, sentido, efecto)};
\node[etape, below=of eje3] (concl) {3. CONCLUSIÓN};
\node[detalle, right=0.6cm of concl] (d3) {balance + respuesta a la problemática};
\node[etape, below=of concl, fill=poporangeL, draw=poporange] (apert) {APERTURA};
\node[detalle, right=0.6cm of apert] (d4) {otro texto, otro arte o un debate actual};
\node[etape, below=of apert, fill=popgreenL, draw=popgreen] (revis) {4. RELECTURA};
\draw[flecha] (intro) -- (eje1);
\draw[flecha] (eje1) -- (eje2);
\draw[flecha] (eje2) -- (eje3);
\draw[flecha] (eje3) -- (concl);
\draw[flecha] (concl) -- (apert);
\draw[flecha] (apert) -- (revis);
\end{tikzpicture}
\end{center}
\legende{Organigramme de la structure du commentaire de poème : de l'introduction à la relecture.}
\end{popfigure}

\courssub{Construire l'introduction type}

L'introduction est la première impression laissée au correcteur : elle doit être complète mais concise (cinq à huit lignes). Elle comporte quatre éléments obligatoires, dans cet ordre.

\begin{center}
\begin{tabularx}{\linewidth}{|l|X|X|}
\hline
\rowcolor{popblueL} \textbf{Élément} & \textbf{Fonction} & \textbf{Formules utiles en español} \\
\hline
Accroche & Situer le thème du poème & \esp{La poesía puede transformar lo cotidiano en algo extraordinario...} \\
\hline
Fiche d'identité & Titre, auteur, thème & \esp{En el poema «...», el poeta... presenta...} \\
\hline
Problématique & Question directrice & \esp{¿Cómo consigue el poeta convertir una comida en una batalla épica?} \\
\hline
Annonce du plan & Deux ou trois axes & \esp{Analizaremos primero..., después... y finalmente...} \\
\hline
\end{tabularx}
\end{center}

\begin{exemplebox}[Introduction type et sa traduction]
\esp{En el poema « El duelo del plato », Juan García transforma una comida en una batalla épica donde los alimentos se convierten en guerreros. ¿Cómo consigue el poeta mezclar lo cotidiano y lo heroico para criticar nuestra manera de comer? Analizaremos primero la poesía de lo cotidiano, después el combate visual y sonoro, y finalmente la dimensión crítica del poema.}

« Dans le poème "Le duel de l'assiette", Juan García transforme un repas en une bataille épique où les aliments deviennent des guerriers. Comment le poète parvient-il à mêler le quotidien et l'héroïque pour critiquer notre façon de manger ? Nous analyserons d'abord la poésie du quotidien, puis le combat visuel et sonore, et enfin la portée critique du poème. »
\end{exemplebox}

\begin{attention}[Pièges de l'introduction]
N'écrivez jamais \esp{« Este poema es muy bonito »} comme accroche : c'est un jugement vide. Évitez aussi le calque du français « dans ce poème » traduit par \esp{« en este poema »} si vous n'avez pas encore cité le titre : citez d'abord le titre complet entre guillemets. Attention enfin à l'accent : \esp{poema} ne prend pas d'accent, mais \esp{poesía}, \esp{crítica} et \esp{épico} en prennent un. Méfiez-vous enfin du faux ami \esp{portada}, qui signifie « couverture » (d'un livre, d'une revue) : pour dire « la portée » d'un texte, employez \esp{la dimensión}, \esp{el alcance} ou \esp{el mensaje}.
\end{attention}

\courssub{Organiser le développement par axes}

Le développement ne suit pas l'ordre des vers : il suit des \textbf{idées}. Chaque axe est un petit raisonnement autonome. Pour \esp{« El duelo del plato »}, on peut proposer trois axes :

\begin{itemize}
\item \textbf{Axe A — La poésie du quotidien} : le poème choisit des objets banals (une assiette, des tomates, du riz) et les élève au rang de sujet poétique.
\item \textbf{Axe B — Le combat visuel et sonore} : les couleurs, les bruits (allitérations) et les métaphores guerrières créent une véritable scène de bataille.
\item \textbf{Axe C — La portée critique} : derrière le jeu, le poème dénonce notre rapport à la nourriture (excès, gaspillage, malbouffe).
\end{itemize}

Chaque axe suit le schéma : \textbf{affirmation} $\rightarrow$ \textbf{citation} $\rightarrow$ \textbf{analyse} $\rightarrow$ \textbf{transition} vers l'axe suivant.

\begin{propriete}[Règles de citation dans un commentaire]
\begin{enumerate}
\item La citation se place entre les \textbf{guillemets espagnols} : \esp{« el tomate, guerrero rojo »}.
\item Elle est \textbf{intégrée} à la phrase, jamais jetée seule : \esp{El poeta presenta el tomate como un « guerrero rojo » que defiende su territorio.}
\item On indique le \textbf{vers} entre parenthèses : (v. 3) ; pour une strophe : (est. 2).
\item Toute citation est \textbf{immédiatement analysée} : nommez le procédé (métaphore, allitération...) et dites son effet. Une citation sans analyse est une citation perdue.
\item La citation doit être \textbf{exacte} : respectez les mots, les accents et la ponctuation du poème.
\end{enumerate}
\end{propriete}

\begin{attention}[L'erreur capitale : le résumé vers par vers]
Le piège classique du francophone consiste à expliquer le poème vers après vers (\esp{« en el verso 1..., en el verso 2... »}). C'est un \textbf{résumé}, pas un commentaire. Privilégiez une \textbf{analyse thématique transversale} : regroupez les vers qui traitent de la même idée, même s'ils sont éloignés dans le poème. Et n'oubliez jamais d'analyser la \textbf{forme} (versification, rimes, figures de style) \textbf{au service du sens} : une allitération n'est pas là pour décorer, elle crée un effet (bruit de bataille, rythme, tension).
\end{attention}

\courssub{Rédiger la conclusion et l'ouverture}

La conclusion se construit en trois temps :

\begin{enumerate}
\item \textbf{Synthèse} : rappelez brièvement ce que les axes ont montré, sans répéter mot à mot.
\item \textbf{Réponse à la problématique} : dites clairement comment le poème atteint son but.
\item \textbf{Ouverture} : élargissez vers un autre poème, un autre art (peinture, chanson) ou un débat actuel — par exemple : \esp{¿somos lo que comemos?} (« sommes-nous ce que nous mangeons ? »), qui annonce le débat sur l'alimentation saine et la malbouffe.
\end{enumerate}

\begin{exemplebox}[Formules de conclusion]
\esp{En conclusión, el poema convierte un acto cotidiano en un espectáculo épico gracias a las metáforas guerreras y a los juegos sonoros.} — « En conclusion, le poème transforme un acte quotidien en un spectacle épique grâce aux métaphores guerrières et aux jeux sonores. »

\esp{Esta lectura nos invita a reflexionar sobre nuestra alimentación actual, dominada por la comida basura.} — « Cette lecture nous invite à réfléchir sur notre alimentation actuelle, dominée par la malbouffe. »
\end{exemplebox}

\courssub{Commentaire modèle complet}

Voici un commentaire modèle d'environ 250 mots, rédigé selon la méthode. Les balises indiquent les parties : vous ne les écrivez pas dans votre copie, elles servent ici à visualiser la structure.

\begin{textebox}[Comentario modelo de « El duelo del plato »]
[Introduction] La poesía es capaz de convertir lo más cotidiano en una aventura extraordinaria. En el poema « El duelo del plato », Juan García presenta una simple comida que se transforma en una batalla épica. ¿Cómo mezcla el poeta lo cotidiano y lo heroico para criticar nuestra relación con la comida? Estudiaremos la poesía de lo cotidiano, el combate visual y sonoro, y la dimensión crítica del texto.

[Axe 1] En primer lugar, el poema eleva objetos humildes a la categoría de héroes. La metáfora « el tomate, guerrero rojo » (v. 3) transforma un alimento común en un soldado valiente. Esta personificación de los alimentos demuestra que la poesía no necesita temas grandiosos : la cocina basta.

[Axe 2] Además, el poeta construye un verdadero combate. Los colores violentos — el rojo del tomate, el blanco del arroz — crean un campo de batalla visual, mientras que la aliteración en « s » del verso 7 imita el ruido de la sartén. La hipérbole « un ejército de sabores » (v. 9) exagera la escena con humor épico.

[Axe 3] Finalmente, detrás del juego aparece una crítica : el « duelo » del plato refleja nuestra lucha diaria entre la comida sana y la comida basura. El poeta nos invita a elegir nuestro bando.

[Conclusión] Así, gracias a las metáforas guerreras, a los sonidos y a la exageración, el poema convierte un acto banal en un espectáculo con mensaje.

[Ouverture] Esta reflexión nos conecta con el debate actual : ¿somos lo que comemos?
\end{textebox}

\textbf{Glossaire et remarques de langue :}
\begin{itemize}
\item \esp{lo cotidiano} : le quotidien, ce qui est banal ; \esp{humilde} : humble, modeste.
\item \esp{la sartén} : la poêle ; \esp{el sabor} : la saveur ; \esp{el bando} : le camp.
\item \esp{detrás del juego} : derrière le jeu ; \esp{así} : ainsi.
\item \esp{la dimensión crítica} : la portée critique (et non \esp{portada}, qui signifie « couverture »).
\item Notez l'emploi du \textbf{présent de narration} (\esp{presenta}, \esp{transforma}, \esp{convierte}) : en espagnol comme en français, on commente un texte au présent.
\item Connecteurs utilisés : \esp{en primer lugar}, \esp{además}, \esp{finalmente}, \esp{así} : apprenez-les par cœur, ils structurent votre devoir.
\end{itemize}

\begin{savaistu}[Le barème du baccalauréat]
Au baccalauréat camerounais, le commentaire de texte littéraire est noté sur 10 points : environ \textbf{4 points} pour la compréhension et la pertinence de l'analyse, \textbf{3 points} pour l'organisation (introduction, plan, conclusion) et \textbf{3 points} pour la qualité de la langue espagnole (orthographe, accents, grammaire, richesse du vocabulaire). Une copie bien structurée avec une langue correcte vaut donc souvent plus qu'une copie brillante mais désorganisée.
\end{savaistu}

\courssub{Exercice résolu et auto-évaluation}

\begin{exoresolu}[Réparer un début de commentaire]
\textbf{Énoncé.} Un élève écrit : \esp{« En el verso 1 hay un plato. En el verso 2 hay un tomate. En el verso 3 el tomate es un guerrero. El poema es bonito. »} Transformez ces quatre phrases en un début d'axe d'analyse correct, avec citation localisée et analyse d'une figure de style.

\tcblower
\textbf{Solution détaillée.} Il faut regrouper les vers autour d'une idée (la transformation du quotidien), citer exactement et analyser : \esp{« Desde los primeros versos, el poeta transforma los elementos humildes de la comida en personajes heroicos : la metáfora « el tomate, guerrero rojo » (v. 3) convierte un simple alimento en un soldado, lo que anuncia el tono épico del poema. »} On remarque : 1) une idée directrice (transformer le quotidien) ; 2) la citation exacte entre guillemets espagnols, localisée (v. 3) ; 3) le nom du procédé (\esp{metáfora}) ; 4) l'effet produit (ton épique). Le jugement vide \esp{« el poema es bonito »} a disparu.
\end{exoresolu}

\begin{methode}[Grille d'auto-évaluation de votre commentaire]
Avant de rendre votre copie, vérifiez chaque critère et notez-vous :
\begin{enumerate}
\item \textbf{Structure} (3 pts) : introduction complète (accroche, fiche d'identité, problématique, plan) ; deux ou trois axes distincts ; conclusion avec bilan et ouverture.
\item \textbf{Pertinence} (4 pts) : chaque axe répond à la problématique ; aucune paraphrase vers par vers ; les figures de style sont nommées et leurs effets expliqués.
\item \textbf{Exemples} : au moins une citation exacte et localisée (v. n) par axe, toujours analysée.
\item \textbf{Langue} (3 pts) : accents corrects, présent de narration, connecteurs variés (\esp{en primer lugar}, \esp{además}, \esp{sin embargo}, \esp{en conclusión}), aucun calque du français.
\end{enumerate}
Si un critère manque, corrigez avant de rendre : la relecture est la quatrième étape de la méthode, pas une option.
\end{methode}

Vous disposez maintenant de tous les outils pour rédiger un commentaire de poème complet et bien noté. La section suivante vous fournira le lexique nécessaire pour alimenter l'ouverture de vos conclusions et participer au débat : alimentation saine contre malbouffe.

\coursec{Lexique du débat : alimentation saine vs malbouffe}

Après avoir appris à commenter un poème — à en dégager le thème, la structure et les figures de style —, nous gardons le même fil conducteur : l'assiette. Le poème \esp{El duelo del plato} opposait déjà deux mondes, celui des produits frais et celui des produits industriels. Pour transformer cette intuition poétique en \cle{argumentation} solide, il faut maintenant disposer d'un vocabulaire précis, nuancé et adapté au registre du débat. C'est l'objectif de cette section : construire votre arsenal lexical pour défendre une position sur l'alimentation, à l'oral comme à l'écrit.

\courssub{Observation : un texte pour faire émerger le vocabulaire}

Lisons d'abord ce court article de presse (texte original), rédigé dans le registre journalistique que l'on attend d'un élève de Terminale. Soulignez mentalement tous les mots liés à l'alimentation.

\begin{textebox}[La salud en la mochila (texto original)]
En los mercados de Yaoundé y de Douala, los puestos rebosan de alimentos frescos: plátanos, ñames, verduras de hoja verde, frutas tropicales y pescado. Sin embargo, cada vez más jóvenes prefieren comprar comida rápida cerca de los institutos: patatas fritas, refrescos azucarados y productos ultraprocesados llenos de aditivos, grasas saturadas y azúcares añadidos.

Los nutricionistas advierten de que esta tendencia puede provocar problemas de salud como la obesidad, la diabetes y la hipertensión. «No se trata de prohibir, sino de educar», explica una especialista. «Una dieta equilibrada, rica en proteínas magras, fibra y una buena hidratación, previene enfermedades y mejora el rendimiento escolar.»

Desde mi punto de vista, el debate no es nuevo, pero sí urgente. Estoy a favor de que los comedores escolares ofrezcan menús equilibrados y me opongo a la publicidad de comida basura dirigida a los niños. Por tanto, es necesario actuar en la escuela, en la familia y en los medios de comunicación. Además, conviene recordar que la malnutrición no significa solo falta de alimentos: también existe la malnutrición por exceso de productos de mala calidad. En conclusión, comer sano no es un lujo: es un derecho y una responsabilidad compartida.
\end{textebox}

\textbf{Glossaire :} \esp{los puestos} = les étals ; \esp{rebosar de} = déborder de ; \esp{el ñame} = l'igname ; \esp{el refresco azucarado} = la boisson sucrée ; \esp{el rendimiento escolar} = les résultats scolaires ; \esp{el comedor escolar} = la cantine scolaire ; \esp{no se trata de} = il ne s'agit pas de.

\textbf{Questions de compréhension :}
\begin{enumerate}
\item ¿Qué alimentos frescos se mencionan en el texto?
\item ¿Qué problemas de salud provoca la comida basura, según los nutricionistas?
\item ¿Cuál es la posición del autor? ¿Qué expresiones usa para indicarla?
\end{enumerate}

\courssub{Deux champs lexicaux : alimentación sana / comida basura}

Le texte s'organise autour de deux \cle{champs lexicaux} opposés. Mémorisez-les par familles grammaticales (noms, adjectifs, verbes, expressions), car un bon débatteur passe facilement d'une catégorie à l'autre.

\begin{center}
\begin{tabularx}{\linewidth}{|X|X|X|}
\hline
\rowcolor{popblueL} Español & Français & Catégorie \\
\hline
alimentos frescos & aliments frais & nom + adj. \\
\hline
verduras de hoja verde & légumes à feuilles vertes & nom \\
\hline
proteínas magras & protéines maigres & nom + adj. \\
\hline
la hidratación & l'hydratation & nom \\
\hline
la fibra & les fibres & nom \\
\hline
una dieta equilibrada & un régime équilibré & expression \\
\hline
los nutrientes / nutritivo & les nutriments / nutritif & nom / adj. \\
\hline
prevenir enfermedades & prévenir les maladies & verbe + nom \\
\hline
\rowcolor{poporangeL} la comida chatarra / basura & la malbouffe & nom \\
\hline
\rowcolor{poporangeL} alimentos ultraprocesados & aliments ultra-transformés & nom + adj. \\
\hline
\rowcolor{poporangeL} grasas saturadas & graisses saturées & nom + adj. \\
\hline
\rowcolor{poporangeL} azúcares añadidos & sucres ajoutés & nom + adj. \\
\hline
\rowcolor{poporangeL} los aditivos / conservantes & les additifs / conservateurs & nom \\
\hline
\rowcolor{poporangeL} los refrescos azucarados & les sodas sucrés & nom \\
\hline
\rowcolor{poporangeL} la obesidad / la diabetes & l'obésité / le diabète & nom \\
\hline
\rowcolor{poporangeL} la malnutrición & la malnutrition & nom \\
\hline
\end{tabularx}
\end{center}

\begin{definition}[Trois termes clés à ne pas confondre]
\begin{itemize}
\item \esp{Comida basura} = \esp{comida chatarra} : termes courants et journalistiques pour « malbouffe ». Le terme technique, utilisé par l'OMS, est \esp{alimentos ultraprocesados}.
\item \esp{Dieta equilibrada} = régime alimentaire équilibré (attention : \esp{dieta} désigne l'alimentation en général, pas seulement un « régime » pour maigrir).
\item \esp{Malnutrición} = malnutrition : elle désigne aussi bien la \emph{carence} (falta de nutrientes) que l'\emph{excès} de produits de mauvaise qualité. Ne la réduisez pas à la sous-alimentation.
\end{itemize}
\end{definition}

\begin{popfigure}
\begin{tikzpicture}[
  box/.style={rounded corners, draw=#1, fill=#1!12, text width=4.6cm, align=center, font=\small, inner sep=5pt},
  head/.style={rounded corners, draw=#1, fill=#1!35, text width=4.6cm, align=center, font=\small\bfseries, inner sep=5pt}]
\node[head=popgreen] (h1) at (0,0) {ALIMENTACIÓN SANA};
\node[box=popgreen, align=center] (a1) at (0,-1.4){alimentos frescos\\ verduras de hoja verde};
\node[box=popgreen, align=center] (a2) at (0,-2.8){proteínas magras\\ fibra e hidratación};
\node[box=popgreen, align=center] (a3) at (0,-4.2){dieta equilibrada\\ prevenir enfermedades};
\node[head=poporange] (h2) at (7,0) {COMIDA BASURA};
\node[box=poporange, align=center] (b1) at (7,-1.4){comida chatarra\\ ultraprocesados};
\node[box=poporange, align=center] (b2) at (7,-2.8){grasas saturadas\\ azúcares añadidos};
\node[box=poporange, align=center] (b3) at (7,-4.2){aditivos y conservantes\\ obesidad y diabetes};
\draw[-{Stealth}, thick, popmuted] (h1) -- (h2) node[midway, above, font=\small\itshape, text=popmuted] {versus};
\draw[thick, popgreen] (h1) -- (a1); \draw[thick, popgreen] (a1) -- (a2); \draw[thick, popgreen] (a2) -- (a3);
\draw[thick, poporange] (h2) -- (b1); \draw[thick, poporange] (b1) -- (b2); \draw[thick, poporange] (b2) -- (b3);
\end{tikzpicture}
\legende{Carte mentale : les deux champs lexicaux du débat alimentaire.}
\end{popfigure}

\courssub{Familias de palabras : construire le vocabulaire en réseau}

Un excellent moyen d'enrichir rapidement son lexique est de maîtriser les \cle{familles de mots} : à partir d'une même racine, l'espagnol forme verbe, nom et adjectifs. Observez :

\begin{center}
\begin{tabular}{|l|l|l|}
\hline
\rowcolor{popgreenL} Verbo & Sustantivo & Adjetivo \\
\hline
nutrir (nourrir) & la nutrición / el nutriente & nutritivo, -a \\
\hline
procesar (transformer) & el procesamiento & procesado, -a \\
\hline
hidratarse (s'hydrater) & la hidratación & hidratado, -a \\
\hline
engordar (grossir) & el engorde & gordo, -a / obeso, -a \\
\hline
prevenir (prévenir) & la prevención & preventivo, -a \\
\hline
\end{tabular}
\end{center}

\begin{exemplebox}[Les familles en contexte]
\esp{Las frutas \textbf{nutren} el organismo.} « Les fruits nourrissent l'organisme. »

\esp{Una buena \textbf{nutrición} es fundamental para los adolescentes.} « Une bonne nutrition est fondamentale pour les adolescents. »

\esp{Los alimentos \textbf{procesados} contienen muchos aditivos.} « Les aliments transformés contiennent beaucoup d'additifs. »

\esp{El \textbf{procesamiento} industrial destruye muchas vitaminas.} « La transformation industrielle détruit beaucoup de vitamines. »
\end{exemplebox}

\courssub{Expresiones para tomar posición}

Un débat exige d'\emph{assumer} sa position. Voici les formules canoniques, classées par degré d'engagement.

\begin{propriete}[Estoy a favor de / en contra de + nom ou infinitif]
\textbf{Forme :} \esp{estar a favor de / estar en contra de / oponerse a} + \textbf{nom} ou \textbf{infinitif} (jamais de subjonctif après ces prépositions).

\esp{Estoy a favor de \textbf{comer} frutas todos los días.} « Je suis favorable au fait de manger des fruits tous les jours. »

\esp{Me opongo a la comida basura.} « Je m'oppose à la malbouffe. »

\esp{Estoy en contra de prohibir totalmente los refrescos.} « Je suis contre l'interdiction totale des sodas. »

\textbf{Attention :} \esp{oponerse a} est pronominal : \esp{me opongo, te opones, se opone...}
\end{propriete}

\begin{aretenir}[Boîte à outils de l'opinion]
\begin{itemize}
\item \esp{Considero que...} / \esp{Pienso que...} / \esp{Creo que...} + indicatif : « Je considère que... »
\item \esp{Desde mi punto de vista...} / \esp{En mi opinión...} / \esp{A mi juicio...} : « De mon point de vue... »
\item \esp{Estoy convencido de que...} : « Je suis convaincu que... »
\item \esp{Me parece que...} : « Il me semble que... » (plus nuancé)
\end{itemize}
\end{aretenir}

\begin{exemplebox}[Exemples traduits]
\esp{Considero que una dieta rica en vegetales previene enfermedades.} « Je considère qu'un régime riche en légumes prévient les maladies. »

\esp{Sin embargo, la comida rápida es más accesible económicamente.} « Cependant, la restauration rapide est plus accessible économiquement. »

\esp{Desde mi punto de vista, los colegios deben educar en hábitos saludables.} « De mon point de vue, les collèges doivent éduquer aux habitudes saines. »
\end{exemplebox}

\courssub{Conectores lógicos : articular la argumentación}

Les connecteurs sont la charpente invisible du débat. Sans eux, vos arguments restent une liste ; avec eux, ils deviennent un raisonnement.

\begin{center}
\begin{tabularx}{\linewidth}{|X|X|X|}
\hline
\rowcolor{poppurpleL} Conector & Función & Ejemplo \\
\hline
además & ajouter un argument & Además, la fruta local es barata. \\
\hline
sin embargo & opposer & Sin embargo, muchos jóvenes la ignoran. \\
\hline
en cambio & contraster & En cambio, los refrescos son caros y dañinos. \\
\hline
por ejemplo & illustrer & Por ejemplo, el ñame es muy nutritivo. \\
\hline
por tanto / por eso & conclure partiellement & Por tanto, debemos consumir productos locales. \\
\hline
en conclusión & conclure & En conclusión, comer sano es posible para todos. \\
\hline
\end{tabularx}
\end{center}

\begin{attention}[Faux amis et pièges de registre]
\begin{itemize}
\item \esp{Preparado} ne signifie pas « préparé » (cuisiné) dans le contexte alimentaire : \esp{alimentos preparados} désigne souvent des plats \emph{industriels} (connotation négative). Pour « fait maison », dites \esp{casero}.
\item \esp{Conservante} = conservateur (additif chimique), pas « conservé ». Un aliment conservé = \esp{un alimento en conserva}.
\item \esp{Por ejemplo} ne se traduit jamais par « par exemple » avec \esp{para} : on dit \esp{por ejemplo}, sans faute possible.
\item Ne confondez pas \esp{en cambio} (au contraire, contraste) et \esp{sin embargo} (cependant, concession).
\end{itemize}
\end{attention}

\courssub{Del registro familiar al registro formal}

Dans un débat scolaire ou un article d'opinion, le registre familier doit être « traduit » en registre soutenu. C'est une compétence d'examen.

\begin{center}
\begin{tabularx}{\linewidth}{|X|X|}
\hline
\rowcolor{popgoldL} Registro familiar & Registro formal / técnico \\
\hline
comida chatarra / basura & alimentos ultraprocesados \\
\hline
comida sana & dieta equilibrada / alimentación saludable \\
\hline
hace daño & perjudica la salud / provoca enfermedades \\
\hline
engorda mucho & favorece la obesidad \\
\hline
es bueno para ti & resulta beneficioso para el organismo \\
\hline
\end{tabularx}
\end{center}

\begin{exoresolu}[Reformulation vers le registre soutenu]
Énoncé : reformulez en registre formel la phrase suivante : \esp{La comida chatarra engorda mucho y hace daño a los chicos.}
\tcblower
Solution : on remplace chaque élément familier par son équivalent technique : \esp{comida chatarra} $\rightarrow$ \esp{alimentos ultraprocesados} ; \esp{engorda mucho} $\rightarrow$ \esp{favorecen la obesidad} ; \esp{hace daño} $\rightarrow$ \esp{perjudican la salud} ; \esp{los chicos} $\rightarrow$ \esp{los jóvenes}.

\esp{Los alimentos ultraprocesados favorecen la obesidad y perjudican la salud de los jóvenes.} « Les aliments ultra-transformés favorisent l'obésité et nuisent à la santé des jeunes. »
\end{exoresolu}

\begin{exercice}[Classement et reformulation]
1. Classez ces mots dans la colonne qui convient (\esp{alimentación sana} / \esp{comida basura}) : \esp{aditivos, fibra, grasas saturadas, verduras de hoja verde, azúcares añadidos, proteínas magras, refrescos, hidratación}.

2. Reformulez en registre soutenu : \esp{La comida rápida es mala para la salud pero a los alumnos les gusta mucho.}
\end{exercice}

\begin{corrige}[Exercice : classement et reformulation]
1. \esp{Alimentación sana} : fibra, verduras de hoja verde, proteínas magras, hidratación. \esp{Comida basura} : aditivos, grasas saturadas, azúcares añadidos, refrescos.

2. Proposition : \esp{La comida rápida resulta perjudicial para la salud; sin embargo, goza de gran popularidad entre los alumnos.} « La restauration rapide s'avère nuisible pour la santé ; cependant, elle jouit d'une grande popularité parmi les élèves. »
\end{corrige}

\courssub{Tu ficha de vocabulario personal}

Pour préparer le débat final, constituez une fiche personnelle réutilisable. Elle doit contenir, pour chaque camp : cinq noms, trois adjectifs, trois verbes, deux expressions d'opinion et trois connecteurs.

\begin{methode}[Construire sa fiche de vocabulaire de débat]
\begin{enumerate}
\item Choisissez votre thèse : \esp{Estoy a favor de...} ou \esp{Me opongo a...}
\item Sélectionnez 8 à 10 mots clés de chaque champ lexical (section 5.2), avec leur traduction française.
\item Ajoutez deux familles de mots (ex. \esp{nutrir / nutrición / nutritivo}) pour varier les formulations.
\item Préparez trois phrases-types : une opinion (\esp{Considero que...}), une concession (\esp{Sin embargo...}), une conclusion (\esp{En conclusión...}).
\item Vérifiez chaque faux ami (section 5.5) avant de figer votre fiche.
\end{enumerate}
\end{methode}

\begin{experience}[Mini-débat éclair]
Par binômes : l'élève A défend \esp{la alimentación sana}, l'élève B défend l'accessibilité de \esp{la comida rápida}. Chacun dispose d'une minute et doit utiliser au moins deux connecteurs (\esp{además, sin embargo, por tanto...}) et une expression d'opinion. La classe vote ensuite pour l'argumentation la mieux construite — pas forcément la plus convaincante sur le fond !
\end{experience}

\begin{savaistu}[Et au Cameroun ?]
En Guinée équatoriale, pays hispanophone voisin du Cameroun, l'alimentation traditionnelle repose sur le manioc, le plantain et le poisson : une \esp{dieta equilibrada} naturelle. Comme à Yaoundé ou à Douala, l'arrivée des \esp{alimentos ultraprocesados} y transforme les habitudes, surtout chez les jeunes urbains. Le débat que vous préparez est donc un vrai débat de société dans tout le monde hispanophone... et chez nous.
\end{savaistu}

\begin{aretenir}[L'essentiel de la section]
\begin{itemize}
\item Deux champs lexicaux opposés : \esp{alimentación sana} (frescos, fibra, dieta equilibrada) vs \esp{comida basura} (ultraprocesados, grasas saturadas, aditivos).
\item Familles de mots : \esp{nutrir / nutrición / nutritivo} ; \esp{procesar / procesamiento / procesado}.
\item Prise de position : \esp{Estoy a favor de / Me opongo a} + nom ou infinitif ; \esp{Considero que} + indicatif.
\item Connecteurs : \esp{además, sin embargo, en cambio, por ejemplo, por tanto, en conclusión}.
\item Registre soutenu : \esp{comida chatarra} $\rightarrow$ \esp{alimentos ultraprocesados}.
\end{itemize}
\end{aretenir}

\coursec{Grammaire \& Traduction : structures d'argumentation et d'opinion}

Dans la section précédente, nous avons enrichi notre vocabulaire du débat : \esp{alimentación sana}, \esp{comida basura}, \esp{nutrientes}, \esp{grasas saturadas}... Mais posséder des mots ne suffit pas pour convaincre. Pour \cle{prendre position} et \cle{justifier} dans un débat, il faut maîtriser trois outils grammaticaux : le \cle{subjonctif} après les verbes d'opinion, les \cle{connecteurs logiques} qui articulent le raisonnement, et les \cle{comparatifs} qui opposent les deux types d'alimentation. C'est l'objet de cette section, directement utile pour l'épreuve de thème et pour l'expression orale du Baccalauréat.

\courssub{Texte de départ : un éditorial sur l'alimentation}

Observons d'abord un court article de presse. Soulignez mentalement tous les verbes conjugués au subjonctif et tous les connecteurs : ils seront la matière de notre analyse.

\begin{textebox}[Editorial : « La mochila y el menú »]
No creo que los jóvenes cameruneses coman peor que sus padres ; sin embargo, es evidente que consumen más comida rápida que antes. En los barrios de Yaoundé y Douala, las hamburgueserías se multiplican porque son baratas y rápidas. Puesto que un menú cuesta menos que un plato de \emph{ndolé} bien preparado, muchos estudiantes eligen la opción fácil.

Es importante que las familias recuperen la cocina tradicional, tan rica como variada : el \emph{eru}, el \emph{okok}, las frutas tropicales. Aunque la comida rápida sea práctica, no es tan nutritiva como la comida casera. Dudo que un adolescente pueda estudiar bien si solo bebe refrescos azucarados. Por tanto, recomiendo que los colegios ofrezcan menús equilibrados y que los padres cocinen más en casa. Es esencial que todos comprendan que comer sano no es un lujo, sino una inversión en el futuro.
\end{textebox}

\textbf{Glossaire :} \esp{la mochila} : le sac à dos ; \esp{la hamburguesería} : le fast-food ; \esp{la opción fácil} : la solution de facilité ; \esp{la comida casera} : la cuisine maison ; \esp{el refresco azucarado} : la boisson sucrée ; \esp{un lujo} : un luxe ; \esp{la inversión} : l'investissement ; \esp{recuperar} : retrouver, remettre au goût du jour.

\textbf{Questions de compréhension :}
\begin{enumerate}
\item ¿Por qué los estudiantes eligen la comida rápida, según el autor ?
\item ¿Qué recomienda el autor a los colegios y a los padres ?
\item Relevez trois verbes au subjonctif et dites après quelle expression ils apparaissent.
\end{enumerate}

\textbf{Éléments de réponse (question 3) :} \esp{coman} après \esp{No creo que} ; \esp{recuperen} après \esp{Es importante que} ; \esp{sea} après \esp{Aunque} ; \esp{pueda} après \esp{Dudo que} ; \esp{ofrezcan} et \esp{cocinen} après \esp{recomiendo que} ; \esp{comprendan} après \esp{Es esencial que}. Toutes ces expressions ont un point commun : elles n'affirment pas un fait, elles expriment un doute, un souhait ou un jugement. Retenez bien cette observation, elle fonde la règle de la section 6.3.

\courssub{Révision du subjonctif présent : formation}

\begin{propriete}[Formation du subjonctif présent]
On part du radical de la 1\textsuperscript{re} personne du singulier du présent de l'indicatif (\esp{yo}), auquel on ajoute les terminaisons « inversées » :
\begin{itemize}
\item verbes en \esp{-ar} : \textbf{-e, -es, -e, -emos, -éis, -en} (\esp{hable, hables, hable, hablemos, habléis, hablen}) ;
\item verbes en \esp{-er} / \esp{-ir} : \textbf{-a, -as, -a, -amos, -áis, -an} (\esp{coma, comas, coma, comamos, comáis, coman} ; \esp{viva, vivas, viva, vivamos, viváis, vivan}).
\end{itemize}
Le radical suit les mêmes irrégularités que la forme \esp{yo} de l'indicatif : \esp{tener} $\rightarrow$ \esp{tenga} ; \esp{hacer} $\rightarrow$ \esp{haga} ; \esp{poner} $\rightarrow$ \esp{ponga} ; \esp{salir} $\rightarrow$ \esp{salga} ; \esp{venir} $\rightarrow$ \esp{venga} ; \esp{decir} $\rightarrow$ \esp{diga} ; \esp{conocer} $\rightarrow$ \esp{conozca}.
\end{propriete}

\begin{center}
\begin{tabular}{|l|l|l|l|}
\hline
\rowcolor{popblueL} \textbf{Infinitivo} & \textbf{yo (indic.)} & \textbf{Subjuntivo presente} & \textbf{Français} \\
\hline
tener & tengo & tenga, tengas, tenga, tengamos, tengáis, tengan & avoir \\
\hline
venir & vengo & venga, vengas, venga, vengamos, vengáis, vengan & venir \\
\hline
poner & pongo & ponga, pongas, ponga, pongamos, pongáis, pongan & mettre \\
\hline
salir & salgo & salga, salgas, salga, salgamos, salgáis, salgan & sortir \\
\hline
haber & he & haya, hayas, haya, hayamos, hayáis, hayan & avoir (aux.) \\
\hline
ser & soy & sea, seas, sea, seamos, seáis, sean & être \\
\hline
ir & voy & vaya, vayas, vaya, vayamos, vayáis, vayan & aller \\
\hline
saber & sé & sepa, sepas, sepa, sepamos, sepáis, sepan & savoir \\
\hline
\end{tabular}
\end{center}

\begin{attention}[Verbes totalement irréguliers]
\esp{ser} $\rightarrow$ \esp{sea}, \esp{ir} $\rightarrow$ \esp{vaya}, \esp{saber} $\rightarrow$ \esp{sepa}, \esp{haber} $\rightarrow$ \esp{haya} : ces formes ne se devinent pas, elles s'apprennent par cœur. Erreur fréquente des francophones : *\esp{No creo que es sano} au lieu de \esp{No creo que sea sano}. Autre piège : ne confondez pas \esp{haya} (subjonctif de \esp{haber}), \esp{halla} (de \esp{hallar}, trouver) et \esp{aya} (une gouvernante).
\end{attention}

\courssub{Opinion et subjonctif : indicatif ou subjonctif ?}

Reprenons le texte : \esp{No creo que los jóvenes coman peor} (subjonctif), mais on dirait \esp{Creo que los jóvenes comen peor} (indicatif). Pourquoi ce changement ? Parce que la négation transforme une certitude en doute : dès que le locuteur n'affirme plus un fait mais exprime son doute, son souhait ou son jugement, l'espagnol passe au subjonctif.

\begin{propriete}[Règle d'or : certitude ou subjectivité]
\begin{itemize}
\item Après \esp{Creo que}, \esp{Pienso que}, \esp{Me parece que}, \esp{Estoy seguro de que}, \esp{Es verdad que}, \esp{Es evidente que}, \esp{Es cierto que} $\rightarrow$ \textbf{INDICATIF} (le locuteur affirme, il est certain).
\item Après \esp{No creo que}, \esp{No pienso que}, \esp{Dudo que}, \esp{Es posible que}, \esp{Quizás / Tal vez}, \esp{Es necesario que}, \esp{Es importante que}, \esp{Es esencial que}, \esp{Recomiendo que}, \esp{Quiero que}, \esp{Ojalá} $\rightarrow$ \textbf{SUBJONCTIF} (doute, possibilité, volonté, jugement de valeur).
\end{itemize}
\end{propriete}

\begin{popfigure}
\begin{center}
\begin{tikzpicture}[node distance=0.4cm]
\node[draw=popblue, fill=popblueL, rounded corners, align=center, font=\bfseries] (titre) {¿Indicativo o subjuntivo?};
\node[draw=popgreen, fill=popgreenL, rounded corners, align=center, below left=0.7cm and -1.2cm of titre, text width=5.4cm] (ind) {\textbf{INDICATIVO}\\ Creo que...\\ Pienso que...\\ Es evidente que...\\ \emph{« Creo que la fruta es sana. »}\\ (Je crois que le fruit est sain.)};
\node[draw=poporange, fill=poporangeL, rounded corners, align=center, below right=0.7cm and -1.2cm of titre, text width=5.4cm] (sub) {\textbf{SUBJUNTIVO}\\ No creo que...\\ Dudo que...\\ Es importante que...\\ \emph{« Dudo que sea sano. »}\\ (Je doute que ce soit sain.)};
\draw[-{Stealth}, thick, popgreen] (titre) -- (ind);
\draw[-{Stealth}, thick, poporange] (titre) -- (sub);
\end{tikzpicture}
\legende{Expressions d'opinion : indicatif (certitude) ou subjonctif (subjectivité)}
\end{center}
\end{popfigure}

\begin{exemplebox}[Exemples traduits]
\esp{Creo que la comida casera es más sana.} « Je pense que la cuisine maison est plus saine. » (indicatif)

\esp{No creo que los refrescos sean buenos para la salud.} « Je ne pense pas que les sodas soient bons pour la santé. » (subjonctif)

\esp{Es importante que los niños beban agua.} « Il est important que les enfants boivent de l'eau. » (subjonctif)

\esp{Recomiendo que comáis menos sal.} « Je recommande que vous mangiez moins de sel. » (subjonctif)

\esp{Ojalá los colegios ofrezcan frutas en el recreo.} « Pourvu que les écoles offrent des fruits à la récréation. » (subjonctif, souhait)
\end{exemplebox}

\begin{attention}[Pièges de traduction]
En français, « il faut que je mange » prend le subjonctif des deux côtés : facile. Mais attention à \esp{Es verdad que} + \textbf{indicatif} : \esp{Es verdad que la comida basura es barata} « Il est vrai que la malbouffe est bon marché ». La négation change tout : \esp{No es verdad que sea barata}. De même, \esp{quizás} et \esp{tal vez} (« peut-être ») exigent le subjonctif : \esp{Quizás tengan razón.} « Peut-être ont-ils raison. »
\end{attention}

\courssub{Les connecteurs logiques : l'ossature du débat}

Un débat sans connecteurs est une liste de slogans. Voici l'armature complète à réutiliser dans vos prises de parole et vos articles d'opinion :

\begin{center}
\begin{tabularx}{\linewidth}{|l|X|X|}
\hline
\rowcolor{popblueL} \textbf{Relation} & \textbf{Español} & \textbf{Français} \\
\hline
Cause & porque, ya que, puesto que, como & parce que, puisque, comme \\
\hline
Conséquence & por tanto, por consiguiente, así que, entonces & donc, par conséquent, ainsi \\
\hline
Opposition & pero, sin embargo, no obstante, en cambio & mais, cependant, néanmoins, en revanche \\
\hline
Concession & aunque (+ ind. / + subj.), a pesar de & bien que, quoique, malgré \\
\hline
But & para, para que (+ subj.), a fin de que (+ subj.) & pour, pour que, afin que \\
\hline
Exemple & por ejemplo, como, así & par exemple, comme, ainsi \\
\hline
Addition & además, también, encima & de plus, aussi, en outre \\
\hline
Conclusion & en conclusión, en resumen, para terminar & en conclusion, en résumé, pour finir \\
\hline
\end{tabularx}
\end{center}

\begin{exemplebox}[Connecteurs en contexte]
\esp{Como la comida rápida es barata, muchos jóvenes la compran.} « Comme la restauration rapide est bon marché, beaucoup de jeunes l'achètent. »

\esp{Las frutas son caras ; sin embargo, son necesarias.} « Les fruits sont chers ; cependant, ils sont nécessaires. »

\esp{Come demasiada grasa ; por tanto, debe hacer deporte.} « Il mange trop de gras ; par conséquent, il doit faire du sport. »

\esp{Traigo fruta al colegio para que mis amigos coman sano.} « J'apporte des fruits au collège pour que mes amis mangent sainement. » (\esp{para que} + subjonctif)
\end{exemplebox}

\begin{attention}[Le cas spécial de \esp{aunque}]
\esp{Aunque} + \textbf{indicatif} = fait réel, constaté : \esp{Aunque llueve, salgo.} « Bien qu'il pleuve (et il pleut vraiment), je sors. »

\esp{Aunque} + \textbf{subjonctif} = hypothèse, concession : \esp{Aunque llueva, saldré.} « Même s'il pleuvait, je sortirais. »

De même : \esp{Aunque la comida rápida es barata...} (c'est un fait) vs \esp{Aunque la comida rápida sea barata...} (admettons qu'elle le soit). Le français « bien que » impose toujours le subjonctif ; l'espagnol, lui, choisit selon la réalité du fait. Attention aussi à la place de \esp{como} : en tête de phrase, il exprime la cause (« comme » = « puisque ») ; au milieu, il exprime la comparaison.
\end{attention}

\courssub{Les comparatifs : confronter les deux alimentations}

\begin{propriete}[Comparatifs d'inégalité et d'égalité]
\begin{itemize}
\item Supériorité : \esp{más + adjectif + que} : \esp{La fruta es más sana que el dulce.} « Le fruit est plus sain que le bonbon. »
\item Infériorité : \esp{menos + adjectif + que} : \esp{El agua es menos dulce que el refresco.} « L'eau est moins sucrée que le soda. »
\item Égalité : \esp{tan + adjectif + como} : \esp{El ndolé es tan rico como el eru.} « Le ndolé est aussi riche que l'eru. »
\item Avec un nom : \esp{más/menos + nom + que} ; \esp{tanto/tanta + nom + como} : \esp{Como tanta fruta como verdura.} « Je mange autant de fruits que de légumes. »
\end{itemize}
\textbf{Irréguliers :} \esp{bueno} $\rightarrow$ \esp{mejor} (meilleur) ; \esp{malo} $\rightarrow$ \esp{peor} (pire) ; \esp{grande} $\rightarrow$ \esp{mayor} (plus grand, plus âgé) ; \esp{pequeño} $\rightarrow$ \esp{menor} (plus petit, plus jeune).
\end{propriete}

\begin{attention}[Faux amis et calques]
Ne dites jamais *\esp{más bueno} : c'est \esp{mejor}. De même *\esp{más malo} $\rightarrow$ \esp{peor}. Attention aussi : « fast food » ne se traduit pas seulement par \esp{comida rápida} (neutre) ; pour le sens péjoratif du débat, préférez \esp{comida basura} ou \esp{comida chatarra} (Amérique latine). Enfin, « sain » se dit \esp{sano} (pour les aliments et les personnes) ; ne calquez pas *\esp{saludablemente sano}.
\end{attention}

\courssub{Méthode de traduction : thème et version}

\begin{methode}[Traduire une phrase argumentative en 5 étapes]
\begin{enumerate}
\item \textbf{Analyser} la structure de la phrase française (proposition principale, subordonnée, connecteur).
\item \textbf{Identifier} le temps et le mode requis : opinion affirmative $\rightarrow$ indicatif ; doute, souhait, jugement $\rightarrow$ subjonctif.
\item \textbf{Choisir} le vocabulaire précis (éviter les calques : « malbouffe » = \esp{comida basura}, pas *\esp{mala comida}).
\item \textbf{Vérifier} les accords (genre, nombre) et l'ordre des mots (adjectif après le nom en général).
\item \textbf{Relire} à voix haute : la phrase doit sonner espagnol, pas français traduit.
\end{enumerate}
\end{methode}

\begin{exoresolu}[Thème résolu : subjonctif + comparatif]
Énoncé : traduisez « Il est essentiel que les jeunes mangent plus de fruits et moins de graisses. »
\tcblower
\textbf{Étape 1 :} structure = « il est essentiel que » + subordonnée. \textbf{Étape 2 :} jugement de valeur impersonnel $\rightarrow$ subjonctif : \esp{Es esencial que... coman}. \textbf{Étape 3 :} « plus de fruits » = \esp{más frutas} ; « moins de graisses » = \esp{menos grasas}. \textbf{Étape 4 :} \esp{los jóvenes} (masculin pluriel), verbe \esp{coman} (3\textsuperscript{e} pers. plur.).

\textbf{Solution :} \esp{Es esencial que los jóvenes coman más frutas y menos grasas.}
\end{exoresolu}

\begin{exoresolu}[Version résolue : concession + opposition]
Énoncé : traduisez \esp{Aunque la comida rápida es barata, no es tan nutritiva como la comida casera.}
\tcblower
\textbf{Étape 1 :} \esp{aunque} + indicatif (\esp{es}) = fait réel $\rightarrow$ « bien que » ou « même si ». \textbf{Étape 2 :} \esp{comida rápida} = « restauration rapide » ; \esp{barata} = « bon marché ». \textbf{Étape 3 :} \esp{tan nutritiva como} = comparatif d'égalité : « aussi nutritive que » ; \esp{comida casera} = « cuisine maison ».

\textbf{Solution :} « Bien que la restauration rapide soit bon marché, elle n'est pas aussi nutritive que la cuisine maison. »
\end{exoresolu}

\courssub{Exercices de thème (français $\rightarrow$ espagnol)}

\begin{exercice}[Thème : 5 phrases]
Traduisez en espagnol :
\begin{enumerate}
\item Je ne pense pas que les hamburgers soient sains.
\item Il est important que les élèves boivent plus d'eau que de sodas.
\item Bien que les fruits soient chers à Douala, ils sont meilleurs que les bonbons.
\item La cuisine maison est aussi savoureuse que la restauration rapide ; cependant, elle est plus saine.
\item Je recommande que les cantines offrent des repas équilibrés, parce que les adolescents ont besoin d'énergie.
\end{enumerate}
\end{exercice}

\begin{corrige}[Exercice de thème]
\begin{enumerate}
\item \esp{No creo que las hamburguesas sean sanas.} (\esp{no creo que} + subjonctif \esp{sean})
\item \esp{Es importante que los alumnos beban más agua que refrescos.} (subjonctif \esp{beban} + comparatif)
\item \esp{Aunque las frutas son caras en Douala, son mejores que los dulces.} (\esp{aunque} + indicatif car fait réel ; \esp{mejores}, pas *\esp{más buenas})
\item \esp{La comida casera es tan sabrosa como la comida rápida ; sin embargo, es más sana.}
\item \esp{Recomiendo que las cantinas ofrezcan comidas equilibradas, porque los adolescentes necesitan energía.}
\end{enumerate}
\end{corrige}

\courssub{Exercices de version (espagnol $\rightarrow$ français)}

\begin{exercice}[Version : 5 phrases de presse]
Traduisez en français :
\begin{enumerate}
\item \esp{Puesto que la obesidad infantil aumenta, es necesario que los colegios enseñen nutrición.}
\item \esp{No creo que la comida chatarra sea tan barata como parece : cuesta salud.}
\item \esp{Aunque beba refrescos cada día, no engordaré, piensan muchos jóvenes.}
\item \esp{Las verduras son menos caras que la carne ; por tanto, conviene comerlas a diario.}
\item \esp{Es evidente que una dieta equilibrada es mejor que cualquier producto milagro.}
\end{enumerate}
\end{exercice}

\begin{corrige}[Exercice de version]
\begin{enumerate}
\item « Puisque l'obésité infantile augmente, il est nécessaire que les écoles enseignent la nutrition. »
\item « Je ne pense pas que la malbouffe soit aussi bon marché qu'elle en a l'air : elle coûte la santé. »
\item « Même si je bois des sodas chaque jour, je ne grossirai pas, pensent beaucoup de jeunes. » (\esp{aunque} + subjonctif = hypothèse)
\item « Les légumes sont moins chers que la viande ; par conséquent, il convient d'en manger tous les jours. »
\item « Il est évident qu'un régime équilibré est meilleur que n'importe quel produit miracle. »
\end{enumerate}
\end{corrige}

\begin{savaistu}[Nutrition et étiquetage dans le monde hispanique]
L'Organisation mondiale de la santé recommande de consommer \esp{menos de 5 gramos de sal al día} (moins de 5 grammes de sel par jour). En Espagne, l'étiquetage nutritionnel \emph{Nutri-Score}, adopté volontairement à partir de 2021, classe les produits de A (vert, le plus sain) à E (rouge) pour aider les consommateurs à choisir. Au Mexique et au Chili, des étiquettes noires d'avertissement signalent les produits \esp{altos en azúcares} (riches en sucres). De quoi alimenter vos arguments dans le débat !
\end{savaistu}

\begin{aretenir}[L'essentiel de la section]
\begin{itemize}
\item Subjonctif présent : radical de \esp{yo} + terminaisons inversées ; irréguliers : \esp{tenga, venga, ponga, salga, haya, sea, vaya, sepa}.
\item Certitude (\esp{creo que, es evidente que}) $\rightarrow$ indicatif ; doute, volonté, jugement (\esp{no creo que, dudo que, es importante que, ojalá}) $\rightarrow$ subjonctif.
\item Connecteurs : \esp{porque / por tanto / sin embargo / aunque / para que / por ejemplo} : l'ossature de toute argumentation.
\item Comparatifs : \esp{más... que}, \esp{menos... que}, \esp{tan... como} ; irréguliers : \esp{mejor, peor, mayor, menor}.
\item \esp{Aunque} + indicatif = fait réel ; \esp{aunque} + subjonctif = hypothèse.
\end{itemize}
\end{aretenir}

\coursec{Production guidée : débat oral et rédaction d'un article d'opinion}

Dans la section précédente, nous avons appris à exprimer notre opinion et à argumenter grâce aux structures grammaticales essentielles (\esp{creo que + indicativo}, \esp{no creo que + subjuntivo}, les comparatifs, les connecteurs logiques). Il est temps de passer de la théorie à la pratique : nous allons d'abord organiser un véritable débat oral en classe, puis transformer nos arguments en un article d'opinion digne du journal du lycée. C'est exactement le type de compétence exigée à l'oral du baccalauréat : prendre position, justifier, convaincre.

\courssub{Mise en place du débat : \esp{¿A dieta estricta o libertad con moderación?}}

\textbf{Observation.} Lisez d'abord ces deux prises de position opposées, extraites d'un forum de lycéens :

\esp{A mí me parece que deberíamos seguir una alimentación sana y estricta: la comida chatarra es un veneno dulce que destruye nuestra salud poco a poco.}

« À mon avis, nous devrions suivre une alimentation saine et stricte : la malbouffe est un poison doux qui détruit notre santé petit à petit. »

\esp{Yo no creo que sea necesario prohibir nada: lo importante es la moderación. Unos \emph{tentempiés} de vez en cuando no matan a nadie.}

« Je ne crois pas qu'il soit nécessaire d'interdire quoi que ce soit : l'important, c'est la modération. Quelques en-cas de temps en temps ne tuent personne. »

Remarquez les marqueurs de prise de position (\esp{a mí me parece que}, \esp{yo no creo que + subjuntivo}) et l'image poétique (\esp{un veneno dulce}) : un bon débatteur combine grammaire correcte et style expressif.

\begin{propriete}[Règles du débat (obligatoires)]
\begin{itemize}
\item La classe est divisée en \cle{deux groupes} : le groupe A défend une alimentation saine et stricte ; le groupe B défend la liberté de choisir, avec modération.
\item Chaque intervention est limitée à \cle{une minute}. On lève la main et on attend la parole du président (l'enseignant).
\item Chaque prise de parole doit contenir \cle{au minimum} : \textbf{3 connecteurs} (\esp{además, sin embargo, por tanto, en primer lugar...}), \textbf{2 comparatifs} (\esp{más sano que, menos graso que, tan importante como...}) et \textbf{1 subjonctif} (\esp{no creo que sea..., es importante que comamos...}).
\item Le vocabulaire spécifique du module est obligatoire : \esp{alimentación sana, comida basura, comida chatarra, ultraprocesados, nutrientes, obesidad, dieta equilibrada}.
\item Respect absolu de l'adversaire : on réfute les idées, jamais les personnes (\esp{respeto tu opinión, pero...}).
\end{itemize}
\end{propriete}

\begin{attention}[Anglicismes à bannir du débat]
Les francophones et les anglophones mélangent souvent les termes. En espagnol correct :
\begin{center}
\begin{tabularx}{\linewidth}{|X|X|X|}\hline
\rowcolor{popblueL} Terme à éviter & Espagnol correct & Français \\ \hline
\esp{fast food} & \esp{comida rápida / comida basura} & fast-food \\ \hline
\esp{junk food} & \esp{comida chatarra} & malbouffe \\ \hline
\esp{healthy food} & \esp{comida saludable / sana} & nourriture saine \\ \hline
\esp{snacks} & \esp{aperitivos / tentempiés} & en-cas \\ \hline
\esp{fitness} & \esp{estar en forma} & être en forme \\ \hline
\end{tabularx}
\end{center}
Dire \esp{comida chatarra} est le terme le plus courant en Amérique latine ; en Espagne on préfère \esp{comida basura}.
\end{attention}

\courssub{Préparation en binômes : la fiche d'arguments}

Avant le débat, chaque binôme remplit une fiche d'arguments. Un argument convaincant repose sur trois piliers : une \cle{citation du poème} étudié (\esp{El duelo del plato}), une \cle{donnée chiffrée} (ici simulée pour l'entraînement) et un \cle{exemple personnel} ancré dans la réalité camerounaise.

\begin{exemplebox}[Fiche d'arguments modèle (groupe A, pro-alimentation saine)]
\begin{itemize}
\item \textbf{Citation du poème :} \esp{« el plato llora cuando la grasa manda »} — le poème lui-même personnifie l'assiette victime de la graisse.
\item \textbf{Donnée (simulée) :} \esp{Según un estudio escolar, siete de cada diez adolescentes consumen bebidas azucaradas a diario.} « Selon une étude scolaire, sept adolescents sur dix consomment des boissons sucrées chaque jour. »
\item \textbf{Exemple personnel :} \esp{En mi barrio de Yaoundé, las frutas del mercado son más baratas que los dulces industriales.} « Dans mon quartier de Yaoundé, les fruits du marché sont moins chers que les bonbons industriels. »
\end{itemize}
\end{exemplebox}

\begin{exemplebox}[Fiche d'arguments modèle (groupe B, pro-liberté modérée)]
\begin{itemize}
\item \textbf{Citation du poème :} \esp{« la alegría también se cocina »} — le plaisir fait partie de la cuisine.
\item \textbf{Donnée (simulée) :} \esp{Un estudio simulado indica que prohibir totalmente ciertos alimentos aumenta el deseo de consumirlos.} « Une étude simulée indique qu'interdire totalement certains aliments augmente le désir de les consommer. »
\item \textbf{Exemple personnel :} \esp{En las fiestas familiares de Douala, compartir un pastel es un gesto de amor, no un crimen.} « Lors des fêtes familiales à Douala, partager un gâteau est un geste d'amour, pas un crime. »
\end{itemize}
\end{exemplebox}

\courssub{Déroulement du débat}

\begin{experience}[Le débat parlementaire]
\begin{enumerate}
\item \textbf{Installation (5 min) :} l'enseignant joue le rôle de \esp{presidente del debate}. Les deux groupes se font face. Le président rappelle la question : \esp{¿Debemos imponer una alimentación sana estricta o defender la libertad de elegir con moderación?}
\item \textbf{Premier tour (10 min) :} chaque groupe présente ses deux arguments principaux, une minute par orateur. Le président note au tableau les \cle{arguments forts} dans deux colonnes : \esp{A favor de la dieta estricta} / \esp{A favor de la moderación}.
\item \textbf{Réfutations (10 min) :} chaque groupe attaque les arguments adverses avec \esp{sin embargo}, \esp{no estoy de acuerdo porque...}, \esp{es cierto que..., pero...}.
\item \textbf{Conclusion (5 min) :} chaque groupe résume sa position en une phrase-choc. La classe vote à main levée ; le président déclare le groupe vainqueur en félicitant les deux camps.
\end{enumerate}
\end{experience}

\begin{exoresolu}[Transformer une opinion simple en argument de débat]
Énoncé : transformez cette opinion simple en argument complet (3 connecteurs + 2 comparatifs + 1 subjonctif) : \esp{La comida sana es buena.}
\tcblower
\textbf{Solution détaillée.}
\begin{enumerate}
\item Ajoutons un connecteur d'ordre : \esp{En primer lugar...}
\item Ajoutons un comparatif : \esp{...la comida sana es mucho más beneficiosa que la comida chatarra...}
\item Ajoutons un connecteur d'addition et un second comparatif : \esp{...además, contiene menos grasas saturadas que los ultraprocesados...}
\item Ajoutons un connecteur de conséquence et un subjonctif après expression de nécessité : \esp{...por tanto, es fundamental que los jóvenes la consuman a diario.}
\end{enumerate}
\textbf{Argument final :} \esp{En primer lugar, la comida sana es mucho más beneficiosa que la comida chatarra; además, contiene menos grasas saturadas que los ultraprocesados; por tanto, es fundamental que los jóvenes la consuman a diario.}

« En premier lieu, la nourriture saine est bien plus bénéfique que la malbouffe ; de plus, elle contient moins de graisses saturées que les ultra-transformés ; par conséquent, il est fondamental que les jeunes la consomment chaque jour. »

Vérifions : 3 connecteurs (\esp{en primer lugar}, \esp{además}, \esp{por tanto}), 2 comparatifs (\esp{más beneficiosa que}, \esp{menos grasas saturadas que}), 1 subjonctif (\esp{consuman}). L'argument respecte toutes les règles du débat.
\end{exoresolu}

\courssub{L'article d'opinion : méthode et modèle}

Après l'oral, l'écrit. Chaque élève rédige individuellement un \cle{article d'opinion} de 150 à 200 mots pour le journal du lycée \esp{La Voz del Estudiante}.

\begin{methode}[Structure de l'article d'opinion en 5 étapes]
\begin{enumerate}
\item \textbf{Titre percutant} : une question (\esp{¿Comida basura o tesoro nutricional?}) ou un jeu de mots. Le titre doit donner envie de lire.
\item \textbf{Introduction} : une accroche (fait, citation, question) puis la \cle{thèse} clairement annoncée (\esp{Estoy convencido de que...}).
\item \textbf{Développement} : deux arguments, chacun suivi d'un exemple concret, reliés par des connecteurs (\esp{en primer lugar, además}).
\item \textbf{Paragraphe de concession et réfutation} : reconnaître l'argument adverse (\esp{es cierto que...}) puis le réfuter (\esp{sin embargo...}).
\item \textbf{Conclusion} : synthèse brève + appel à l'action (\esp{¡Cambiemos nuestro plato hoy!}).
\end{enumerate}
\end{methode}

\begin{popfigure}
\begin{center}
\begin{tikzpicture}[
node distance=0.45cm,
caja/.style={rectangle, rounded corners=4pt, draw=popblue, fill=popblueL, text width=6.2cm, align=center, minimum height=0.85cm, font=\small},
flecha/.style={-{Stealth[length=3mm]}, thick, poporange}]
\node[caja, fill=popgoldL, draw=popgold, align=center] (t){\textbf{1. Título percutante}\\ \esp{¿Pregunta o juego de palabras?}};
\node[caja, below=of t, align=center] (i){\textbf{2. Introducción}\\ accroche + tesis (\esp{estoy convencido de que...})};
\node[caja, below=of i, align=center] (a){\textbf{3. Argumentos}\\ 2 argumentos + ejemplos + conectores};
\node[caja, below=of a, align=center] (c){\textbf{4. Concesión y refutación}\\ \esp{es cierto que... sin embargo...}};
\node[caja, fill=popgreenL, draw=popgreen, below=of c, align=center] (f){\textbf{5. Conclusión}\\ síntesis + llamada a la acción};
\draw[flecha] (t) -- (i);
\draw[flecha] (i) -- (a);
\draw[flecha] (a) -- (c);
\draw[flecha] (c) -- (f);
\end{tikzpicture}
\end{center}
\legende{Organigramme : la structure de l'article d'opinion en cinq blocs}
\end{popfigure}

\begin{textebox}[Modèle d'article d'opinion — \esp{La Voz del Estudiante}]
¿Comida basura o tesoro nutricional? El plato de cada día decide nuestro futuro. Estoy convencido de que debemos elegir una alimentación sana si queremos vivir con energía.

En primer lugar, los productos naturales son mucho más ricos en nutrientes que los ultraprocesados. Una fruta del mercado de Yaoundé alimenta mejor que cualquier dulce industrial. Además, los ultraprocesados contienen aditivos que alteran el metabolismo; por tanto, su consumo habitual aumenta el riesgo de obesidad.

Es cierto que la comida rápida es barata y sabrosa, y que prohibirla sería exagerado. Sin embargo, no creo que el precio justifique el daño: la comida chatarra es un lobo disfrazado de cordero. Seduce, brilla, y destruye. Seduce, brilla, y destruye cada día un poco más.

En conclusión, es fundamental que los jóvenes aprendamos a cocinar y a elegir. No se trata de prohibir, sino de comprender. ¡Cambiemos nuestro plato hoy para cambiar nuestra vida mañana!
\end{textebox}

\textbf{Glossaire :} \esp{el plato} (le plat/l'assiette), \esp{los nutrientes} (les nutriments), \esp{los aditivos} (les additifs), \esp{alterar el metabolismo} (altérer le métabolisme), \esp{sabroso} (savoureux), \esp{el daño} (le tort/le dommage), \esp{un lobo disfrazado de cordero} (un loup déguisé en agneau), \esp{seducir} (séduire), \esp{no se trata de} (il ne s'agit pas de).

\textbf{Questions de compréhension :}
\begin{enumerate}
\item ¿Cuál es la tesis del autor? Cita la frase exacta.
\item ¿Qué dos argumentos defienden la alimentación sana?
\item ¿Qué figura retórica aparece en \esp{« la comida chatarra es un lobo disfrazado de cordero »}? ¿Y en la repetición de \esp{« seduce, brilla, y destruye »}?
\item ¿Cuál es la llamada a la acción final?
\end{enumerate}

\begin{propriete}[Intégrer une figure de style dans l'article]
Pour enrichir le style (et gagner des points à l'évaluation), intégrez \cle{au moins une figure de style} étudiée dans la leçon :
\begin{itemize}
\item \textbf{Métaphore :} \esp{La comida chatarra es un lobo disfrazado de cordero.} « La malbouffe est un loup déguisé en agneau. »
\item \textbf{Anaphore :} \esp{Seduce, brilla, y destruye. Seduce, brilla, y destruye cada día un poco más.} — la répétition en tête de phrase martèle l'idée.
\item \textbf{Hyperbole :} \esp{Un refresco contiene un océano de azúcar.} « Un soda contient un océan de sucre. »
\item \textbf{Personnification :} \esp{El plato llora cuando la grasa manda.} « L'assiette pleure quand la graisse commande. »
\end{itemize}
\end{propriete}

\begin{exemplebox}[Argument type à imiter]
\esp{Además, los ultraprocesados contienen aditivos que alteran el metabolismo; por tanto, su consumo habitual aumenta el riesgo de obesidad.}

« De plus, les ultra-transformés contiennent des additifs qui altèrent le métabolisme ; par conséquent, leur consommation habituelle augmente le risque d'obésité. »

Analyse : connecteur d'addition (\esp{además}) + connecteur de conséquence (\esp{por tanto}) + vocabulaire technique précis. C'est le modèle d'une phrase d'argumentation réussie.
\end{exemplebox}

\courssub{Auto-correction croisée et grille d'évaluation}

Une fois l'article rédigé, les élèves échangent leurs copies par binômes et s'évaluent mutuellement à l'aide de la grille suivante. Chaque critère est noté de 0 à 4.

\begin{center}
\begin{tabularx}{\linewidth}{|X|X|X|}\hline
\rowcolor{popblueL} Critère & Qué evaluar & Puntos \\ \hline
\esp{Contenido} & Tesis clara, 2 argumentos, ejemplos concretos & 0--4 \\ \hline
\esp{Estructura} & 5 partes presentes (título, introducción, argumentos, concesión, conclusión) & 0--4 \\ \hline
\esp{Lengua} & 3 conectores, 2 comparativos, 1 subjuntivo correctos & 0--4 \\ \hline
\esp{Vocabulario} & Léxico del módulo, sin anglicismos & 0--4 \\ \hline
\esp{Figuras retóricas} & Al menos una metáfora o anáfora & 0--4 \\ \hline
\end{tabularx}
\end{center}

Le correcteur souligne en vert les réussites, en rouge les erreurs, et écrit un conseil final (\esp{un consejo para mejorar}). L'auteur corrige ensuite sa copie avant de la remettre à l'enseignant.

\begin{attention}[Pièges fréquents dans la rédaction]
\begin{itemize}
\item Ne traduisez pas « je suis d'accord » par \esp{estoy de acuerdo}... si, c'est correct ! Mais attention à « je ne suis pas d'accord » : \esp{no estoy de acuerdo}, jamais \esp{no soy de acuerdo}.
\item « Il faut que + subjonctif » se dit \esp{es necesario que + subjuntivo} : \esp{Es necesario que comamos mejor.} Ne dites jamais \esp{es necesario que comemos}.
\item Accentuation : \esp{opinión}, \esp{nutrición}, \esp{además}, \esp{también} prennent un accent ; leur oubli est sanctionné à l'écrit.
\item \esp{saludable} s'applique aux aliments ; pour une personne en bonne santé, dites \esp{sano/a} : \esp{una comida saludable}, \esp{un chico sano}.
\end{itemize}
\end{attention}

\begin{savaistu}[L'alimentation, un classique du Bac oral]
À l'oral du baccalauréat, l'exposé suivi de discussion porte très souvent sur un thème de société : l'environnement, les réseaux sociaux, le sport... et l'alimentation est un grand classique. L'examinateur attend trois choses : un vocabulaire précis (\esp{dieta equilibrada, ultraprocesados, obesidad}), des structures d'argumentation variées (connecteurs, comparatifs, subjonctif) et la capacité à réagir spontanément aux objections. Le débat que vous venez de mener en classe est exactement l'entraînement idéal : celui qui sait défendre une position sur la \esp{comida chatarra} saura défendre n'importe quelle position le jour de l'examen.
\end{savaistu}

\begin{exercice}[À vous de jouer]
Rédigez un article d'opinion de 150 à 200 mots intitulé \esp{¿Moderación o disciplina?} pour \esp{La Voz del Estudiante}, en défendant la position inverse de celle de votre groupe au débat. Exigences : titre sous forme de question, thèse en introduction, deux arguments avec exemples camerounais, un paragraphe de concession-réfutation, une métaphore ou une anaphore, au moins trois connecteurs, deux comparatifs et un subjonctif. Terminez par la relecture croisée avec la grille d'évaluation.
\end{exercice}

\begin{corrige}[Exercice — éléments de réussite]
Une copie réussie contient : un titre interrogatif accrocheur ; une thèse explicite (\esp{Estoy convencido de que la moderación es la clave...}) ; deux arguments distincts illustrés (marché de Yaoundé, fêtes familiales de Douala) ; une concession honnête (\esp{Es cierto que la dieta estricta protege la salud...}) suivie d'une réfutation (\esp{sin embargo, prohibir genera frustración...}) ; au moins une figure de style identifiable ; les trois connecteurs, deux comparatifs et un subjonctif exigés ; une conclusion avec appel à l'action (\esp{¡Aprendamos a disfrutar con equilibrio!}). L'enseignant valorise la prise de risque stylistique autant que la correction grammaticale.
\end{corrige}

\newpage

\coursec{Activité d'intégration}

\courssub{Objectifs de l'activité}

À l'issue de cette activité d'intégration, l'élève sera capable de :

\begin{itemize}
\item mobiliser le \cle{lexique de la poésie} (\esp{verso, estrofa, rima, figura retórica}) dans une situation de communication authentique ;
\item identifier et \cle{produire des figures de style} (\esp{metáfora, personificación, hipérbole, anáfora}) dans un texte original ;
\item \cle{prendre position} dans un débat sur l'alimentation et justifier son point de vue avec des arguments structurés ;
\item utiliser correctement les \cle{structures d'opinion et d'argumentation} (\esp{opino que, me parece que, estoy a favor de, sin embargo}...) ;
\item s'exprimer à l'oral de manière claire et expressive devant un public, et évaluer la production de ses camarades à l'aide d'une grille de critères.
\end{itemize}

\courssub{Situation}

Le lycée bilingue de Yaoundé organise sa \emph{Semaine de la santé}. La chaîne locale « Visión Joven TV » a accepté de venir tourner une émission spéciale dans l'établissement, animée par les élèves de Terminale A. Le thème imposé par la rédaction : \esp{« ¿Qué hay en tu plato? Alimentación sana contra comida basura »}. Deux équipes de journalistes et d'experts s'affronteront en direct : l'équipe \esp{« Verde Vida »}, favorable à l'alimentation saine, et l'équipe \esp{« Sabor Rápido »}, défenseure de la restauration rapide. Voici le communiqué officiel de la chaîne :

\begin{textebox}[Comunicado de « Visión Joven TV »]
¡Atención, jóvenes poetas y periodistas!\\
Nuestra cadena busca dos equipos para grabar un debate de quince minutos sobre la alimentación de los adolescentes cameruneses. Cada equipo deberá presentar a un periodista cultural que lea un fragmento del poema \emph{El duelo del plato}, a dos expertos (un nutricionista y un cocinero o sociólogo) que defiendan su posición con argumentos claros, y a un reportero que realice un microtrotuario en el patio del liceo.\\
La emisión terminará con un reto poético : cada equipo compondrá un cuarteto original que resuma su posición, utilizando al menos dos figuras retóricas.\\
Criterios del jurado : contenido, corrección de la lengua, originalidad y uso de figuras de estilo. ¡Que gane el mejor plato!
\end{textebox}

\emph{Traduction :} « Attention, jeunes poètes et journalistes ! Notre chaîne cherche deux équipes pour enregistrer un débat de quinze minutes sur l'alimentation des adolescents camerounais. Chaque équipe devra présenter un journaliste culturel qui lira un fragment du poème \emph{Le Duel de l'assiette}, deux experts (un nutritionniste et un cuisinier ou sociologue) qui défendront leur position avec des arguments clairs, et un reporter qui réalisera un micro-trottoir dans la cour du lycée. L'émission se terminera par un défi poétique : chaque équipe composera un quatrain original résumant sa position, en utilisant au moins deux figures de style. Critères du jury : contenu, correction de la langue, originalité et emploi des figures de style. Que le meilleur plat gagne ! »

\courssub{Consignes}

Travail en deux équipes de six à huit élèves. Chaque élève choisit un rôle : journaliste culturel, nutritionniste, chef cuisinier, sociologue, reporter ou poète. Les tâches se déroulent en quatre étapes.

\begin{enumerate}
\item \textbf{Lecture expressive (3 minutes).} Le journaliste culturel de chaque équipe choisit une strophe du poème \esp{El duelo del plato}, la relit à voix haute et explique en deux ou trois phrases le thème du passage, sa structure (nombre de vers, rimes) et au moins une figure de style qu'il contient. Il doit employer le lexique poétique : \esp{estrofa, verso, rima, metáfora, personificación}.

\item \textbf{Table ronde (6 minutes).} Les experts de chaque équipe présentent au moins \textbf{trois arguments} pour leur position. Chaque argument doit être introduit par une structure d'opinion ou d'argumentation (\esp{opino que, estoy convencido de que, en primer lugar, además, sin embargo, por eso}) et illustré par un exemple concret du contexte camerounais (le marché de Mokolo, les bendskins et leurs snacks, le manioc, les beignets, les fast-foods de Douala...). L'équipe adverse a le droit de répliquer avec \esp{no estoy de acuerdo porque...}.

\item \textbf{Micro-trottoir simulé (3 minutes).} Le reporter de chaque équipe interviewe deux « passants » (élèves de l'autre équipe) dans la cour avec la question : \esp{¿Qué comes normalmente en el recreo y por qué?} Les interviewés doivent répondre en justifiant leur choix avec au moins un connecteur logique (\esp{porque, ya que, por eso, aunque}).

\item \textbf{Conclusion poétique (3 minutes).} Chaque équipe compose un \cle{cuarteto} (quatre vers) qui résume sa position. Le quatrain doit contenir \textbf{au moins deux figures de style} parmi : \esp{metáfora, comparación, personificación, hipérbole, aliteración, anáfora}. Le poète de l'équipe le récite pour clore l'émission.
\end{enumerate}

L'émission est filmée avec un téléphone, puis diffusée en classe. Chaque élève remplit la grille d'évaluation par les pairs :

\begin{center}
\begin{tabularx}{\linewidth}{|X|c|c|c|}
\hline
\rowcolor{popblueL} Critère & Bien (3) & Moyen (2) & À améliorer (1) \\
\hline
Contenu et arguments & & & \\
\hline
Correction de la langue & & & \\
\hline
Originalité & & & \\
\hline
Figures de style identifiées et produites & & & \\
\hline
\end{tabularx}
\end{center}

\courssub{Ressources à mobiliser}

\begin{aretenir}[Boîte à outils de l'émission]
\textbf{Lexique poétique :} \esp{el verso, la estrofa, la rima (asonante/consonante), el poeta, la musa, el cuarteto, la figura retórica}.

\textbf{Figures de style :} \esp{metáfora} (identification : « la hamburguesa es una bomba »), \esp{comparación} (avec \esp{como}), \esp{personificación} (« la manzana me sonríe »), \esp{hipérbole} (exagération), \esp{aliteración} (répétition de sons), \esp{anáfora} (répétition en début de vers).

\textbf{Structures d'opinion :} \esp{opino que + indicativo ; no creo que + subjuntivo ; me parece que ; estoy a favor de / en contra de ; desde mi punto de vista}.

\textbf{Connecteurs :} \esp{en primer lugar, además, sin embargo, por eso, en conclusión}.

\textbf{Lexique du thème :} \esp{la alimentación sana, la comida basura, los nutrientes, las vitaminas, la obesidad, la diabetes, los productos frescos, la grasa, el azúcar, cocinar, la dieta equilibrada}.
\end{aretenir}

\begin{attention}[Pièges à éviter pendant l'émission]
\begin{itemize}
\item Ne dites pas \esp{estoy de acuerdo con que...} : la structure correcte est \esp{estar de acuerdo en que + indicativo} ou \esp{estar de acuerdo con + nom/personne}.
\item Après \esp{no creo que}, n'oubliez pas le \cle{subjonctif} : \esp{No creo que la comida basura sea buena.} (et non \esp{es buena}).
\item Attention au faux ami : \esp{la dieta} signifie « le régime alimentaire » au quotidien, pas seulement « le régime amaigrissant ».
\item Prononcez bien \esp{¿Qué hay en tu plato?} avec le point d'interrogation ouvrant.
\end{itemize}
\end{attention}

\courssub{Proposition de production et corrigé commenté}

\begin{exoresolu}[Émission modèle de l'équipe « Verde Vida »]
Rédigez le script complet de l'équipe « Verde Vida » : présentation du fragment poétique, deux arguments d'experts, une réplique, et le quatrain final avec deux figures de style.
\tcblower
\textbf{Script modèle :}

\esp{Buenas tardes, queridos telespectadores. Soy la periodista cultural del equipo Verde Vida. Voy a leer la segunda estrofa del poema El duelo del plato. Esta estrofa tiene cuatro versos con rima asonante : es un cuarteto. El poeta utiliza una personificación muy expresiva : « la zanahoria grita su verdad », porque da cualidades humanas a un vegetal para defender la comida natural.}

\emph{(« Cette strophe a quatre vers à rime assonante : c'est un quatrain. Le poète utilise une personnification très expressive... » — le journaliste mobilise le lexique poétique : \esp{estrofa, verso, rima asonante, cuarteto, personificación}.)}

\esp{Como nutricionista, opino que los jóvenes cameruneses deben comer más productos frescos del mercado. En primer lugar, las frutas y las verduras contienen vitaminas que protegen el cuerpo. Además, estoy convencida de que la comida basura provoca enfermedades como la obesidad y la diabetes. No creo que una hamburguesa diaria sea inofensiva : es una bomba de grasa y de azúcar.}

\emph{(Notez : \esp{opino que} + indicatif ; \esp{no creo que sea} + subjonctif ; la métaphore \esp{una bomba de grasa} ; les connecteurs \esp{en primer lugar, además}.)}

\esp{Réplica al equipo Sabor Rápido : ustedes dicen que la comida rápida es barata, pero no estoy de acuerdo en que el precio justifique la enfermedad. Un plato de frijoles con plátano cuesta poco y alimenta mucho mejor.}

\emph{(La réplique utilise \esp{no estoy de acuerdo en que} + indicatif, structure correcte, et un exemple local : les haricots et le plantain.)}

\textbf{Quatrain final :}

\esp{El mango es oro del mercado,\\
la hamburguesa, lenta trampa ;\\
como un sol, el plato amado\\
da vida y nunca la apaga.}

\emph{(Figures : \esp{el mango es oro} = métaphore ; \esp{como un sol} = comparaison ; \esp{lenta trampa} = allitération en « l/t ». Rimes assonantes : mercado/amado, trampa/apaga.)}
\end{exoresolu}

\courssub{Bilan de l'activité}

Cette activité d'intégration a permis de réunir, dans une situation de communication réaliste et motivante, l'ensemble des savoir-faire de la leçon :

\begin{itemize}
\item \textbf{Compréhension et analyse poétique :} lire à voix haute, décrire une strophe avec le vocabulaire technique (\esp{verso, rima, cuarteto}) et expliquer l'effet d'une figure de style ;
\item \textbf{Production orale argumentée :} défendre une position sur l'alimentation avec des structures d'opinion, des connecteurs logiques et des exemples ancrés dans la réalité camerounaise ;
\item \textbf{Créativité langagière :} passer de l'identification des figures de style à leur \emph{production} dans un quatrain original ;
\item \textbf{Évaluation par les pairs :} développer l'esprit critique en jugeant le contenu, la langue et l'originalité selon une grille commune.
\end{itemize}

Pour prolonger l'activité à la maison, chaque élève peut rédiger un court article d'opinion (quinze lignes) intitulé \esp{« Mi plato, mi salud »}, en réutilisant au moins trois structures d'argumentation et une figure de style. Cet article pourra alimenter le mur d'expression espagnole de la classe pendant la Semaine de la santé.

\newpage

\coursec{Méthodes et exercices résolus}

\courssub{Découverte du poème support : « El duelo del plato »}

\begin{textebox}[Poème original — Thème : alimentation]
El plato vacío llora sobre la mesa, \\
mientras la fritura ríe con dientes de grasa. \\
La zanahoria, soldado caído en la guerra, \\
espera en vano que alguien la nombre en la fiesta.
\end{textebox}

\begin{exemplebox}[Traduction et glossaire]
\textbf{Traduction :} « L'assiette vide pleure sur la table, / pendant que la friture rit avec des dents de graisse. / La carotte, soldat tombé à la guerre, / attend en vain que quelqu'un la nomme à la fête. »

\textbf{Glossaire :} \esp{el plato} (assiette), \esp{vacío} (vide), \esp{llora} (pleure), \esp{la fritura} (la friture, le frit), \esp{la grasa} (la graisse), \esp{la zanahoria} (la carotte), \esp{el soldado} (le soldat), \esp{caído} (tombé), \esp{la guerra} (la guerre), \esp{espera} (attend), \esp{en vano} (en vain), \esp{nombrar} (nommer, citer), \esp{la fiesta} (la fête, ici : le repas de fête).
\end{exemplebox}

\courssub{Lexique de la poésie et analyse structurelle}

\begin{definition}[Vocabulaire essentiel de la poésie]
\begin{itemize}
\item \cle{El verso} : le vers (unité de base, une ligne).
\item \cle{La estrofa} : la strophe (groupe de vers séparés par une blanche).
\item \cle{La rima} : la rime (identité sonore à la fin des vers).
\item \cle{El poeta / la poeta} : le poète.
\item \cle{La musa} : la muse (source d'inspiration).
\item \cle{La métrica} : la métrique (étude de la mesure des vers).
\item \cle{La sinalefa} : fusion d'une voyelle finale et d'une voyelle initiale du vers suivant (compte pour 1 syllabe).
\end{itemize}
\end{definition}

\begin{propriete}[Mesure des vers (Métrique espagnole)]
\begin{enumerate}
\item Compter les syllabes \emph{phonétiques} (tenez compte de la \cle{sinalefa}).
\item Ajuster selon l'accent du dernier mot :
  \begin{itemize}
  \item Vers \cle{agudo} (accent sur la dernière syllabe : \esp{café}) $\rightarrow$ \textbf{+1} syllabe.
  \item Vers \cle{llano} (accent sur l'avant-dernière : \esp{casa}) $\rightarrow$ \textbf{inchangé}.
  \item Vers \cle{esdrújulo} (accent sur l'antépénultième : \esp{poesía}) $\rightarrow$ \textbf{-1} syllabe.
  \end{itemize}
\item Noms des vers courants : \esp{octosílabo} (8), \cle{endecasílabo} (11), \cle{alejandrino} (14). Vers de $\le$ 8 syllabes = \esp{arte menor} ; $\ge$ 9 = \esp{arte mayor}.
\end{enumerate}
\end{propriete}

\begin{propriete}[Types de rimes et schémas]
\begin{itemize}
\item \cle{Rima consonante} : identité de tous les sons à partir de la dernière voyelle accentuée (\esp{corazón / canción}).
\item \cle{Rima asonante} : identité des voyelles seulement (\esp{mesa / grasa} $\rightarrow$ \emph{e-a} / \emph{a-a}).
\item \cle{Rima libre / verso suelto} : pas de rime.
\item Schéma : minuscules pour \esp{arte menor} (\esp{abab}), majuscules pour \esp{arte mayor} (\esp{ABBA}).
\end{itemize}
\end{propriete}

\begin{exoresolu}[Analyse structurelle de « El duelo del plato »]
\textbf{Énoncé.} Analysez la structure du poème (vers, strophes, mesure, rime).
\tcblower
\textbf{Solution détaillée.}

\textbf{1. Vers et strophes :} 4 vers, une seule strophe (\esp{cuarteto} ou quatrain).

\textbf{2. Mesure (comptage syllabique) :}
\begin{itemize}
\item V1 : \esp{El pla-to vá-cio llo-ra so-bre la me-sa} $\rightarrow$ sinalefa \esp{llora so-bre} (\emph{o-a}), \esp{sobre la} (\emph{e-a}). Comptage : 8 syllabes (\esp{octosílabo}, vers llano).
\item V2 : \esp{mien-tras la fri-tu-ra ríe con dien-tes de gra-sa} $\rightarrow$ sinalefa \esp{tras la} (\emph{a-a}), \esp{fritura ríe} (\emph{a-í}), \esp{dientes de} (\emph{e-e}). Comptage : 8 syllabes (\esp{octosílabo}, vers llano).
\item V3 : \esp{La za-na-ho-ria, sol-da-do ca-í-do en la gue-rra} $\rightarrow$ sinalefa \esp{zanahoria, soldado} (\emph{a-o}? non, consonne), \esp{caído en} (\emph{o-e}), \esp{en la} (\emph{e-a}). Comptage : 11 syllabes (\cle{endecasílabo}, vers llano). \textbf{Attention :} ce vers est plus long (arte mayor), ce qui rompt la régularité populaire pour souligner la gravité.
\item V4 : \esp{es-pe-ra en va-no que al-guien la nom-bre en la fi-es-ta} $\rightarrow$ sinalefa \esp{espera en} (\emph{a-e}), \esp{vano que} (\emph{o-e}), \esp{alguien la} (\emph{e-a}), \esp{nombre en} (\emph{e-e}), \esp{en la} (\emph{e-a}). Comptage : 11 syllabes (\cle{endecasílabo}, vers llano).
\end{itemize}

\textbf{3. Rime :} V1 \esp{mesa} / V2 \esp{grasa} $\rightarrow$ \cle{rima asonante} en \emph{a} (voyelles \emph{e-a} / \emph{a-a}). V3 \esp{guerra} / V4 \esp{fiesta} $\rightarrow$ pas de rime (\esp{versos sueltos}). Schéma : \textbf{abcc} (ou \esp{ab} libre).

\textbf{4. Interprétation :} L'alternance octosyllabes (vers courts, rythme rapide, monde de la friture) / endecasylabes (vers longs, solennels, monde de la carotte/soldat) marque le contraste entre la légèreté trompeuse de la malbouffe et la tragédie silencieuse de la santé négligée.
\end{exoresolu}

\courssub{Figures de style : identification, définition et effets}

\begin{aretenir}[Les 6 figures au programme — Tableau de révision]
\begin{center}
\begin{tabularx}{\linewidth}{|X|X|X|X|}
\hline
\rowcolor{popblueL}
\textbf{Figure} & \textbf{Définition / Test} & \textbf{Marqueurs} & \textbf{Effet principal} \\
\hline
\cle{Métaphore} & Identification directe A = B (sans outil). Test : \esp{ser} & Absence de \esp{como} & Image forte, condensation du sens \\
\hline
\cle{Comparación} (Símil) & Comparaison explicite A comme B & \esp{como}, \esp{cual}, \esp{parecido a} & Mise en parallèle, clarté, parfois humour \\
\hline
\cle{Personificación} (Prosopopeya) & Attributs humains à non-humain & Verbes/adjectifs humains & Émotion, animation du monde, empathie \\
\hline
\cle{Hipérbole} & Exagération volontaire & Mots excessifs (\esp{mil}, \esp{nunca}, \esp{eternidad}) & Intensité, drame ou comique \\
\hline
\cle{Aliteración} & Répétition de sons consonantiques & Sons répétés (\esp{r, s, t, p}) & Musicalité, rythme, suggestion sonore \\
\hline
\cle{Anáfora} & Répétition en \emph{début} de vers/phrase & Même mot/groupe initial & Insistance, rythme martelé, force persuasive \\
\hline
\end{tabularx}
\end{center}
\end{aretenir}

\begin{methode}[Reconnaître et commenter une figure : la règle des 3 temps]
\begin{enumerate}
\item \textbf{Nommer} la figure avec son terme exact en espagnol.
\item \textbf{Citer} le segment précis (guillemets \esp{« ... »}) et justifier (ex : « absence de \esp{como} », « présence de \esp{como} », « verbe humain appliqué à... »).
\item \textbf{Expliquer l'effet} : que \emph{ressent} le lecteur ? que \emph{comprend}-il de plus ? (émotion, image, rythme, insistance, critique).
\end{enumerate}
\textbf{Interdit :} dire « c'est une métaphore » sans citation ni effet. Au Bac, 0 point si l'effet manque.
\end{methode}

\begin{exoresolu}[Identification et commentaire de figures]
\textbf{Énoncé.} Pour chaque vers, identifiez la figure, citez la preuve et expliquez l'effet.

1. \esp{El hambre es un monstruo que devora la ciudad.}
2. \esp{El niño come como un león después del partido.}
3. \esp{La luna miraba triste los platos vacíos.}
4. \esp{He esperado mil años en la cola del restaurante.}
5. \esp{Pan para el pobre, pan para el preso, pan para el peregrino.}
\tcblower
\textbf{Solution détaillée.}

1. \textbf{Métaphore}. Preuve : \esp{El hambre es un monstruo} (verbe \esp{ser}, pas de \esp{como}). Effet : la faim devient une bête terrifiante, vivante, qui menace la collectivité ; elle n'est plus un manque mais un prédateur.

2. \textbf{Comparación (símil)}. Preuve : outil \esp{como} (\esp{come como un león}). Effet : image de voracité animale, puissance, mais aussi naturel ; adoucit le jugement sur l'appétit de l'enfant.

3. \textbf{Personificación}. Preuve : \esp{miraba} (regarder) + \esp{triste} (sentiment) appliqués à \esp{la luna}. Effet : l'univers entier partage la détresse des affamés ; atmosphère mélancolique, solidarité cosmique.

4. \textbf{Hipérbole}. Preuve : \esp{mil años} (impossible physiquement). Effet : exprime l'exaspération, le temps subjectif dilaté par la faim ou l'ennui ; intensité émotionnelle.

5. \textbf{Anáfora}. Preuve : répétition de \esp{Pan para} en tête de chaque proposition. Effet : rythme de litanie, revendication rythmée, universalité du droit au pain (pauvre, prisonnier, pèlerin = humanité entière).
\end{exoresolu}

\courssub{Méthodologie du commentaire de poème et commentaire modèle}

\begin{methode}[Commenter un poème : la démarche en 5 étapes (Bac)]
\begin{enumerate}
\item \textbf{Lecture expressive} : lire 2 fois à voix haute (plaisir + repérage mots-clés, répétitions, images).
\item \textbf{Thème vs Sujet} : \esp{El tema es...} (ce que le poète \emph{dit} du sujet). Ne pas confondre : Sujet = « la nourriture » ; Thème = « critique de la malbouffe et défense du naturel ».
\item \textbf{Structure \& Versification} : nb vers, strophes, mesure (octosílabo, endecasílabo...), type de rime (consonante/asonante/libre), schéma. \emph{Relier forme et sens}.
\item \textbf{Figures de style} : appliquer la règle des 3 temps (Nommer $\rightarrow$ Citer $\rightarrow$ Effet). Minimum 3 figures distinctes.
\item \textbf{Interprétation personnelle \& Contexte} : \esp{En mi opinión...}, \esp{Me parece que...} + justification (\esp{porque}, \esp{ya que}). Lier au contexte camerounais/hispanophone (Yaoundé, Guinée équatoriale, actualité).
\end{enumerate}
\textbf{Rédaction :} Présent de vérité générale (\esp{el poeta denuncia}, \esp{los versos expresan}). Connecteurs : \esp{en primer lugar}, \esp{además}, \esp{por otra parte}, \esp{en conclusión}. Citer \emph{toujours} les vers.
\end{methode}

\begin{exoresolu}[Commentaire guidé complet de « El duelo del plato »]
\textbf{Énoncé.} Rédigez un commentaire structuré du poème (environ 15 lignes).
\tcblower
\textbf{Commentaire modèle.}

\esp{El poema « El duelo del plato » trata de la oposición entre la comida basura y la alimentación natural. El tema principal es la denuncia de los malos hábitos alimenticios: el ser humano prefiere lo grasiento a lo sano.}

\esp{En primer lugar, la estructura refuerza el mensaje: los dos primeros versos son octosílabos con rima asonante (mesa/grasa), ritmo ligero y popular que evoca el mundo de la fritura; los dos siguientes son endecasílabos sin rima, más graves, para el mundo de la zanahoria-soldado.}

\esp{Además, las figuras de estilo dramatizan el conflicto. La personificación « el plato vacío llora » humaniza el objeto para conmover al lector ante el desperdicio. La metáfora « la zanahoria, soldado caído en la guerra » transforma el plato en campo de batalla: la comida sana es una víctima inocente. Finalmente, la hipérbole « guerra » magnifica la banalidad de una comida.}

\esp{En mi opinión, este poema es muy actual en Camerún: en los mercados de Yaoundé o Douala, los jóvenes eligen a menudo frituras y refrescos azucarados. El texto nos invita a reflexionar sobre nuestra salud y me parece una herramienta pedagógica eficaz.}
\end{exoresolu}

\courssub{Lexique du débat : alimentation saine vs malbouffe}

\begin{definition}[Vocabulaire thématique — Alimentation et santé]
\begin{center}
\begin{tabularx}{\linewidth}{|X|X|X|}
\hline
\rowcolor{popgreenL}
\textbf{Français} & \textbf{Español} & \textbf{Exemple / Collocation} \\
\hline
Alimentation équilibrée & \esp{la alimentación equilibrada} & \esp{llevar una alimentación equilibrada} \\
Malbouffe / Nourriture poubelle & \esp{la comida basura} / \esp{la comida chatarra} & \esp{evitar la comida basura} \\
Aliments gras / sucrés / salés & \esp{los alimentos grasos / azucarados / salados} & \esp{reducir el consumo de azucarados} \\
Fruits et légumes frais & \esp{las frutas y verduras frescas} & \esp{comer cinco raciones al día} \\
Obésité / Surpoids & \esp{la obesidad} / \esp{el sobrepeso} & \esp{prevenir la obesidad infantil} \\
Maladies cardiovasculaires / diabète & \esp{las enfermedades cardiovasculares / la diabetes} & \esp{riesgo de diabetes tipo 2} \\
Nuisible / Bénéfique pour la santé & \esp{perjudicial / beneficioso para la salud} & \esp{es perjudicial para el corazón} \\
Habitudes alimentaires & \esp{los hábitos alimenticios / alimentarios} & \esp{cambiar los hábitos alimenticios} \\
Sédentarité & \esp{el sedentarismo} & \esp{Combatir el sedentarismo} \\
Étiquetage nutritionnel & \esp{el etiquetado nutricional} & \esp{leer el etiquetado} \\
\hline
\end{tabularx}
\end{center}
\end{definition}

\begin{savaistu}[Contexte camerounais et Guinée équatoriale]
La Guinée équatoriale (pays hispanophone voisin) fait face à une transition nutritionnelle rapide : coexistence de la malnutrition (zones rurales) et de l'obésité (Malabo, Bata) due à l'importation d'aliments ultra-transformés. Au Cameroun, le \esp{« bendskin »} (danse populaire) brûle des calories, mais la vente de \esp{« beignets-haricots »}, \esp{« poisson braisé »} et boissons sucrées devant les lycées pose un vrai problème de santé publique. Le programme \esp{« NANA »} (Nutrición y Alimentación para Niños y Adolescentes) de l'OMS inspire les cantines scolaires à Yaoundé.
\end{savaistu}

\courssub{Grammaire \& Traduction : structures d'argumentation et d'opinion}

\begin{propriete}[Structures d'opinion : Indicatif vs Subjonctif]
\begin{center}
\begin{tabularx}{\linewidth}{|X|X|X|}
\hline
\rowcolor{poporangeL}
\textbf{Français} & \textbf{Espagnol (Structure)} & \textbf{Mode / Remarque} \\
\hline
Je pense que / Je crois que + ind. & \esp{Pienso que / Creo que + indicativo} & \textbf{Indicatif} (certitude subjective) \\
\hline
Je ne pense pas que / Je ne crois pas que & \esp{No pienso que / No creo que + subjuntivo} & \textbf{Subjonctif} (doute/négation) \\
\hline
Il faut / Il est nécessaire que & \esp{Es necesario que + subjuntivo} & \textbf{Subjonctif} (obligation/nécessité) \\
\hline
Il faut + infinitif (général) & \esp{Hay que + infinitivo} & \textbf{Infinitif} (impersonnel) \\
\hline
Il est bon/mauvais que & \esp{Es bueno/malo que + subjuntivo} & \textbf{Subjonctif} (évaluation) \\
\hline
Je veux que / J'exige que & \esp{Quiero que / Exijo que + subjuntivo} & \textbf{Subjonctif} (volonté/influence) \\
\hline
\end{tabularx}
\end{center}
\end{propriete}

\begin{attention}[Pièges de traduction — Thème/Version]
\begin{enumerate}
\item \textbf{« Bon pour la santé »} $\rightarrow$ \esp{bueno \textbf{para} la salud} (jamais \esp{por}). \esp{Para} = finalité/bénéficiaire.
\item \textbf{« De plus en plus »} $\rightarrow$ \esp{cada vez más} + adjectif/nom (\esp{cada vez más jovenes}, \esp{cada vez más grasa}). Pas \esp{más y más}.
\item \textbf{« Prendre position »} $\rightarrow$ \esp{tomar posición} / \esp{posicionarse} / \esp{decantarse por}. Pas \esp{prendre position} calqué.
\item \textbf{« Malbouffe »} $\rightarrow$ \esp{comida basura} (Espagne) / \esp{comida chatarra} (Amérique). Éviter \esp{comida mala} (incorrect).
\item \textbf{Accords :} \esp{la salud} (fém.) $\rightarrow$ \esp{buena salud} ; \esp{el azúcar} (masc.) mais \esp{las bebidas azucaradas} (fém. pl.).
\item \textbf{Subjonctif après \esp{no creer que} :} \esp{No creo que \textbf{sea} sano} (pas \esp{es}). \textbf{Erreur éliminatoire au Bac.}
\end{enumerate}
\end{attention}

\begin{exoresolu}[Thème : phrases d'opinion sur l'alimentation]
\textbf{Énoncé.} Traduisez en espagnol correct (soignez accents, mode, prépositions).

1. Je pense que la malbouffe est dangereuse pour la santé des adolescents.
2. Je ne crois pas que les frites soient un bon repas.
3. Il faut manger des fruits frais tous les jours.
4. Il faut que le gouvernement protège les consommateurs.
5. Les jeunes de Douala boivent de plus en plus de boissons sucrées.
\tcblower
\textbf{Solution détaillée.}

1. \esp{Pienso que la comida basura \textbf{es} peligrosa \textbf{para} la salud de los adolescentes.} — \esp{Pienso que} + \textbf{indicatif} (\esp{es}). \esp{Peligrosa} (accord fém. sg). \esp{Para} la santé.

2. \esp{No creo que las patatas fritas \textbf{sean} una buena comida.} — \esp{No creo que} + \textbf{subjonctif présent} (\esp{sean}, 3\textsuperscript{e} pl. de \esp{ser}). \esp{Patatas fritas} (standard péninsulaire) ou \esp{papas fritas} (Amérique).

3. \esp{\textbf{Hay que} comer fruta fresca todos los días.} — Impersonnel \esp{hay que} + infinitif. \esp{Fruta fresca} (collectif singulier naturel) ou \esp{frutas frescas}.

4. \esp{Es necesario que el gobierno \textbf{proteja} \textbf{a} los consumidores.} — \esp{Es necesario que} + \textbf{subjonctif} (\esp{proteja}, 3\textsuperscript{e} sg. de \esp{proteger}). \esp{A} personnelle (personnes déterminées).

5. \esp{Los jóvenes de Douala beben \textbf{cada vez más} bebidas \textbf{azucaradas}.} — \esp{Cada vez más} + nom. \esp{Azucaradas} (accord fém. pl., accent sur \emph{á}).
\end{exoresolu}

\courssub{Production guidée : débat oral et rédaction d'un article d'opinion}

\begin{methode}[Débattre et rédiger un article d'opinion : structure type]
\begin{enumerate}
\item \textbf{Accroche / Titre} : clair, percutant (\esp{¿Comida basura en el recreo? ¡Basta ya!}).
\item \textbf{Thèse (Position nette)} : \esp{Estoy a favor de...} / \esp{Me opongo a...} / \esp{Desde mi punto de vista, los colegios deben...}
\item \textbf{Arguments (min. 2)} : \esp{En primer lugar... porque...} ; \esp{Además... ya que...} ; \esp{Por otra parte...}. Appuyer par exemples (\esp{por ejemplo}, \esp{como el caso de...}).
\item \textbf{Concession / Réfutation} : \esp{Es cierto que... pero...} / \esp{Aunque algunos argumenten que..., lo cierto es que...}.
\item \textbf{Conclusion forte} : \esp{En conclusión / Por todo ello / En definitiva}, reformuler la thèse + ouverture (\esp{proteger el futuro}, \esp{educar para la vida}).
\end{enumerate}
\textbf{Connecteurs de haut niveau :} \esp{asimismo} (de plus), \esp{no obstante} (cependant), \esp{por consiguiente} (par conséquent), \esp{en definitiva} (en définitive).
\end{methode}

\begin{exoresolu}[Rédaction d'un article d'opinion pour le journal scolaire]
\textbf{Énoncé.} Écrivez 10--12 lignes : \esp{« Los colegios deben prohibir la venta de comida basura en los recreos »}. 3 connecteurs + 2 structures d'opinion minimum.
\tcblower
\textbf{Proposition de rédaction (niveau Bac excellent).}

\esp{¿Hasta cuándo permitiremos que la comida basura invada nuestros recreos? En mi opinión, los colegios deben prohibir su venta de manera inmediata. En primer lugar, los ultraprocesados son perjudiciales para la salud, ya que disparan la obesidad y la diabetes juvenil. Además, un alumno mal alimentado rinde menos: la concentración y el rendimiento exigen nutrientes reales, no calorías vacías. Es cierto que la comida basura es barata y accesible, pero no obstante, el coste sanitario futuro es inasumible. Por ejemplo, en varios liceos de Yaoundé ya se ofrecen frutas de temporada a precio coste. En conclusión, prohibir la comida basura no es una restricción, es una inversión en el futuro de nuestra juventud.}

\textbf{Analyse grammaticale du modèle :}
\begin{itemize}
\item \esp{¿Hasta cuándo...?} : question rhétorique (accroche).
\item \esp{deben prohibir} : \esp{deber} + inf. (obligation morale).
\item \esp{sean / disparen / rinde / exijan} : indicatif (faits constatés après \esp{ya que}, \esp{porque}).
\item \esp{es cierto que... pero no obstante} : concession + opposition forte (\esp{no obstante} = toutefois).
\item \esp{se ofrecen} : passif réfléchi (\esp{se} + 3\textsuperscript{e} pl.), accord \esp{frutas}.
\item \esp{no es una restricción, es una inversión} : parallélisme, style affirmé.
\end{itemize}
\end{exoresolu}

\begin{experience}[Débat oral guidé : « Alimentación sana en el instituto »]
\textbf{Consigne.} En groupes de 4. Rôles : \textbf{A} (Directeur du lycée), \textbf{B} (Parent d'élève), \textbf{C} (Élève), \textbf{D} (Vendeur de cantine).
\textbf{Sujet :} \esp{« ¿Debe el instituto prohibir totalmente la venta de refrescos y frituras en el recreo? »}.
\textbf{Préparation (5 min) :} Chaque rôle note 2 arguments avec connecteurs (\esp{en primer lugar}, \esp{porque}, \esp{aunque}).
\textbf{Débat (10 min) :} Tour de table, le Directeur synthétise et prend une décision finale.
\textbf{Évaluation :} Utilisation du subjonctif (\esp{es necesario que...}, \esp{no creo que...}), vocabulaire précis (\esp{azucarados}, \esp{obesidad}, \esp{rendimiento}), fluidité.
\end{experience}

\begin{exercice}[Entraînement écrit — Expression d'opinion]
\textbf{Consigne.} Choisissez UN sujet. Rédigez 12--15 lignes. Respectez la structure : Thèse -- 2 Arguments justifiés -- Concession/Réfutation -- Conclusion.
\begin{enumerate}
\item \esp{« El gobierno debe gravar con impuestos altos las bebidas azucaradas. »}
\item \esp{« La educación nutricional debería ser una asignatura obligatoria en la secundaria. »}
\item \esp{« La publicidad de comida basura dirigida a niños debe ser ilegal. »}
\end{enumerate}
\textbf{Critères :} 3 connecteurs variés, 1 subjonctif d'opinion (\esp{no creo que / es necesario que}), 5 mots du lexique thématique, accents parfaits.
\end{exercice}

\begin{corrige}[Exercice — Pistes de correction pour le sujet 1]
\textbf{Thèse :} \esp{Estoy totalmente a favor de gravar las bebidas azucaradas.}
\textbf{Arg 1 :} \esp{En primer lugar, el impuesto reduce el consumo, como demuestra México desde 2014 (bajada del 7\% en dos años).}
\textbf{Arg 2 :} \esp{Además, los ingresos financian campañas de salud y deporte escolar.}
\textbf{Concession :} \esp{Es cierto que afecta a las rentas bajas, pero no obstante, son las que más sufren la diabetes.}
\textbf{Conclusion :} \esp{Por consiguiente, el impuesto es una medida de justicia sanitaria.}
\end{corrige}

\courssub{Erreurs à éviter et points de vigilance (Bac)}

\begin{attention}[Les 5 erreurs qui coûtent des points au Bac]
\begin{enumerate}
\item \textbf{Accents oubliés :} \esp{poesía}, \esp{alegría}, \esp{mañana}, \esp{corazón}, \esp{azucarado}, \esp{opinión}, \esp{país}, \esp{también}. Un accent change le sens (\esp{hablo} vs \esp{habló}) ou marque l'intensité (\esp{fría} vs \esp{fria} incorrect).
\item \textbf{Métaphore vs Comparación :} Présence de \esp{como} = \textbf{comparación} (jamais métaphore). Absence de \esp{como} + \esp{ser} = \textbf{métaphore}. Erreur classique = 0 pt sur la figure.
\item \textbf{Subjonctif après négation d'opinion :} \esp{No creo que \textbf{sea} bueno} (subj.). \esp{Creo que \textbf{es} bueno} (ind.). \textbf{Test systématique} avant d'écrire.
\item \textbf{Calques français :} \textbf{Pour la santé} = \esp{para la salud} (pas \esp{por}). \textbf{De plus en plus} = \esp{cada vez más} (pas \esp{más y más}). \textbf{Prendre position} = \esp{tomar posición / posicionarse}. \textbf{Réaliser} (prendre conscience) = \esp{darse cuenta de} (pas \esp{realizar} = accomplir).
\item \textbf{Accords \& Faux amis :} \esp{la salud} (fém.) $\rightarrow$ \esp{buena}. \esp{el problema} (masc.) $\rightarrow$ \esp{grave}. \esp{actualmente} = « actuellement » (pas « en fait »). \esp{asistir} = « assister à » (pas « aider » = \esp{ayudar}). \esp{éxito} = « succès » (pas « sortie » = \esp{salida}).
\end{enumerate}
\end{attention}

\begin{savaistu}[Poésie et engagement social dans l'espace hispanophone]
De \esp{Antonio Machado} (\esp{« Caminante, no hay camino, se hace camino al andar »}) à \esp{Gioconda Belli} (Nicaragua, poésie féministe et révolutionnaire), la poésie hispanophone a toujours été une arme de conscience sociale. En Guinée équatoriale, \esp{Juan Tomás Ávila Laurel} utilise le vers pour dénoncer la dictature et l'exil. Au Cameroun, les \esp{« slameurs »} bilingues (français/espagnol/anglais) comme \esp{« Koppo »} ou collectifs \esp{« Slamouraï »} investissent les scènes de Yaoundé pour parler nutrition, identité et jeunesse. La poésie n'est pas un ornement : c'est un \cle{acto de comunicación} politique et éthique.
\end{savaistu}

\newpage

\coursec{Exercices}

\courssub{Je vérifie mes connaissances}

\begin{exercice}[QCM : la poésie et ses figures]
Entourez la bonne réponse.
\begin{enumerate}
\item Un \esp{verso} est : a) une strophe ; b) un vers ; c) une rime.
\item La \esp{personificación} consiste à : a) comparer avec \esp{como} ; b) attribuer des qualités humaines à un objet ou à une idée ; c) exagérer.
\item Dans le vers \esp{« las papas fritas, fieras, me persiguen »}, on trouve surtout : a) une aliteración ; b) una hipérbole y una personificación ; c) una anáfora.
\item La répétition d'un même mot en début de plusieurs vers s'appelle : a) la metáfora ; b) la rima ; c) la anáfora.
\item Une \esp{estrofa} de quatre vers s'appelle : a) un terceto ; b) un cuarteto ; c) un soneto.
\end{enumerate}
\end{exercice}

\begin{exercice}[Vrai ou faux, à justifier]
Répondez par vrai ou faux, puis justifiez chaque réponse en une phrase en français.
\begin{enumerate}
\item La \esp{metáfora} utilise toujours l'outil de comparaison \esp{como}.
\item L'\esp{aliteración} est la répétition de sons identiques ou proches dans un même vers.
\item \esp{« Es tan grande como una montaña »} est une \esp{comparación}.
\item Dans un débat, \esp{« estoy a favor de »} permet d'exprimer une opposition.
\item Le poème « El duelo del plato » oppose la nourriture saine et la malbouffe.
\end{enumerate}
\end{exercice}

\begin{exercice}[Vocabulaire : poésie et alimentation]
Complétez chaque phrase avec le mot espagnol qui convient, choisi dans la liste : \esp{poeta} ; \esp{rima} ; \esp{estrofa} ; \esp{comida basura} ; \esp{alimentación sana} ; \esp{musa}.
\begin{enumerate}
\item \esp{El \ldots\ldots\ldots\ldots escribe poemas.}
\item \esp{« corazón » y « canción » forman una \ldots\ldots\ldots\ldots.}
\item \esp{Las hamburguesas y los refrescos azucarados son \ldots\ldots\ldots\ldots.}
\item \esp{Comer frutas, verduras y pescado forma parte de una \ldots\ldots\ldots\ldots.}
\item \esp{El poema tiene cuatro \ldots\ldots\ldots\ldots de cinco versos.}
\end{enumerate}
\end{exercice}

\begin{exercice}[Conjugaison : le subjonctif de l'opinion]
Complétez avec le verbe au subjonctif présent.
\begin{enumerate}
\item \esp{Es importante que nosotros \ldots\ldots\ldots\ldots (comer) más verduras.}
\item \esp{No creo que la comida basura \ldots\ldots\ldots\ldots (ser) buena para la salud.}
\item \esp{Es necesario que el gobierno \ldots\ldots\ldots\ldots (prohibir) esa publicidad.}
\item \esp{Dudo que los niños \ldots\ldots\ldots\ldots (saber) elegir bien.}
\item \esp{Aunque las verduras \ldots\ldots\ldots\ldots (parecer) menos atractivas, hay que comerlas.}
\end{enumerate}
\end{exercice}

\courssub{Je m'entraîne}

\begin{exercice}[Thème : phrases à traduire]
Traduisez en espagnol.
\begin{enumerate}
\item Bien que les légumes soient moins attractifs que les frites, il est indispensable que nous les mangions pour rester en bonne santé.
\item Je ne crois pas que la malbouffe soit moins chère qu'un repas préparé à la maison.
\item Il faut que les écoles interdisent les boissons sucrées dans les cantines.
\item À mon avis, il est urgent que les parents cuisinent avec leurs enfants.
\end{enumerate}
\end{exercice}

\begin{exercice}[Transformations : du style direct à l'argumentation]
Reprenez chaque opinion en la rendant plus nuancée ou plus forte, selon la consigne entre parenthèses. Utilisez la structure imposée.
\begin{enumerate}
\item \esp{La comida basura es mala.} (nuancer avec \esp{a mi juicio} + subjonctif possible)
\item \esp{Los jóvenes comen mal.} (renforcer avec \esp{es evidente que} + indicatif)
\item \esp{Hay que cocinar en casa.} (transformer avec \esp{es necesario que} + subjonctif, sujet « nosotros »)
\item \esp{La publicidad influye en los niños.} (exprimer le doute avec \esp{no creo que})
\end{enumerate}
\end{exercice}

\begin{exercice}[Texte à trous : un article sur la cantine]
Complétez le texte avec les mots suivants : \esp{aunque} ; \esp{sin embargo} ; \esp{por eso} ; \esp{además} ; \esp{en conclusión}.
\begin{textebox}[La cantine du lycée]
\esp{En nuestro liceo de Yaoundé, muchos alumnos compran frituras a la salida, \ldots\ldots\ldots\ldots la cantina ofrece platos equilibrados. Los fritos son baratos y rápidos; \ldots\ldots\ldots\ldots, contienen demasiada grasa. \ldots\ldots\ldots\ldots, no aportan las vitaminas que necesitan los adolescentes. \ldots\ldots\ldots\ldots, el director quiere crear un club de cocina saludable. \ldots\ldots\ldots\ldots, comer bien es también aprender a vivir bien.}
\end{textebox}
\end{exercice}

\begin{exercice}[Appariements : figures de style]
Associez chaque vers du poème « El duelo del plato » à la figure de style dominante : a) metáfora ; b) comparación ; c) personificación ; d) hipérbole ; e) aliteración ; f) anáfora.
\begin{enumerate}
\item \esp{« La hamburguesa me sonríe con dientes de grasa »}
\item \esp{« El brócoli es un soldado verde de mi cuerpo »}
\item \esp{« Dulce, dulce, dulce trampa dorada »}
\item \esp{« Podría comer mil montañas de papas fritas »}
\item \esp{« La manzana cruje como un tambor de la mañana »}
\item \esp{« No quiero azúcar, no quiero grasa, no quiero mentiras »}
\end{enumerate}
\end{exercice}

\courssub{Je me prépare au Bac}

\begin{exercice}[Compréhension et commentaire de poème]
Lisez le poème suivant, puis répondez aux questions en français.
\begin{textebox}[El duelo del plato (extraits)]
\esp{En la esquina del mercado brilla la trampa dorada,\\
la hamburguesa me sonríe con dientes de grasa.\\
Podría comer mil montañas de papas fritas,\\
dulce, dulce, dulce veneno que nunca se agota.\\[0.5em]
Pero el brócoli es un soldado verde de mi cuerpo,\\
la manzana cruje como un tambor de la mañana.\\
No quiero azúcar, no quiero grasa, no quiero mentiras :\\
mi plato es mi escudo, mi cocina, mi bandera.}
\end{textebox}
\begin{enumerate}
\item Présentez le poème : thème, structure (nombre de strophes et de vers), ton général. (4 points)
\item Relevez dans les vers 1 à 4 deux figures de style, nommez-les et expliquez leur effet. (4 points)
\item Analysez les vers 5 à 8 : comment le poète utilise-t-il la \esp{personificación} et la \esp{metáfora} pour valoriser l'alimentation saine ? (4 points)
\item Expliquez le titre du poème : pourquoi un « duelo » ? (2 points)
\item Traduisez en français les deux derniers vers. (2 points)
\end{enumerate}
\end{exercice}

\begin{exercice}[Thème et version]
\textbf{A. Thème.} Traduisez en espagnol. (10 points)
\begin{enumerate}
\item Bien que les légumes soient moins attractifs que les frites, il est indispensable que nous les mangions pour rester en bonne santé.
\item Il est temps que le gouvernement interdise la publicité pour la malbouffe pendant les programmes pour enfants.
\item Plus nous cuisinons à la maison, moins nous tombons malades.
\end{enumerate}
\textbf{B. Version.} Traduisez en français le texte suivant. (10 points)
\begin{textebox}[Article de presse]
\esp{Según los médicos, la obesidad infantil ha aumentado en muchos países durante los últimos años. Los expertos afirman que es urgente que las familias cambien sus hábitos : menos refrescos azucarados, más frutas y más deporte. Aunque las campañas de prevención son cada vez más numerosas, todavía no se ha conseguido que los niños prefieran una manzana a una bolsa de patatas fritas.}
\end{textebox}
\end{exercice}

\begin{exercice}[Expression écrite et orale]
\textbf{A. Expression écrite.} Rédigez en espagnol une lettre ouverte au ministre de la Santé pour proposer l'interdiction de la publicité pour la malbouffe destinée aux enfants. Vous utiliserez un registre formel, au moins quatre connecteurs argumentatifs (\esp{en primer lugar}, \esp{además}, \esp{sin embargo}, \esp{por lo tanto}...), deux structures avec subjonctif et un vocabulaire précis de la santé. Longueur attendue : entre 180 et 220 mots. (12 points)

\textbf{B. Expression orale.} Préparez une prise de parole de deux minutes : présentez votre plat traditionnel camerounais préféré en expliquant pourquoi il représente une alimentation saine. Vous utiliserez au moins trois connecteurs, une métaphore et une structure d'opinion avec subjonctif. (8 points)
\end{exercice}

\newpage

\coursec{Corrigés des exercices}

\coursec{Corrigés des exercices}

\begin{corrige}[Exercice 1 — QCM : La poésie et ses figures]
\begin{enumerate}
\item \textbf{b) un verso.} Un \esp{verso} est une ligne unique de poésie ; une \esp{estrofa} est un ensemble de vers (strophe) et la \esp{rima} est la répétition de sons en fin de vers.
\item \textbf{b) attribuer des qualités humaines à un objet ou à une idée.} La \esp{personificación} (prosopopée) humanise l'inanimé ou l'abstrait. La comparaison utilise \esp{como} (a) et l'exagération est l'\esp{hipérbole} (c).
\item \textbf{b) una hipérbole y una personificación.} Le vers \esp{« Podría comer mil montañas de papas fritas »} (v. 3 du poème) contient une \esp{hipérbole} (\esp{mil montañas}) et, par extension, le vers \esp{« las papas fritas, fieras, me persiguen »} (v. 2) présente une \esp{personificación} (\esp{fieras}, \esp{me persiguen}). L'item 3 du QCM porte sur l'ensemble de la première strophe où ces deux figures dominent.
\item \textbf{c) la anáfora.} Répétition d'un mot ou groupe de mots en début de vers successifs. La \esp{metáfora} est une identification, la \esp{rima} une identité sonore finale.
\item \textbf{b) un cuarteto.} \esp{Terceto} = 3 vers, \esp{cuarteto} = 4 vers, \esp{soneto} = poème de 14 vers (structure fixe).
\end{enumerate}
\end{corrige}

\begin{corrige}[Exercice 2 — Vrai ou faux, à justifier]
\begin{enumerate}
\item \textbf{Faux.} La \esp{metáfora} identifie deux réalités sans outil de comparaison (ni \esp{como}, ni \esp{parece}, ni \esp{es como}). C'est la \esp{comparación} qui utilise \esp{como}.
\item \textbf{Vrai.} L'\esp{aliteración} est la répétition de sons consonantiques identiques ou similaires au début des mots ou à l'intérieur d'un même vers ou groupe de vers (ex. \esp{« dulce, dulce, dulce »}).
\item \textbf{Vrai.} La présence de \esp{como} établit explicitement une relation de similitude entre deux termes : c'est la définition de la \esp{comparación}.
\item \textbf{Faux.} \esp{« Estoy a favor de »} exprime un \textbf{accord}, un soutien. Pour exprimer une opposition, on utilise \esp{« estoy en contra de »}.
\item \textbf{Vrai.} Le poème met en scène un conflit (\esp{duelo}) entre l'attrait de la malbouffe (\esp{hamburguesa, papas fritas}) et les bienfaits de l'alimentation saine (\esp{brócoli, manzana}), le « je » poétique devant choisir.
\end{enumerate}
\end{corrige}

\begin{corrige}[Exercice 3 — Vocabulaire : poésie et alimentation]
\begin{enumerate}
\item \esp{El \textbf{poeta} escribe poemas.} (Auteur du poème).
\item \esp{« corazón » y « canción » forman una \textbf{rima}.} (Identité sonore finale \textit{-ón}).
\item \esp{Las hamburguesas y los refrescos azucarados son \textbf{comida basura}.} (Synonyme de \esp{comida chatarra}, malbouffe).
\item \esp{Comer frutas, verduras y pescado forma parte de una \textbf{alimentación sana}.} (Équilibre nutritionnel).
\item \esp{El poema \textbf{« El duelo del plato »} tiene dos \textbf{estrofas} de cuatro \textbf{versos} (dos cuartetos).} (Structure exacte du poème support étudié en Exercice 9).
\end{enumerate}
\end{corrige}

\begin{corrige}[Exercice 4 — Conjugaison : le subjonctif de l'opinion]
\textbf{Rappel :} Après les expressions d'opinion, de doute, de nécessité, d'émotion (\esp{es importante que, no creo que, es necesario que, dudo que, aunque} + subjonctif pour concession non réalisée), le verbe qui suit \esp{que} se met au \textbf{subjonctif présent}.
\begin{enumerate}
\item \esp{Es importante que nosotros \textbf{comamos} más verduras.} (\esp{comer} $\rightarrow$ 1\textsuperscript{re} pers. pl. subj. \esp{-emos}).
\item \esp{No creo que la comida basura \textbf{sea} buena para la salud.} (\esp{ser} $\rightarrow$ 3\textsuperscript{e} pers. sg. subj. irrégulier \esp{sea}).
\item \esp{Es necesario que el gobierno \textbf{prohíba} esa publicidad.} (\esp{prohibir} $\rightarrow$ 3\textsuperscript{e} pers. sg. subj. \esp{-a} ; accent sur \esp{í} pour garder la tonicité).
\item \esp{Dudo que los niños \textbf{sepan} elegir bien.} (\esp{saber} $\rightarrow$ 3\textsuperscript{e} pers. pl. subj. irrégulier \esp{sepan}).
\item \esp{Aunque las verduras \textbf{parezcan} menos atractivas, hay que comerlas.} (\esp{parecer} $\rightarrow$ 3\textsuperscript{e} pers. pl. subj. \esp{-an} ; \esp{aunque} + subjonctif = concession hypothétique/non réalisée).
\end{enumerate}
\end{corrige}

\begin{corrige}[Exercice 5 — Thème : phrases à traduire]
\begin{enumerate}
\item \esp{Aunque las verduras sean menos atractivas que las patatas fritas, es indispensable que las comamos para mantenernos sanos.}
\begin{itemize}
\item \esp{Aunque} + subjonctif (\esp{sean}) pour la concession réelle (fait avéré).
\item \esp{Es indispensable que} + subjonctif (\esp{comamos}, 1\textsuperscript{re} pl.).
\item \esp{Patatas fritas} (Espagne) / \esp{papas fritas} (Amérique) : les deux acceptés.
\item \esp{Mantenerse sanos} / \esp{gozar de buena salud}.
\end{itemize}
\item \esp{No creo que la comida basura sea más barata que una comida casera.}
\begin{itemize}
\item \esp{No creo que} + subjonctif (\esp{sea}, 3\textsuperscript{e} sg.).
\item \esp{Comida casera} = repas fait maison.
\item Comparatif : \esp{más barata que}.
\end{itemize}
\item \esp{Es necesario que las escuelas prohíban las bebidas azucaradas en los comedores.}
\begin{itemize}
\item \esp{Es necesario que} + subjonctif (\esp{prohíban}, 3\textsuperscript{e} pl.).
\item \esp{Bebidas azucaradas} = boissons sucrées. \esp{Comedores} / \esp{cantinas} = cantines.
\end{itemize}
\item \esp{A mi juicio, es urgente que los padres cocinen con sus hijos.}
\begin{itemize}
\item \esp{A mi juicio} / \esp{En mi opinión} / \esp{Para mí}.
\item \esp{Es urgente que} + subjonctif (\esp{cocinen}, 3\textsuperscript{e} pl.).
\end{itemize}
\end{enumerate}
\end{corrige}

\begin{corrige}[Exercice 6 — Transformations : du style direct à l'argumentation]
\begin{enumerate}
\item \esp{A mi juicio, no es bueno que los jóvenes coman comida basura.} (Ou : \esp{A mi juicio, la comida basura no debería ser tan accesible.}) — Nuance par l'opinion personnelle + subjonctif (\esp{coman}, 3\textsuperscript{e} pl.) ou conditionnel.
\item \esp{Es evidente que los jóvenes comen mal.} — Renforcement par \esp{es evidente que} + \textbf{indicatif} (\esp{comen}, 3\textsuperscript{e} pl.) car c'est présenté comme un fait certain.
\item \esp{Es necesario que nosotros cocinemos en casa.} — Structure d'obligation/nécessité impersonnelle + subjonctif (\esp{cocinemos}, 1\textsuperscript{re} pl.) avec sujet explicite \esp{nosotros}.
\item \esp{No creo que la publicidad influya en los niños.} — Doute \esp{no creo que} + subjonctif (\esp{influya}, 3\textsuperscript{e} sg. car \esp{la publicidad} est singulier).
\end{enumerate}
\end{corrige}

\begin{corrige}[Exercice 7 — Texte à trous : un article sur la cantine]
\begin{textebox}[La cantine du lycée — Corrigé]
\esp{En nuestro liceo de Yaoundé, muchos alumnos compran frituras a la salida, \textbf{aunque} la cantina ofrece platos equilibrados. Los fritos son baratos y rápidos; \textbf{sin embargo}, contienen demasiada grasa. \textbf{Además}, no aportan las vitaminas que necesitan los adolescentes. \textbf{Por eso}, el director quiere crear un club de cocina saludable. \textbf{En conclusión}, comer bien es también aprender a vivir bien.}
\end{textebox}
\textbf{Justifications :}
\begin{itemize}
\item \esp{aunque} : concession (malgré le fait que la cantine offre des plats équilibrés).
\item \esp{sin embargo} : opposition/adversative (mais ils contiennent trop de graisse).
\item \esp{además} : addition d'argument (de plus, pas de vitamines).
\item \esp{por eso} : conséquence (c'est pourquoi le directeur agit).
\item \esp{en conclusión} : synthèse finale.
\end{itemize}
\end{corrige}

\begin{corrige}[Exercice 8 — Appariements : figures de style]
\begin{enumerate}
\item \esp{« La hamburguesa me sonríe con dientes de grasa »} $\rightarrow$ \textbf{c) personificación.} La hamburguesa (objet) sourit (action humaine) et a des dents.
\item \esp{« El brócoli es un soldado verde de mi cuerpo »} $\rightarrow$ \textbf{a) metáfora.} Identification directe : le brocoli = un soldat (pas de \esp{como}). Il défend le corps.
\item \esp{« Dulce, dulce, dulce trampa dorada »} $\rightarrow$ \textbf{e) aliteración.} Répétition du son \esp{/d/} et \esp{/l/} (\esp{dulce, dulce, dulce, dorada}) et aussi anaphore (\esp{dulce} répété), mais la figure \textbf{sonore} dominante est l'allitération.
\item \esp{« Podría comer mil montañas de papas fritas »} $\rightarrow$ \textbf{d) hipérbole.} Exagération démesurée (\esp{mil montañas}) pour exprimer l'envie addictive.
\item \esp{« La manzana cruje como un tambor de la mañana »} $\rightarrow$ \textbf{b) comparación.} Comparaison explicite avec \esp{como} entre le bruit de la pomme et le tambour (fraîcheur, vitalité).
\item \esp{« No quiero azúcar, no quiero grasa, no quiero mentiras »} $\rightarrow$ \textbf{f) anáfora.} Répétition de \esp{« No quiero »} en début de segments (hémistiches/vers) pour insister sur le refus radical.
\end{enumerate}
\end{corrige}

\begin{corrige}[Exercice 9 — Compréhension et commentaire de poème]
\textbf{Barème indicatif : 16 points}

\textbf{1. Présentation du poème (4 pts)}
\begin{itemize}
\item \textbf{Thème :} Le conflit intérieur (\esp{duelo}) face au choix alimentaire : attraction de la malbouffe (gras, sucre, rapidité) vs conscience des bienfaits de l'alimentation saine (légumes, fruits, vitalité). C'est un poème engagé pour la santé.
\item \textbf{Structure :} 2 strophes (\esp{estrofas}) de 4 vers (\esp{cuartetos}). Vers de mesure variable (art majeur, environ 9-11 syllabes), rimes assonantes ou libres (vers libre moderne).
\item \textbf{Ton général :} D'abord séducteur/tentateur (v. 1-4), puis résolu/engagé (v. 5-8). Ton personnel (\esp{me, mi}) et direct.
\end{itemize}

\textbf{2. Figures de style v. 1-4 (4 pts — 2 figures $\times$ 2 pts)}
\begin{itemize}
\item \textbf{Personificación (v. 2) :} \esp{« la hamburguesa me sonríe con dientes de grasa »}. La hamburguesa est dotée d'intention (sourire) et d'anatomie humaine (dents de graisse). \textbf{Effet :} Rend la malbouffe vivante, séduisante, presque prédatrice ; la graisse devient un sourire carnassier.
\item \textbf{Hipérbole (v. 3) :} \esp{« Podría comer mil montañas de papas fritas »}. \textbf{Effet :} Montre l'addiction, l'absence de satiété, le pouvoir hypnotique de la « comfort food ».
\item \textit{(Autres possibles : Métaphore v. 1 \esp{« trampa dorada »} = piège attractif ; Anaphore/Allitération v. 4 \esp{« dulce, dulce, dulce »} = insistance sur le goût sucré).}
\end{itemize}

\textbf{3. Analyse v. 5-8 : Personificación et Metáfora (4 pts)}
\begin{itemize}
\item \textbf{Métaphore (v. 5) :} \esp{« el brócoli es un soldado verde de mi cuerpo »}. Identification : le brocoli = un soldat (humain) qui défend l'organisme. \textbf{Effet :} Valorise le légume : il devient un allié actif, un protecteur, non plus une contrainte.
\item \textbf{Comparación (v. 6) :} \esp{« la manzana cruje como un tambor de la mañana »}. Comparaison explicite (\esp{como}). \textbf{Effet :} Le bruit du croquant évoque la vitalité, le rythme sain de la vie, l'éveil.
\item \textbf{Métaphore (v. 8) :} \esp{« mi plato es mi escudo, mi cocina, mi bandera »}. Identification multiple : assiette = bouclier (protection), cuisine = lieu de pouvoir/soin, drapeau = identité/fierté. \textbf{Effet :} Transforme l'acte de manger en acte citoyen, identitaire et défensif. L'alimentation saine devient un engagement de vie.
\end{itemize}

\textbf{4. Titre : « El duelo del plato » (2 pts)}
\esp{Duelo} = duel, combat, mais aussi deuil. Le « plato » (assiette) est le champ de bataille. C'est un duel entre deux modèles alimentaires (malbouffe vs sain), entre pulsion et raison, entre maladie et santé. Le « je » doit trancher. Le titre résume la tension dramatique du poème.

\textbf{5. Traduction des deux derniers vers (2 pts)}
\esp{« No quiero azúcar, no quiero grasa, no quiero mentiras : mi plato es mi escudo, mi cocina, mi bandera. »}
\textbf{Traduction :} « Je ne veux pas de sucre, je ne veux pas de graisse, je ne veux pas de mensonges : mon assiette est mon bouclier, ma cuisine, mon drapeau. »
\begin{itemize}
\item \esp{No quiero} : Je ne veux pas (répétition anaphorique conservée).
\item \esp{Mentiras} : Mensonges (métaphore de la publicité trompeuse).
\item \esp{Escudo / bandera} : Bouclier / Drapeau (symboles forts).
\end{itemize}
\end{corrige}

\begin{corrige}[Exercice 10 — Thème et version]
\textbf{Barème : 20 points (10 + 10)}

\textbf{A. Thème (10 pts — $\approx$ 3,33 pts par phrase)}

\begin{enumerate}
\item \esp{Aunque las verduras sean menos atractivas que las patatas fritas, es indispensable que las comamos para mantenernos sanos.}
\begin{itemize}
\item \esp{Aunque} + subj. \esp{sean} (concession réelle).
\item \esp{Menos atractivas que} (comparatif d'infériorité).
\item \esp{Es indispensable que} + subj. \esp{comamos} (1\textsuperscript{re} pl.).
\item \esp{Mantenernos sanos} / \esp{gozar de buena salud}.
\end{itemize}
\item \esp{Es hora de que el gobierno prohíba la publicidad de la comida basura durante los programas infantiles.}
\begin{itemize}
\item \esp{Es hora de que} + \textbf{subjonctif} (\esp{prohíba}, 3\textsuperscript{e} sg.) — \textbf{Piège :} \esp{Es hora de que} exige le subjonctif (souhait/nécessité non réalisée).
\item \esp{Publicidad de la comida basura} / \esp{publicidad de comida chatarra}.
\item \esp{Durante los programas infantiles} / \esp{para niños}.
\end{itemize}
\item \esp{Cuanto más cocinamos en casa, menos enfermamos.} (Ou : \esp{A medida que cocinamos más en casa, enfermamos menos.})
\begin{itemize}
\item Structure correlative : \esp{Cuanto más [indicatif], menos [indicatif]} (faits constatés, indicatif).
\item \esp{Cocinamos} (indicatif présent), \esp{enfermamos} (indicatif présent).
\end{itemize}
\end{enumerate}

\textbf{B. Version (10 pts)}

\textbf{Texte source :}
\begin{textebox}[Texte à traduire]
\esp{Según los médicos, la obesidad infantil ha aumentado en muchos países durante los últimos años. Los expertos afirman que es urgente que las familias cambien sus hábitos : menos refrescos azucarados, más frutas y más deporte. Aunque las campañas de prevención son cada vez más numerosas, todavía no se ha conseguido que los niños prefieran una manzana a una bolsa de patatas fritas.}
\end{textebox}

\textbf{Traduction modèle :}
« Selon les médecins, l'obésité infantile a augmenté dans de nombreux pays au cours des dernières années. Les experts affirment qu'il est urgent que les familles changent leurs habitudes : moins de boissons gazeuses sucrées, plus de fruits et plus de sport. Bien que les campagnes de prévention soient de plus en plus nombreuses, on n'a pas encore réussi à faire en sorte que les enfants préfèrent une pomme à un sachet de chips. »

\textbf{Points de vigilance (version) :}
\begin{itemize}
\item \esp{Según los médicos} : Selon les médecins (pas « d'après » mais « selon »).
\item \esp{Ha aumentado} : Passé composé (\esp{pretérito perfecto compuesto}) $\rightarrow$ « a augmenté » (récent/résultat présent).
\item \esp{Es urgente que las familias cambien} : Subjonctif \esp{cambien} (3\textsuperscript{e} pl.) $\rightarrow$ « que les familles changent » (subjonctif fr. après « il est urgent que »).
\item \esp{Menos refrescos... más frutas} : Ellipse de verbe $\rightarrow$ « moins de... plus de... ».
\item \esp{Aunque las campañas... son} : \esp{son} (indicatif) car fait réel $\rightarrow$ « Bien que les campagnes soient » (subjonctif fr. après \esp{bien que}).
\item \esp{No se ha conseguido que... prefieran} : Passif réfléchi \esp{se ha conseguido} + subjonctif \esp{prefieran} (3\textsuperscript{e} pl.) $\rightarrow$ « on n'a pas réussi à ce que... préfèrent / on n'a pas réussi à faire préférer ».
\item \esp{Una bolsa de patatas fritas} : Un sachet de chips / frites.
\end{itemize}
\end{corrige}

\begin{corrige}[Exercice 11 — Expression écrite et orale]
\textbf{Barème total : 20 points (12 + 8)}

\textbf{A. Expression écrite — Lettre ouverte au Ministre (12 pts)}

\textbf{Critères d'évaluation :}
\begin{itemize}
\item \textbf{Registre formel (2 pts) :} Formules d'appel (\esp{Estimado Sr. Ministro}), de politesse (\esp{Atentamente}), vouvoiement (\esp{usted}), syntaxe soignée.
\item \textbf{Connecteurs argumentatifs (2 pts) :} Minimum 4 utilisés correctement (\esp{En primer lugar, Además, Sin embargo, Por lo tanto, En consecuencia, Para concluir}).
\item \textbf{Structures avec subjonctif (2 pts) :} Minimum 2 (\esp{Es necesario que... prohíba, Exijo que... se retire, Recomiendo que... se regule}).
\item \textbf{Vocabulaire précis santé (2 pts) :} \esp{obesidad infantil, diabetes tipo 2, hábitos nutricionales, publicidad engañosa, alimentos ultraprocesados, etiquetado nutricional}.
\item \textbf{Cohérence/Longueur (2 pts) :} 180-220 mots, structure lettre (intro, arguments, conclusion, appel à l'action).
\item \textbf{Correction linguistique (2 pts) :} Accords, accents, conjugaisons, orthographe.
\end{itemize}

\textbf{Proposition de rédaction modèle (environ 200 mots) :}

\begin{textebox}[Modèle de lettre ouverte]
\esp{Estimado Sr. Ministro de Salud:}\\[0.5em]
\esp{Le escribo como ciudadano preocupado por el alarmante aumento de la obesidad infantil en nuestro país. \textbf{En primer lugar}, los niños están expuestos a una avalancha de publicidad engañosa que presenta los alimentos ultraprocesados como premios de felicidad. \textbf{Además}, estos productos, ricos en azúcares y grasas saturadas, causan diabetes tipo 2 e hipertensión cada vez más tempranas.}\\[0.5em]
\esp{\textbf{Sin embargo}, la legislación actual permite que estos anuncios inunden los horarios infantiles. \textbf{Por lo tanto}, \textbf{es imperativo que su ministerio prohíba} inmediatamente la publicidad de comida basura dirigida a menores, tal como ya han hecho Chile o México con resultados positivos. \textbf{Es fundamental que se establezca} un etiquetado frontal de advertencia obligatorio y se promueva la educación nutricional en las escuelas.}\\[0.5em]
\esp{No podemos permitir que el beneficio económico de unas pocas empresas prime sobre la salud de las futuras generaciones. \textbf{En conclusión}, le exijo que actúe con valentía y urgencia. La salud de nuestros niños no es negociable.}\\[0.5em]
\esp{Atentamente,}\\[0.5em]
\esp{[Nombre y Apellidos]}\\
\esp{Estudiante de Terminale A, Liceo de Yaoundé}
\end{textebox}

\textbf{Points de vigilance (rédaction) :}
\begin{itemize}
\item Ne pas confondre \esp{comida basura} / \esp{comida chatarra} / \esp{alimentos ultraprocesados}.
\item Subjonctif après \esp{es imperativo que, es fundamental que, exijo que, recomiendo que}.
\item \esp{Dirigida a menores} (participe accordé au féminin avec \esp{publicidad}), pas \esp{dirigido}.
\item \esp{Prime} (verbe \esp{primar}, indicatif présent 3\textsuperscript{e} sg.) ou \esp{prevalezca} (subj. de \esp{prevalecer} après \esp{no podemos permitir que}).
\end{itemize}

\textbf{B. Expression orale — Plat traditionnel camerounais (8 pts)}

\textbf{Critères :}
\begin{itemize}
\item \textbf{Contenu/Vocabulaire (2 pts) :} Nom du plat, ingrédients, préparation, valeur nutritionnelle.
\item \textbf{Connecteurs (2 pts) :} Min. 3 (\esp{En primer lugar, Además, Por otro lado, Por eso, En definitiva}).
\item \textbf{Métaphore (1 pt) :} Ex. \esp{« El ndolé es el corazón verde de nuestra mesa »}, \esp{« El kgwem es un tesoro del bosque »}.
\item \textbf{Structure opinion + subjonctif (1 pt) :} Ex. \esp{Creo que es vital que comamos...}, \esp{Me parece esencial que se valore...}.
\item \textbf{Prononciation/Fluence (2 pts) :} 2 minutes, débit naturel, intonation expressive.
\end{itemize}

\textbf{Proposition de canevas oral (exemple : \esp{Ndolé} ou \esp{Kpwem/Kgwem}) :}

\end{corrige}

\begin{exemplebox}[Canevas oral : Le Ndolé]
\esp{« Buenas tardes. Hoy quiero hablarles del \textbf{ndolé}, mi plato favorito de Camerún. \textbf{En primer lugar}, el ndolé no es solo comida, \textbf{es el corazón verde de nuestra mesa}: una metáfora de la vida y la salud. \textbf{Además}, se prepara con hojas de \textit{Vernonia amygdalina} (ndolé), ricas en hierro, fibra y antioxidantes, cocinadas con pescado ahumado, cacahuetes y especias locales, sin aceite excesivo. \textbf{Por otro lado}, a diferencia de las frituras industriales, el ndolé nutre profundamente y previene la anemia. \textbf{Por eso}, \textbf{creo que es fundamental que las familias recuperen} estas recetas ancestrales. \textbf{En definitiva}, comer ndolé es comer identidad, es elegir la vida. ¡Gracias! »}
\end{exemplebox}

\textbf{Points de vigilance (oral) :}
\begin{itemize}
\item Bien prononcer les voyelles ouvertes/fermées (\esp{ndolé, corazón, vida}).
\item Respecter l'accent tonique (\esp{fundamEntal, recupEren, identidAd}).
\item Utiliser le subjonctif déclenché par \esp{es fundamental que} (\esp{recuperen}, 3\textsuperscript{e} pl.).
\item La métaphore doit être explicite (\esp{es el corazón verde...}).
\end{itemize}


\newpage

\coursec{Fiche bilan}

\begin{popfigure}
\begin{tikzpicture}[mindmap, grow cyclic, every node/.style={concept}, root concept/.append style={concept color=popblue, font=\small\bfseries, text=white}, level 1 concept/.append style={level distance=3.4cm, sibling angle=51.4, font=\scriptsize}, level 2 concept/.append style={level distance=2.1cm, font=\tiny}]
\node[root concept, align=center]{E11\\ Poema y debate}
  child[concept color=poporange]{ node {Lectura del poema}
    child { node {tema} }
    child { node {estructura} }
  }
  child[concept color=popgreen]{ node {Léxico poético}
    child { node {verso, estrofa} }
    child { node {rima, poeta} }
  }
  child[concept color=poppurple]{ node {Versificación}
    child { node {métrica} }
    child { node {rima asonante} }
  }
  child[concept color=poppink]{ node {Figuras retóricas}
    child { node {metáfora} }
    child { node {anáfora} }
  }
  child[concept color=popgold]{ node {Comentario}
    child { node {intro, análisis} }
    child { node {conclusión} }
  }
  child[concept color=popblue]{ node {Debate}
    child { node {comida sana} }
    child { node {comida basura} }
  }
  child[concept color=popgreen!70!black]{ node {Opinión}
    child { node {pienso que} }
    child { node {estoy a favor} }
  };
\end{tikzpicture}
\legende{Carte mentale de la leçon E11 : poème, figures de style et débat sur l'alimentation.}
\end{popfigure}

\begin{bilan}
\textbf{1. Lexique clé de la poésie}

\begin{center}
\begin{tabularx}{\linewidth}{|X|X|}\hline
\rowcolor{popblueL} \textbf{Español} & \textbf{Français} \\ \hline
el verso & le vers \\ \hline
la estrofa & la strophe \\ \hline
la rima (asonante / consonante) & la rime (assonance / rime riche) \\ \hline
el poeta / la poetisa & le poète / la poétesse \\ \hline
la musa & la muse \\ \hline
el poema / la poesía & le poème / la poésie \\ \hline
el título, el tema, el tono & le titre, le thème, le ton \\ \hline
el hablante lírico & le locuteur lyrique (le « je » du poème) \\ \hline
\end{tabularx}
\end{center}

\textbf{2. Les six figures de style à connaître}

\begin{center}
\begin{tabularx}{\linewidth}{|X|X|X|}\hline
\rowcolor{popblueL} \textbf{Figura} & \textbf{Définition} & \textbf{Ejemplo} \\ \hline
metáfora & identification directe, sans outil de comparaison & \esp{El plato es un campo de batalla.} \\ \hline
comparación (símil) & comparaison avec \esp{como}, \esp{cual}, \esp{parece} & \esp{duro como una piedra} \\ \hline
personificación & donner des traits humains à un objet & \esp{El tomate llora en la mesa.} \\ \hline
hipérbole & exagération volontaire & \esp{He comido mil hamburguesas.} \\ \hline
aliteración & répétition d'un même son & \esp{el ruido roto del río} \\ \hline
anáfora & répétition en début de vers & \esp{Quiero pan, quiero fruta, quiero vida.} \\ \hline
\end{tabularx}
\end{center}

\textbf{3. Méthode express du commentaire de poème}
\begin{enumerate}
\item \textbf{Introducción} : présenter le poème (titre, thème, situation) et annoncer le plan.
\item \textbf{Análisis temático} : dégager le thème principal et les thèmes secondaires.
\item \textbf{Análisis formal} : structure (estrofas, versos), versificación, rima.
\item \textbf{Figuras retóricas} : citer chaque figure, la nommer, expliquer son effet.
\item \textbf{Conclusión} : bilan + ouverture (opinión personal).
\end{enumerate}

\textbf{4. Lexique du débat : alimentación}

\begin{center}
\begin{tabularx}{\linewidth}{|X|X|}\hline
\rowcolor{popblueL} \textbf{Español} & \textbf{Français} \\ \hline
la alimentación sana / equilibrada & l'alimentation saine / équilibrée \\ \hline
la comida basura & la malbouffe \\ \hline
los alimentos naturales / procesados & les aliments naturels / transformés \\ \hline
la obesidad, la diabetes & l'obésité, le diabète \\ \hline
las grasas, los azúcares, la sal & les graisses, les sucres, le sel \\ \hline
las frutas y verduras frescas & les fruits et légumes frais \\ \hline
cocinar en casa & cuisiner à la maison \\ \hline
llevar una vida sana & mener une vie saine \\ \hline
\end{tabularx}
\end{center}

\textbf{5. Règles de grammaire à connaître par cœur}

\begin{center}
\begin{tabularx}{\linewidth}{|X|X|}\hline
\rowcolor{popblueL} \textbf{Structure d'opinion} & \textbf{Mode} \\ \hline
\esp{Pienso / Creo / Opino que + frase} & indicatif \\ \hline
\esp{No pienso / No creo que + frase} & subjonctif : \esp{No creo que sea sano.} \\ \hline
\esp{Me parece que + frase} & indicatif \\ \hline
\esp{Es importante / Es necesario que + frase} & subjonctif : \esp{Es necesario que comas verduras.} \\ \hline
\esp{Estoy a favor de / Estoy en contra de + infinitivo o nombre} & \esp{Estoy a favor de cocinar en casa.} \\ \hline
\esp{En mi opinión, / Desde mi punto de vista, / Sin embargo, / Además, / Por eso,} & connecteurs du débat \\ \hline
\end{tabularx}
\end{center}

\textbf{6. Méthode express du débat oral}
\begin{enumerate}
\item Prendre position clairement : \esp{Yo defiendo la alimentación sana porque...}
\item Donner au moins deux arguments + un exemple concret.
\item Réagir à l'adversaire : \esp{Entiendo tu punto de vista, pero...} / \esp{No estoy de acuerdo porque...}
\item Conclure : \esp{En conclusión, pienso que...}
\end{enumerate}
\end{bilan}

\begin{aretenir}[Checklist avant le Bac]
\begin{itemize}
\item[$\square$] Je sais définir et citer en espagnol : \esp{verso}, \esp{estrofa}, \esp{rima}, \esp{poeta}, \esp{musa}.
\item[$\square$] Je sais compter les vers et identifier les strophes d'un poème.
\item[$\square$] Je reconnais et je définis les six figures : \esp{metáfora}, \esp{comparación}, \esp{personificación}, \esp{hipérbole}, \esp{aliteración}, \esp{anáfora}.
\item[$\square$] Je sais citer un vers du poème pour justifier chaque figure identifiée.
\item[$\square$] Je connais le plan du commentaire : introducción, análisis temático, análisis formal, figuras, conclusión.
\item[$\square$] Je maîtrise le vocabulaire du débat : \esp{alimentación sana}, \esp{comida basura}, \esp{obesidad}, \esp{alimentos procesados}.
\item[$\square$] J'emploie correctement l'indicatif après \esp{pienso que} et le subjonctif après \esp{no creo que} et \esp{es importante que}.
\item[$\square$] Je sais exprimer mon accord et mon désaccord : \esp{estoy a favor de}, \esp{estoy en contra de}, \esp{no estoy de acuerdo}.
\item[$\square$] J'utilise les connecteurs : \esp{además}, \esp{sin embargo}, \esp{por eso}, \esp{en conclusión}.
\item[$\square$] Je fais attention aux faux amis : \esp{la comida} (la nourriture, pas « la comédie »), \esp{sano} (sain).
\item[$\square$] Je vérifie les accents écrits : \esp{opinión}, \esp{alimentación}, \esp{metáfora}, \esp{análisis}.
\item[$\square$] Je peux rédiger un court article d'opinion structuré en trois paragraphes.
\end{itemize}
\end{aretenir}

\findecours

\end{document}
